%% 本文件由 工程文件/scripts/assemble.py 生成；第二轮校订（r-ch04）已把 units 内容内联于此，以便直接修订；源码/units/ 保持译者原件不动。

% ===================== begin 源码/units/ch04_u01.tex =====================
% ch04_u01 —— 原书 p.122–p.131（4.0 全部；4.1 自 4.1.1 起，至 p.131 末句中断处）
\bdchap{ch:pfaff}{Pfaffian 及其相关}
\SourcePage{122}

\bdsec{4.0. 引言}

\bdp{4.0.1} 考虑「归一化」的典范丛
\begin{equation}\tag{146}
\omega^{\sharp}_{\operatorname{Bun}_G}:=\omega_{\operatorname{Bun}_G}\otimes\omega_0^{\otimes -1},
\end{equation}
其中 $\omega_0$ 是 $\omega_{\operatorname{Bun}_G}$ 在 $\operatorname{Bun}_G$ 中对应于 $X$ 上平凡 $G$-丛那一点处的纤维。本节将给每个 ${}^{L}G$-oper $\mathfrak{F}$ 配置 $\operatorname{Bun}_G$ 上的可逆层 $\lambda_{\mathfrak{F}}$，即 0.2(d) 中提到的那一个。层 $\lambda_{\mathfrak{F}}$ 上会装一个同构 $\lambda_{\mathfrak{F}}^{\otimes 2n}\xrightarrow{\ \sim\ }(\omega^{\sharp}_{\operatorname{Bun}_G})^{\otimes n}$，其中 $n\ne 0$。该同构诱导出 0.2(d) 所要求的、$\lambda_{\mathfrak{F}}$ 上的扭 $D$-模结构。

由 3.4.3 的公式 (57)，$\operatorname{Op}_{{}^{L}G}(X)=\operatorname{Op}_{{}^{L}\mathfrak{g}}(X)\times Z\operatorname{tors}_{\theta}(X)$，其中 $Z$ 是 ${}^{L}G$ 的中心。实际上 $\lambda_{\mathfrak{F}}$ 只依赖 $\mathfrak{F}$ 在 $Z\operatorname{tors}_{\theta}(X)$ 中的像。于是目标就是构造一个典范函子
\begin{equation}\tag{147}
\lambda:Z\operatorname{tors}_{\theta}(X)\to\mu_\infty\operatorname{tors}_{\theta}(\operatorname{Bun}_G),
\end{equation}
其中 $\mu_\infty\operatorname{tors}_{\theta}(\operatorname{Bun}_G)$ 是如下群胚：对象为 $\operatorname{Bun}_G$ 上的线丛 $\mathcal{A}$，连同同构 $\mathcal{A}^{\otimes 2n}\xrightarrow{\ \sim\ }(\omega^{\sharp}_{\operatorname{Bun}_G})^{\otimes n}$（对某个 $n\ne 0$）。

\bdp{4.0.2} 若 $G$ 单连通，(147) 的构造相当简单。此时 $Z$ 平凡，只需构造 $\mu_\infty\operatorname{tors}_{\theta}(\operatorname{Bun}_G)$ 的一个对象。由于 $G$ 单连通，$\operatorname{Bun}_G$ 连通且单连通（把 $X$ 上的 $G$-丛解释为 $X$ 对应 $C^\infty$ 流形上的 $G$-丛，并带上一个 $\bar\partial$-联络）。于是问题化为说明 $\omega^{\sharp}_{\operatorname{Bun}_G}$ 有平方根（此时 $\mu_\infty\operatorname{tors}_{\theta}(\operatorname{Bun}_G)$ 有唯一对象，该对象在 $\operatorname{Bun}_G$ 中对应于平凡 $G$-丛那一点处的纤维被平凡化）。为解决此问题，我们使用 Pfaffian（普法夫式）的概念。

对带非退化对称型 $\mathcal{Q}\otimes\mathcal{Q}\to\omega_X$ 的向量丛 $\mathcal{Q}$，Laszlo 与 Sorger 在 [La-So] 中给 $\mathcal{Q}$ 配置其 Pfaffian $\operatorname{Pf}(\mathcal{Q})$，即 $\det R\Gamma(X,\mathcal{Q})$ 的典范平方根。4.2 中我们给出 Pfaffian 的另一个定义，它应当与 [La-So] 中的定义等价。

\SourcePage{123}

固定 $\mathcal{L}\in\omega^{1/2}(X)$（即 $\mathcal{L}$ 是 $\omega_X$ 的一个平方根）。则 $\operatorname{Bun}_G$ 上以
\begin{equation}\tag{148}
\operatorname{Pf}(\mathfrak{g}_{\mathfrak{F}}\otimes\mathcal{L})\otimes\operatorname{Pf}(\mathfrak{g}\otimes\mathcal{L})^{\otimes -1}
\end{equation}
为 $\mathfrak{F}\in\operatorname{Bun}_G$ 处纤维的线丛，是 $\omega^{\sharp}_{\operatorname{Bun}_G}$ 的一个平方根（细节见 4.3.1）。

于是要理解 $G$ 单连通的情形，只需通读 4.2 与 4.3.1。一般情形下 (147) 的构造更复杂。要点在于：(148) 所定义的 $\omega^{\sharp}_{\operatorname{Bun}_G}$ 的平方根依赖 $\mathcal{L}\in\omega^{1/2}(X)$。

\bdp{4.0.3} 下面概述对任意半单 $G$ 构造 (148) 的思路。

如 3.4.6 所述，$Z\operatorname{tors}_{\theta}(X)$ 是 Picard 范畴 $Z\operatorname{tors}(X)$ 上的一个 torsor，而 $\mu_\infty\operatorname{tors}_{\theta}(\operatorname{Bun}_G)$ 是 Picard 范畴
\begin{equation}\tag{149}
\mu_\infty\operatorname{tors}(\operatorname{Bun}_G):=\lim_{\substack{\longrightarrow\\ n}}\mu_n\operatorname{tors}(\operatorname{Bun}_G)
\end{equation}
上的一个 torsor。

% [Tutor Note: 原书 p.123 式 (149) 的 lim 下方印的是右向箭头（正极限的写法），但 μ∞-torsor 应为 μ_n-torsor 的逆极限；此处照录原书箭头，未改。]

我们要构造的函子 (147)，对某个 Picard 函子 $\ell:Z\operatorname{tors}(X)\to\mu_\infty\operatorname{tors}(\operatorname{Bun}_G)$ 而言，是 3.4.6 意义下的 $\ell$-仿射函子。$\ell$ 在 4.1 中定义。torsor $Z\operatorname{tors}_{\theta}(X)$ 由 $\omega^{1/2}(X)$ 经某个 Picard 函子 $\mu_2\operatorname{tors}(X)\to Z\operatorname{tors}(X)$ 诱导而来（见 3.4.6）。因此要构造 $\lambda$，只需构造一个 $\ell'$-仿射函子 $\lambda':\omega^{1/2}(X)\to\mu_\infty\operatorname{tors}_{\theta}(X)$，其中 $\ell'$ 是复合 $\mu_2\operatorname{tors}(X)\to Z\operatorname{tors}(X)\xrightarrow{\ \ell\ }\mu_\infty\operatorname{tors}(\operatorname{Bun}_G)$。我们以 $\mathcal{L}\mapsto\lambda'_{\mathcal{L}}$ 定义 $\lambda'$，其中 $\lambda'_{\mathcal{L}}$ 是 $\operatorname{Bun}_G$ 上的线丛，它在 $\mathfrak{F}\in\operatorname{Bun}_G$ 处的纤维等于 (148)。$\lambda'$ 是 $\ell'$-仿射这一事实，在 4.4 中由 4.3.10 推出；4.3.10 是关于 $\mathrm{SO}_n$-丛的一般论断\footnote{事实上，4.3.10 是 [BLaSo] 中命题 5.2 的精细化。}。实际上，在专门讨论 Pfaffian 的 4.2 与 4.3 两小节里，群 $G$ 完全不出现。

\bdp{4.0.4} 本节构造的 $\operatorname{Bun}_G$ 上每个线丛都带有如下附加结构：对每个 $x\in X$，$G(K_x)$ 的某个中心扩张作用在它到 2.3.1 中概形 $\operatorname{Bun}_{G,x}$ 的拉回上。这一结构在 4.3 中用到，第 5 章还会用到。

\SourcePage{124}

% 节标题进 \addcontentsline/书签：unicode-math 的希腊字母在 PDF 字符串里会触发
% 「Improper alphabetic constant」，故只对 $\mu_\infty$ 加 \texorpdfstring，正文排版不变。
\bdsec{4.1. $\operatorname{Bun}_G$ 上的 \texorpdfstring{$\mu_\infty$}{μ∞}-torsor}

\bdp{4.1.1} 设 $G$ 是连通仿射代数群，$\Pi$ 是有限交换群，$0\to\Pi(1)\to\tilde G\to G\to 0$ 是 $G$ 的一个扩张。我们要构造典范 Picard 函子 $\ell:\Pi^{\vee}\operatorname{tors}(X)\to\mu_\infty\operatorname{tors}(\operatorname{Bun}_G)$，其中 $\Pi^{\vee}:=\operatorname{Hom}(\Pi,\mu_\infty)$。

\emph{注.} 若 $G$ 半单、$\tilde G$ 是 $G$ 的泛覆叠，则 $\Pi=\pi_1(G)$，而 $\Pi^{\vee}$ 典范同构于 ${}^{L}G$ 的中心 $Z$（该同构由 $G$ 与 ${}^{L}G$ 的嘉当环面之间的对偶诱导）。故此时 $\ell$ 是 Picard 函子 $Z\operatorname{tors}(X)\to\mu_\infty\operatorname{tors}(\operatorname{Bun}_G)$，正如 4.0.3 所许诺的。

我们在 4.1.2–4.1.4 中构造 $\ell$。4.1.5 对上述构造给出「解释」，4.1.6 稍作改写。4.1.7–4.1.9 考虑 $G(K_x)$ 的某个中心扩张的作用。4.1.10–4.1.11 描述 $\operatorname{Bun}_G$ 的基本群胚，这使 $\operatorname{Bun}_G$ 上 torsor 的构造变得清楚。读者可以略过 4.1.5 与 4.1.10–4.1.11。

\bdp{4.1.2} 对 $\mathfrak{F}\in\operatorname{Bun}_G$，记 $\tilde{\mathfrak{F}}$ 为 $X$ 上 $\mathfrak{F}$ 的 $\tilde G$-提升构成的 $\Pi(1)$-gerbe。其类 $c(\mathfrak{F})$ 是 $cl(\mathfrak{F})$ 在边缘态射 $H^1(X,G)\to H^2(X,\Pi(1))=\Pi$ 下的像。对有限非空子集 $S\subset X$，gerbe $\tilde{\mathfrak{F}}_{X\setminus S}$ 平凡。因此 $|\tilde{\mathfrak{F}}(X\setminus S)|$（即 $\tilde{\mathfrak{F}}(X\setminus S)$ 的对象的同构类集合）是非空的 $H^1(X\setminus S,\Pi(1))$-torsor。记之为 $\phi_{S,\mathfrak{F}}$。$\mathfrak{F}$ 变动时，$\phi_{S,\mathfrak{F}}$ 成为 $\operatorname{Bun}_G$ 上某个 $H^1(X\setminus S,\Pi(1))$-torsor $\phi_S$ 的纤维。

\bdp{4.1.3} 对任意 $x\in X$，集合 $|\tilde{\mathfrak{F}}(\operatorname{Spec} O_x)|$ 只有一个元素。用它可平凡化 $\Pi$-torsor $|\tilde{\mathfrak{F}}(\operatorname{Spec} K_x)|$（注意 $\Pi=H^1(\operatorname{Spec} K_x,\Pi(1))$）。于是限制到 $\operatorname{Spec} K_s$（$s\in S$）给出 $\operatorname{Res}_s$-仿射映射 $\operatorname{Res}_{s,\mathfrak{F}}:\phi_{S,\mathfrak{F}}\to\Pi$，其中 $\operatorname{Res}_s:H^1(X\setminus S,\Pi(1))\to\Pi$ 是 $s$ 处的留数。对 $c\in\Pi$，令 $\Pi^S_c:=\{\pi_S=(\pi_s):\sum\pi_s=c\}\subset\Pi^S$。映射 $\operatorname{Res}_{S,\mathfrak{F}}:=(\operatorname{Res}_{s,\mathfrak{F}}):\phi_{S,\mathfrak{F}}\to\Pi^S$ 的像是 $\Pi^S_{c(\mathfrak{F})}$。

\SourcePage{125}

\bdp{4.1.4} 回忆 $\Pi^{\vee}$ 是对偶于 $\Pi$ 的群，故有非退化配对 $(\ ,\ ):\Pi\times\Pi^{\vee}\to\mu_\infty$。

设 $\mathcal{E}$ 是 $X$ 上的 $\Pi^{\vee}$-torsor。置 $\mathcal{E}_S:=\prod_{s\in S}\mathcal{E}_s$，即 $\mathcal{E}$ 在 $S$ 处的平凡化集合；它是 $(\Pi^{\vee})^S$-torsor。对任意 $e\in\mathcal{E}_S$ 有类 $cl(\mathcal{E},e)\in H^1_c(X\setminus S,\Pi^{\vee})$。记 $\ell_{S,\mathcal{E},\mathfrak{F}}$ 为一个带映射
\begin{equation}\tag{150}
(\ ,\ )_{\ell}:\phi_{S,\mathfrak{F}}\times\mathcal{E}_S\to\ell_{S,\mathcal{E},\mathfrak{F}}
\end{equation}
的 $\mu_\infty$-torsor，使得对 $\varphi\in\phi_{S,\mathfrak{F}}$、$e=(e_s)\in\mathcal{E}_S$、$h\in H^1(X\setminus S,\Pi(1))$、$\chi=(\chi_s)\in(\Pi^{\vee})^S$，有
\begin{equation}\tag{151}
\begin{aligned}
(\varphi+h,e)_{\ell}&=(h,cl(\mathcal{E},e))_P(\varphi,e)_{\ell},\\
(\varphi,\chi e)_{\ell}&=(\operatorname{Res}_S\varphi,\chi)(\varphi,e)_{\ell}.
\end{aligned}
\end{equation}
这里 $(\ ,\ )_P:H^1(X\setminus S,\Pi(1))\times H^1_c(X\setminus S,\Pi^{\vee})\to\mu_\infty$ 是 Poincaré 配对，而 $(\operatorname{Res}_S\varphi,\chi):=\prod_{s\in S}(\operatorname{Res}_s\varphi,\chi_s)\in\mu_\infty$。这样的 $(\ell_{S,\mathcal{E},\mathfrak{F}},(\ ,\ )_{\ell})$ 存在且唯一。若 $S'\supset S$，则有显然的映射 $\phi_{S,\mathfrak{F}}\hookrightarrow\phi_{S',\mathfrak{F}}$、$\mathcal{E}_{S'}\to\mathcal{E}_S$，且存在唯一的 $\mu_\infty$-torsor 等同 $\ell_{S,\mathcal{E},\mathfrak{F}}=\ell_{S',\mathcal{E},\mathfrak{F}}$，使这些映射关于 $(\ ,\ )_{\ell}$ 互为伴随。于是这个 $\mu_\infty$-torsor 不依赖 $S$，我们简记之为 $\ell_{\mathcal{E},\mathfrak{F}}$。

$\mathfrak{F}$ 变动时，$\ell_{\mathcal{E},\mathfrak{F}}$ 成为 $\operatorname{Bun}_G$ 上 $\mu_\infty$-torsor $\ell_{\mathcal{E}}$ 的纤维。函子
\begin{equation}\tag{152}
\ell=\ell^{\tilde G}:\Pi^{\vee}\operatorname{tors}(X)\to\mu_\infty\operatorname{tors}(\operatorname{Bun}_G),
\end{equation}
$\mathcal{E}\mapsto\ell_{\mathcal{E}}$，带有显然的 Picard 函子结构。相应的自同构群同态 $\Pi^{\vee}\to\Gamma(\operatorname{Bun}_G,\mu_\infty)$ 是 $\chi\mapsto(c,\chi)$。

\emph{注.} 事实上 $\ell$ 是函子 $\Pi^{\vee}\operatorname{tors}(X)\to\mu_m\operatorname{tors}(\operatorname{Bun}_G)$，其中 $m$ 是 $\Pi$ 的阶。这可由构造看出，也可由 (152) 是 Picard 函子推出。

\SourcePage{126}

\bdp{4.1.5} 对交换群 $A$，记 $A\operatorname{gerbes}(X)$ 为与 $X$ 上 $A$-gerbe 的 2-范畴相联系的范畴（故 $A\operatorname{gerbes}(X)$ 是群胚，其对象是 $X$ 上的 $A$-gerbe，其态射是 1-态射模去 2-同构）。4.1.2–4.1.4 实际上构造了双 Picard 函子
\begin{equation}\tag{153}
\Pi^{\vee}\operatorname{tors}(X)\times\Pi(1)\operatorname{gerbes}(X)\to\mu_\infty\operatorname{tors},
\end{equation}
其中 $\mu_\infty\operatorname{tors}$ 记一点上 $\mu_\infty$-torsor 的范畴。本小节（读者可以略过）给出该构造的「科学解释」。

在 [Del73] 的 §1.4.11 中，Deligne 给满足 $K^i=0$（$i\ne 0,-1$）的交换群复形 $K^{\cdot}$ 配置一个 Picard 范畴。记这个 Picard 范畴为 $P(K^{\cdot})$。其对象是 $K^0$ 的元素；从 $x\in K^0$ 到 $y\in K^0$ 的态射是满足 $df=y-x$ 的元素 $f\in K^{-1}$。

4.1.4 中我们隐含地把 $\Pi^{\vee}\operatorname{tors}(X)$ 解释为 $P(K^{\cdot}_S)$，其中 $K^0_S=H^1_c(X\setminus S,\Pi^{\vee})$ 是 $X$ 上在 $S$ 上平凡化的 $\Pi^{\vee}$-torsor 的同构类集合，$K^{-1}_S=H^0(S,\Pi^{\vee})$。4.1.3 中我们隐含地把 $\Pi(1)\operatorname{gerbes}(X)$ 解释为 $P(L^{\cdot}_S)$，其中 $L^0_S=H^2_S(X,\Pi(1))=\Pi^S$，$L^{-1}_S=H^1(X\setminus S,\Pi(1))$（$L^0_S$ 参数化 $X$ 上在 $X\setminus S$ 上固定了一个对象的 $\Pi(1)$-gerbe）。4.1.4 中给出的双 Picard 函子 (153) 的构造，只用到典范配对 $K^{\cdot}_S\times L^{\cdot}_S\to\mu_\infty[1]$。

对 $S'\supset S$ 有典范拟同构 $K^{\cdot}_{S'}\to K^{\cdot}_S$ 与 $L^{\cdot}_S\to L^{\cdot}_{S'}$。相应的范畴等价 $P(K^{\cdot}_{S'})\to P(K^{\cdot}_S)$ 与 $P(L^{\cdot}_S)\to P(L^{\cdot}_{S'})$，与我们关于 $P(K^{\cdot}_S)$、$P(K^{\cdot}_{S'})$ 同 $\Pi^{\vee}\operatorname{tors}(X)$ 的等同相容，也与关于 $P(L^{\cdot}_S)$、$P(L^{\cdot}_{S'})$ 同 $\Pi(1)\operatorname{gerbes}(X)$ 的等同相容。关于配对 $K^{\cdot}_S\times L^{\cdot}_S\to\mu_\infty[1]$ 与 $K^{\cdot}_{S'}\times L^{\cdot}_{S'}\to\mu_\infty[1]$，态射 $L^{\cdot}_S\to L^{\cdot}_{S'}$ 是 $K^{\cdot}_{S'}\to K^{\cdot}_S$ 的伴随。因此 (153) 不依赖 $S$。

\emph{注.}
% 「Remarks」在原书排在 p.126 末行，(i)(ii) 两项在 p.127；故页锚点放在标题之后。
\SourcePage{127}
\begin{enumerate}[label=(\roman*),leftmargin=2.2em]
\item 比起 $K^{\cdot}_S$ 与 $L^{\cdot}_S$，用它们在导出范畴中的像更自然，即 $(\tau_{\le 1}R\Gamma(X,\Pi^{\vee}))[1]$ 与 $(\tau_{\ge 1}R\Gamma(X,\Pi(1)))[2]$。但通常的导出范畴不够用：按 [Del73] 的 §§1.4.13–1.4.14，$K^{\cdot}$ 在导出范畴中的像只能把 $P(K^{\cdot})$ 确定到等价为止，而这个等价又只到相差一个同构为止，且该同构并不唯一。因此需要一个精细化了的导出范畴概念，它大概在文献中还找不到。
\item 由配对 $K^{\cdot}_S\times L^{\cdot}_S\to\mu_\infty[1]$ 非退化，容易推出 (153) 诱导 $\Pi^{\vee}\operatorname{tors}(X)$ 与 Picard 函子 $\Pi(1)\operatorname{gerbes}(X)\to\mu_\infty\operatorname{tors}$ 的范畴之间的等价（这是 [Del73] 中等价 (1.4.18.1) 的特例）。
\end{enumerate}

\bdp{4.1.6} 4.1.4 中 $\ell_{\mathcal{E}}$ 的定义可改写如下。设 $S\subset X$ 有限非空。对固定的 $e\in\mathcal{E}_S$，有类 $c=cl(\mathcal{E},e)\in H^1_c(X\setminus S,\Pi^{\vee})$，从而有态射 $\lambda_e:H^1(X\setminus S,\Pi(1))\to\mu_\infty$，由 $\lambda_e(h)=(h,c)_P$ 定义。记 $\ell_{\mathcal{E},e}$ 为 4.1.2 中 $H^1(X\setminus S,\Pi(1))$-torsor $\phi_S$ 沿 $\lambda_e$ 的前推。对不同 $e\in\mathcal{E}_S$ 的 torsor $\ell_{\mathcal{E},e}$，按下述方式等同。

设 $\tilde e=\chi e$，$\chi\in(\Pi^{\vee})^S$。则 $\lambda_{\tilde e}(h)/\lambda_e(h)=(\operatorname{Res}_S(h),\chi)$，其中 $\operatorname{Res}_S$ 是边缘态射 $H^1(X\setminus S,\Pi(1))\to H^2_S(X,\Pi(1))=\Pi^S$。于是 $\ell_{\mathcal{E},\tilde e}/\ell_{\mathcal{E},e}$ 是 $\Pi^S$-torsor $(\operatorname{Res}_S)_*\phi_S$ 沿 $\chi:\Pi^S\to\mu_\infty$ 的前推。4.1.3 中的映射 $\operatorname{Res}_{S,\mathfrak{F}}:\phi_{S,\mathfrak{F}}\to\Pi^S$ 诱导 $(\operatorname{Res}_S)_*\phi_S$ 的典范平凡化，从而诱导典范同构 $\ell_{\mathcal{E},e}\xrightarrow{\ \sim\ }\ell_{\mathcal{E},\tilde e}$。于是可把不同 $e\in\mathcal{E}_S$ 的 $\ell_{\mathcal{E},e}$ 等同起来，得到 $\operatorname{Bun}_G$ 上的 $\mu_\infty$-torsor，它不依赖 $e\in\mathcal{E}_S$。显然它也不依赖 $S$。这就是 $\ell_{\mathcal{E}}$。

\bdp{4.1.7} 设 $S\subset X$ 是非空有限集，$O_S:=\prod_{x\in S}O_x$，$K_S:=\prod_{x\in S}K_x$，其中 $O_x$ 是 $x$ 的完全化局部环，$K_x$ 是其分式域。记 $\underline{S}$ 为 $S$ 的形式邻域，$\operatorname{Bun}_{G,\underline{S}}$ 为 $X$ 上在 $\underline{S}$ 上平凡化的 $G$-丛的模概形（2.3.1 中引入的 $\operatorname{Bun}_{G,x}$ 对应 $S=\{x\}$）。在 $\operatorname{Bun}_{G,\underline{S}}$ 上定义 $G(K_S)$ 的作用，
% 原书 p.127 末句与 p.128 首句是同一句（"One defines an action ... extending the action of G(O_S) by interpreting ..."），
% 故此处不空行：页锚点落在同一段中间。
\SourcePage{128}
它延拓 $G(O_S)$ 的作用：把 $X$ 上的 $G$-丛解释为 $X\setminus S$ 上的 $G$-丛，连同其拉回到 $\operatorname{Spec} K_S$ 的一个平凡化（见 2.3.4 与 2.3.7）。

设 $\ell_{\mathcal{E}}$ 是 $\operatorname{Bun}_G$ 上对应于 $X$ 上 $\Pi^{\vee}$-torsor $\mathcal{E}$ 的 $\mu_\infty$-torsor（见 4.1.4、4.1.6）。记 $\ell^S_{\mathcal{E}}$ 为 $\ell_{\mathcal{E}}$ 在 $\operatorname{Bun}_{G,\underline{S}}$ 上的原像。$G(O_S)$ 在 $\operatorname{Bun}_{G,\underline{S}}$ 上的作用典范地提升为它在 $\ell^S_{\mathcal{E}}$ 上的作用。我们断言：$\mathcal{E}$ 在 $S$ 上的平凡化定义了 $G(K_S)$ 在 $\ell^S_{\mathcal{E}}$ 上的作用，它延拓上述 $G(O_S)$ 的作用，并与 $G(K_S)$ 在 $\operatorname{Bun}_{G,\underline{S}}$ 上的作用相容。事实上，一旦取定 $e\in\mathcal{E}_S$，就可把 $\ell^S_{\mathcal{E}}$ 与 $\ell^S_{\mathcal{E},e}=(\lambda_e)_*\tilde\phi_S$ 等同，其中 $\tilde\phi_S$ 是 $\phi_S$ 到 $\operatorname{Bun}_{G,\underline{S}}$ 的拉回，$\lambda_e$ 已在 4.1.6 中定义。$G(K_S)$ 作用在 $\tilde\phi_S$ 上，因为 $\phi_{S,\mathfrak{F}}$ 只依赖 $\mathfrak{F}$ 到 $X\setminus S$ 的限制。故 $G(K_S)$ 作用在 $\ell^S_{\mathcal{E},e}$ 上。

由 4.1.6 的同构 $\ell_{\mathcal{E},e}\xrightarrow{\ \sim\ }\ell_{\mathcal{E},\tilde e}$ 诱导的同构 $\ell^S_{\mathcal{E},e}\xrightarrow{\ \sim\ }\ell^S_{\mathcal{E},\tilde e}$ 并\emph{不}是 $G(K_S)$-等变的。事实上，若 $\tilde e=\chi e$，$\chi\in(\Pi^{\vee})^S$，则按 4.1.6，$\ell^S_{\mathcal{E},\tilde e}/\ell^S_{\mathcal{E},e}$ 是 $\Pi^S$-torsor $(\operatorname{Res})_*\tilde\phi_S$ 沿 $\chi:\Pi^S\to\mu_\infty$ 的前推。4.1.6 中的等同 $(\operatorname{Res})_*\tilde\phi_S=\operatorname{Bun}_{G,\underline{S}}\times\Pi^S$，在 $G(K_S)$ 经边缘态射 $\varphi:G(K_S)\to H^1(\operatorname{Spec} K_S,\Pi(1))=\Pi^S$ 作用在 $\Pi^S$ 上时，成为 $G(K_S)$-等变的（我们应当核对符号！！！）。因此平凡的 $\mu_\infty$-torsor $\ell^S_{\mathcal{E},\tilde e}/\ell^S_{\mathcal{E},e}$ 带上了 $G(K_S)$ 的非平凡作用：它按 $\chi\varphi:G(K_S)\to\mu_\infty$ 作用。

于是每个 $e\in\mathcal{E}_S$ 对应 $G(K_S)$ 在 $\tilde\phi_S$ 上的一个作用；若把 $e$ 换成 $\chi e$（$\chi\in(\Pi^{\vee})^S=\operatorname{Hom}(\Pi^S,\mu_\infty)$），则作用乘以 $\chi\varphi:G(K_S)\to\mu_\infty$。

\emph{注.} 顺便我们已证明：上边缘态射 $\varphi:G(K_S)\to H^1(\operatorname{Spec} K_S,\Pi(1))=\Pi^S$ 局部常值\footnote{另见 4.5.4。}（的确，$G(K_S)$ 作为群 ind-概形作用在 $(\operatorname{Res})_*\tilde\phi_S$ 上，故 $\varphi$ 是 ind-概形的态射，即 $\varphi$ 局部常值。证明可改述如下。不妨设 $S$ 只由一点 $x$ 组成。群 ind-概形 $G(K_x)$ 作用在 $\operatorname{Bun}_{G,x}$ 上（见 2.3.3--%
% 原书 p.128 末行与 p.129 首行是同一括注（"（see 2.3.3 –" / "2.3.4）"），故不空行。
\SourcePage{129}%
2.3.4），故它作用在 $\pi_0(\operatorname{Bun}_{G,x})=\pi_0(\operatorname{Bun}_G)$ 上。有「第一 Chern 类」映射 $c:\pi_0(\operatorname{Bun}_G)\to\Pi$。容易证明，对 $u\in\pi_0(\operatorname{Bun}_G)$、$g\in G(K_x)$，有 $c(gu)=\varphi(g)c(u)$，其中 $\varphi:G(K_x)\to H^1(K_x,\Pi(1))=\Pi$ 是上边缘态射。故 $\varphi$ 局部常值。

\bdp{4.1.8} 记 $\widetilde{G(K_S)_{\mathcal{E}}}$ 为由 $\mu_\infty$ 与元素 $\langle g,e\rangle$（$g\in G(K_S)$，$e\in\mathcal{E}_S$）生成的群，其定义关系为
\[
\begin{aligned}
\langle g_1g_2,e\rangle&=\langle g_1,e\rangle\langle g_2,e\rangle,\\
\langle g_1,\chi e\rangle&=\chi(\varphi(g))\cdot\langle g,e\rangle,\qquad \chi\in(\Pi^{\vee})^S=\operatorname{Hom}(\Pi^S,\mu_\infty),\\
\alpha\langle g,e\rangle&=\langle g,e\rangle\alpha,\qquad \alpha\in\mu_\infty.
\end{aligned}
\]
% [Tutor Note: 原书 p.129 第二条关系左端两处下标印作 g_1 与 g，不一致（左为 ⟨g_1,χe⟩，右为 χ(φ(g))·⟨g,e⟩）；疑右端下标应为 g_1，底本如此，保留原样。]

$\widetilde{G(K_S)_{\mathcal{E}}}$ 是 $G(K_S)$ 被 $\mu_\infty$ 的中心扩张。该扩张平凡：取定 $e\in\mathcal{E}_S$ 便给出分裂
\begin{equation}\tag{154}
\sigma_e:G(K_S)\to\widetilde{G(K_S)_{\mathcal{E}}},\qquad g\mapsto\langle g,e\rangle.
\end{equation}

由 4.1.7 可知，$\widetilde{G(K_S)_{\mathcal{E}}}$ 作用在 $\ell^S_{\mathcal{E}}$ 上：$\mu_\infty\subset\widetilde{G(K_S)_{\mathcal{E}}}$ 按显然方式作用，而对应于 $e\in\mathcal{E}_S$ 的 $G(K_S)$ 在 $\ell^S_{\mathcal{E}}$ 上的作用（见 4.1.7）来自分裂 (154)。

\bdp{4.1.9} 考虑 $\operatorname{Bun}_{G,\underline{S}}$ 中对应于 $X$ 上平凡 $G$-丛（连同其在 $\underline{S}$ 上的显然平凡化）的点。让 $G(K_S)$ 作用于此点，得到态射 $f:G(K_S)\to\operatorname{Bun}_{G,\underline{S}}$。设 $G$ 半单。则 $f$ 诱导同构
\begin{equation}\tag{155}
G(K_S)/G(A_S)\xrightarrow{\ \sim\ }\operatorname{Bun}_{G,\underline{S}},
\end{equation}
其中 $A_S:=H^0(X\setminus S,\mathcal{O}_X)$（见 [La-So] 的定理 1.3，以及该书 §3 中的证明）。关键在于把 $G(K_S)$ 与 $G(A_S)$ 看作群 ind-概形，并把 $G(K_S)/G(A_S)$ 看作 fppf 商；故 (155) 不止是 $\mathbb{C}$-点集合之间的双射。还有同构
\begin{equation}\tag{156}
G(O_S)\backslash G(K_S)/G(A_S)\xrightarrow{\ \sim\ }\operatorname{Bun}_G.
\end{equation}

\SourcePage{130}

容易看出，4.1.4 与 4.1.7 中定义的 $\mu_\infty$-torsor $\ell_{\mathcal{E}}$ 与 $\ell^S_{\mathcal{E}}$ 可描述为
\begin{equation}\tag{157}
\ell^S_{\mathcal{E}}=\widetilde{G(K_S)_{\mathcal{E}}}/G(A_S),
\end{equation}
\begin{equation}\tag{158}
\ell_{\mathcal{E}}=G(O_S)\backslash\widetilde{G(K_S)_{\mathcal{E}}}/G(A_S),
\end{equation}
\looseness=-1 其中 $\widetilde{G(K_S)_{\mathcal{E}}}$ 是 4.1.8 中的中心扩张。这里嵌入 $i:G(O)\to\widetilde{G(K_S)_{\mathcal{E}}}$ 与 $j:G(A_S)\to\widetilde{G(K_S)_{\mathcal{E}}}$ 由下式定义
\begin{equation}\tag{159}
i(g)=\langle g,e\rangle,\qquad e\in\mathcal{E}_S,
\end{equation}
\begin{equation}\tag{160}
j(g)=\langle g,e\rangle\cdot(\psi(g),cl(\mathcal{E},e))_P^{-1},\qquad e\in\mathcal{E}_S
\end{equation}
% [Tutor Note: 原书 p.130 此处印作 $i:G(O)\to\cdots$，与同一段 (158) 中的 $G(O_S)$ 不一致；疑为 $G(O_S)$ 之误，底本如此，保留原样。]
（我们应当核对符号！！！）其中 $\psi$ 是边缘态射 $G(A_S)\to H^1(X\setminus S,\Pi(1))$，$cl(\mathcal{E},e)\in H^1_c(X\setminus S,\Pi^{\vee})$ 是 $(\mathcal{E},e)$ 的类（(159) 与 (160) 的右端不依赖 $e$）。

\emph{注.} 态射 $\varphi:G(K_S)\to\Pi^S$ 与 $\psi:G(A_S)\to H^1(X\setminus S,\Pi(1))$ 诱导态射
\begin{equation}\tag{161}
\operatorname{Bun}_G=G(O_S)\backslash G(K_S)/G(A_S)\to\Pi^S/H^1(X\setminus S,\Pi(1)),
\end{equation}
其中 (161) 的右端理解为商栈。显然 $\ell_{\mathcal{E}}$ 是栈 $\Pi^S/H^1(X\setminus S,\Pi(1))$ 上某个 $\mu_\infty$-torsor 的拉回。

\bdp{4.1.10} 读者可以略过 4.1 的剩余部分。

设 $C$ 是群胚。记 $\underline{C}$ 为带 fppf 拓扑的 $\mathbb{C}$-概形范畴上对应的常群胚层。若 $C$ 中对象的自同构群有限，则 $\underline{C}$ 是代数栈。记号滥用之下，我们常把 $\underline{C}$ 写成 $C$（例如 $C$ 是集合时，$\underline{C}=C\times\operatorname{Spec}\mathbb{C}$ 常与 $C$ 等同）。

\emph{例.}
\begin{enumerate}[label=\arabic*),leftmargin=2.2em]
\item 若 $C$ 只有一个对象、$G$ 是其自同构群，则 $\underline{C}$ 是 $G$ 的分类栈。
\SourcePage{131}
\item 若 $C=P(K^{\cdot})$（见 4.1.5），则 $\underline{C}$ 是 $K^0$ 关于 $K^{-1}$ 作用的商栈。因此按 4.1.5，(161) 的右端是对应于群胚 $\Pi(1)\operatorname{gerbes}(X)$ 的栈。
\item 若 $C=A\operatorname{gerbes}(X)$，则 $\underline{C}$ 是与预层 $S\mapsto A\operatorname{gerbes}(X\times S)$ 相联系的群胚层。
\end{enumerate}

把群胚 $\Pi(1)\operatorname{gerbes}(X)$ 看作代数栈。4.1.2 中我们定义了典范态射
\begin{equation}\tag{162}
\tilde c:\operatorname{Bun}_G\to\Pi(1)\operatorname{gerbes}(X),
\end{equation}
它把 $G$-丛 $\mathfrak{F}$ 对应到 $\mathfrak{F}$ 的 $\tilde G$-提升构成的 $\Pi(1)$-gerbe（顺便说，对半单 $G$ 定义的态射 (161) 与 $\tilde c$ 一致）。$\tilde c$ 是 Chern 类映射 $c:\operatorname{Bun}_G\to H^2(X,\Pi(1))=\Pi$ 的精细化；更确切地说，$c$ 是 $\tilde c$ 与典范态射 $\Pi(1)\operatorname{gerbes}(X)\to H^2(X,\Pi(1))$ 的复合，其中 $H^2(X,\Pi(1))=\Pi$ 就是 $\Pi(1)\operatorname{gerbes}(X)$ 的同构类集合。

4.1.4 中构造的 $\operatorname{Bun}_G$ 上的 $\mu_\infty$-torsor，来自 $\Pi(1)\operatorname{gerbes}(X)$ 上的 $\mu_\infty$-torsor。下述命题表明：若 $\tilde G$ 是 $G$ 的泛覆叠，则 $\operatorname{Bun}_G$ 上任一局部系统都来自 $\Pi(1)\operatorname{gerbes}(X)$ 上唯一的局部系统。

\bdp{4.1.11} \emph{命题.} 设 $\tilde G$ 是 $G$ 的泛覆叠（故 $\Pi=\pi_1(G)$）。则态射 (162) 诱导 $\operatorname{Bun}_G$ 的基本群胚与 $\Pi(1)\operatorname{gerbes}(X)$ 之间的等价。

我们勾画一个超越的证明（既然是超越的，就不再区分 $\Pi$ 与 $\Pi(1)$）。记 $X^{\mathrm{top}}$ 为 $X$ 对应的 $C^\infty$ 流形；对 $X$ 上的 $G$-丛 $\mathfrak{F}$，记 $\mathfrak{F}^{\mathrm{top}}$ 为 $X^{\mathrm{top}}$ 上对应的 $G$-丛。考虑群胚 $\operatorname{Bun}_G^{\mathrm{top}}$，其对象是 $X^{\mathrm{top}}$ 上的 $G$-丛，其态射是 $G$-丛之间 $C^\infty$ 同构的同痕类。容易证明，自然函子 $\operatorname{Bun}_G^{\mathrm{top}}\to\Pi\operatorname{gerbes}(X^{\mathrm{top}})=\Pi\operatorname{gerbes}(X)$ 是等价。于是须证明：对 $X^{\mathrm{top}}$ 上的 $G$-丛 $\xi$，由 $X$ 上配备……的 $G$-丛 $\mathfrak{F}$ 构成的栈
% [Tutor Note: 原书 p.131 末行止于 "equipped"；该句的其余部分（"with an isotopy class F^top ≅ ξ, is non-empty, connected and simply connected"）在 p.132，见下一单元首行。译文在此以「……」标记断点。]
% ===================== end 源码/units/ch04_u01.tex =====================

% ===================== begin 源码/units/ch04_u02.tex =====================
% ch04_u02 —— 原书 p.132–p.141（4.1 末、4.2.1–4.2.11）
% 首句承接 ch04_u01 的末句（p.131 末，ch04_u01 以「由 $X$ 上配备……的 $G$-丛 $\mathfrak{F}$ 构成的栈」收尾）；
% 本单元以括注补出 p.132 的省略数据，再接谓语，使拼装后的句子通顺。
\SourcePage{132}
（这里「配备」指的是同构 $\mathcal{F}^{\mathrm{top}}\xrightarrow{\ \sim\ }\xi$ 的同痕类），它非空、连通且单连通的。若把 $X$ 上的 $G$-丛解释为 $X^{\mathrm{top}}$ 上带 $\bar\partial$-联络的 $G$-丛，这一点就很清楚。

\emph{注.} 在 4.1.2 中，我们定义了 $H^{1}(X\setminus S,\Pi(1))$-torsor（主齐性空间）$\phi_{S}\to\operatorname{Bun}_{G}$。若对某个 $x\in X$ 有 $S=\{x\}$，则 $H^{1}(X\setminus S,\Pi(1))=H^{1}(X,\Pi(1))$。故 $\phi_{\{x\}}\to\operatorname{Bun}_{G}$ 是一个 $H^{1}(X,\Pi(1))$-torsor。命题 4.1.11 可以改述如下：若 $\tilde{G}$ 是 $G$ 的泛覆叠，则 Chern 类映射 $\pi_{0}(\operatorname{Bun}_{G})\to\Pi$ 是双射，而 $\phi_{\{x\}}\to\operatorname{Bun}_{G}$ 限制到每个连通分支上是泛覆叠。这确实是改述，因为 $x$ 的选取定义了一个等价。

\begin{equation}\tag{163}
\Pi(1)\text{ gerbes}(X)\xrightarrow{\ \sim\ }\Pi\times H^{1}(X,\Pi(1))\text{ tors}
\end{equation}
（对 $X$ 上的一个 $\Pi(1)$-gerbe，对应它在 $H^{2}(X,\Pi(1))=\Pi$ 中的类，以及它的对象在 $X\setminus\{x\}$ 上的同构类所构成的 $H^{1}(X,\Pi(1))$-torsor）。

\bdsec{4.2. Pfaffian（一）}

本小节我们假定：对 $(\mathbb{Z}/2\mathbb{Z})$-分次向量空间 $A$ 与 $B$，$A\otimes B$ 与 $B\otimes A$ 的等同由 $a\otimes b\mapsto(-1)^{p(a)p(b)}b\otimes a$ 定义，其中 $p(a)$ 是 $a$ 的奇偶性。仿照 [Kn-Mu]，对维数 $n<\infty$ 的向量空间 $V$，我们把 $\det V$ 看作次数为 $n\bmod 2$ 的 $(\mathbb{Z}/2\mathbb{Z})$-分次空间。

\bdp{4.2.1} 设 $X$ 是 $\mathbb{C}$ 上的光滑完备曲线。$X$ 上的 \emph{$\omega$-正交丛}是指带非退化对称配对 $\mathcal{Q}\otimes\mathcal{Q}\to\omega_{X}$ 的向量丛 $\mathcal{Q}$。记 $\omega\text{-Ort}$ 为 $X$ 上 $\omega$-正交丛的栈。在 $\omega\text{-Ort}$ 上有一个熟知的线丛 $\det R\Gamma$（它在 $\mathcal{Q}$ 上的纤维是 $\det R\Gamma(X,\mathcal{Q})$）。Laszlo 与 Sorger [La-So] 在 $\omega\text{-Ort}$ 上构造了一个 $(\mathbb{Z}/2\mathbb{Z})$-分次线丛（他们称之为 \emph{Pfaffian}），并证明 Pfaffian 的张量平方是 $\det R\Gamma$。为我们的目的，使用 Pfaffian 的另一个定义更方便。它当然应该与 [La-So] 中的定义等价，但我们没有验证这一点。

\SourcePage{133}
我们将构造 $\omega\text{-Ort}$ 上的一个线丛 $\operatorname{Pf}$，称之为 \emph{Pfaffian}。它在 $\omega$-正交丛 $\mathcal{Q}$ 上的纤维记作 $\operatorname{Pf}(\mathcal{Q})$。$-1\in\operatorname{Aut}\mathcal{Q}$ 在 $\operatorname{Pf}(\mathcal{Q})$ 上的作用定义了 $\operatorname{Pf}$ 上的一个 $(\mathbb{Z}/2\mathbb{Z})$-分次。由于 $\operatorname{Pf}$ 是线丛，「分次」只意味着存在局部常值函数 $p:(\omega\text{-Ort})\to\mathbb{Z}/2\mathbb{Z}$，使得 $\operatorname{Pf}(\mathcal{Q})$ 的次数为 $p(\mathcal{Q})$。实际上 $p(\mathcal{Q})=\dim H^{0}(\mathcal{Q})\bmod 2$。$\dim H^{0}(\mathcal{Q})\bmod 2$ 局部常值这一事实由 M. Atiyah 与 D. Mumford [At, Mu] 证明。

对 $\omega$-正交丛 $\mathcal{Q}$，记 $\mathcal{Q}^{-}$ 为同一个丛 $\mathcal{Q}$ 配备相反的配对 $\mathcal{Q}\otimes\mathcal{Q}\to\omega_{X}$。置 $\operatorname{Pf}^{-}(\mathcal{Q}):=\operatorname{Pf}(\mathcal{Q}^{-})$。我们将定义一个典范同构 $\operatorname{Pf}\otimes\operatorname{Pf}^{-}\xrightarrow{\ \sim\ }\det R\Gamma$。定义同构 $f_{\pm i}:\operatorname{Pf}(\mathcal{Q})\xrightarrow{\ \sim\ }\operatorname{Pf}(\mathcal{Q}^{-})$ 为 $f_{\pm i}:=(\varphi_{\pm i})_{*}$，其中 $i=\sqrt{-1}$，而 $\varphi_{i}:\mathcal{Q}\xrightarrow{\ \sim\ }\mathcal{Q}^{-}$ 是乘以 $i$。借助 $f_{\pm i}$ 把 $\operatorname{Pf}$ 与 $\operatorname{Pf}^{-}$ 等同起来，我们得到同构 $c_{\pm i}:\operatorname{Pf}^{\otimes 2}\xrightarrow{\ \sim\ }\det R\Gamma$，使得 $(c_{i})^{-1}c_{-i}:\operatorname{Pf}(\mathcal{Q})^{\otimes 2}\xrightarrow{\ \sim\ }\operatorname{Pf}(\mathcal{Q})^{\otimes 2}$ 是乘以 $(-1)^{p(\mathcal{Q})}$。

\emph{注.}
\begin{enumerate}[label=(\roman*),leftmargin=2.2em]
\item 若 $\mathcal{Q}$ 是 $\omega$-正交丛，则由 Serre 对偶 $H^{1}(X,\mathcal{Q})=(H^{0}(X,\mathcal{Q}))^{*}$，故 $\det R\Gamma(X,\mathcal{Q})=\det H^{0}(X,\mathcal{Q})^{\otimes 2}$。朴素的定义会是 $\operatorname{Pf}^{?}(\mathcal{Q}):=\det H^{0}(X,\mathcal{Q})$，但它对 $\mathcal{Q}$ 的族没有意义，因为 $\dim H^{0}(X,\mathcal{Q})$ 可能跳跃。
\item 设 $\mathcal{Q}$ 是 $\mathcal{Q}_{1}$ 与 $\mathcal{Q}_{2}$ 的正交直和。则 $\det R\Gamma(X,\mathcal{Q})=\det R\Gamma(X,\mathcal{Q}_{1})\otimes\det R\Gamma(X,\mathcal{Q}_{2})$。由 $\operatorname{Pf}$ 与 $\operatorname{Pf}\otimes\operatorname{Pf}^{-}\xrightarrow{\sim}\det R\Gamma$ 的定义将会清楚：存在典范同构 $\operatorname{Pf}(\mathcal{Q})\xrightarrow{\ \sim\ }\operatorname{Pf}(\mathcal{Q}_{1})\otimes\operatorname{Pf}(\mathcal{Q}_{2})$，并且图表
\[
\begin{tikzcd}[column sep=3.2em, row sep=2.6em]
  \operatorname{Pf}(\mathcal{Q})\otimes\operatorname{Pf}(\mathcal{Q}^{-})
    \arrow[r, "\sim"]
    \arrow[d, "\wr"']
  & \operatorname{Pf}(\mathcal{Q}_{1})\otimes\operatorname{Pf}(\mathcal{Q}_{1}^{-})\otimes\operatorname{Pf}(\mathcal{Q}_{2})\otimes\operatorname{Pf}(\mathcal{Q}_{2}^{-})
    \arrow[d, "\wr"] \\
  \det R\Gamma(X,\mathcal{Q})
    \arrow[r, "\sim"]
  & \det R\Gamma(X,\mathcal{Q}_{1})\otimes\det R\Gamma(X,\mathcal{Q}_{2})
\end{tikzcd}
\]
交换。因此同构 $c_{\pm i}:\operatorname{Pf}(\mathcal{Q})^{\otimes 2}\xrightarrow{\ \sim\ }\det R\Gamma(X,\mathcal{Q})$ 与分解 $\mathcal{Q}=\mathcal{Q}_{1}\oplus\mathcal{Q}_{2}$ 相容。
\SourcePage{134}
\item 可以定义 $c_{\pm}:\operatorname{Pf}(\mathcal{Q})^{\otimes 2}\xrightarrow{\ \sim\ }\det R\Gamma(X,\mathcal{Q})$ 为 $c_{\pm}=i^{\pm p(\mathcal{Q})^{2}}c_{i}$，其中 $p(\mathcal{Q})^{2}$ 看作 $\mathbb{Z}/4\mathbb{Z}$ 的元素。于是把 $i$ 换成 $-i$ 时 $c_{\pm}$ 不变。然而 $c_{\pm}$ 似乎不是自然的对象，例如它们与分解 $\mathcal{Q}=\mathcal{Q}_{1}\oplus\mathcal{Q}_{2}$ 不相容（「误差」是 $(-1)^{p(\mathcal{Q}_{1})p(\mathcal{Q}_{2})}$）。
\item 若把 $\mathbb{C}$ 换成满足 $\operatorname{char}k\ne 2$ 的任意域 $k$，$\operatorname{Pf}(\mathcal{Q})$ 的构造仍然有效。$\operatorname{char}k=2$ 的情形在 4.2.16 中讨论。
\end{enumerate}

\bdp{4.2.2} \emph{Lagrange 三元组}由以下数据构成：偶维向量空间 $V$，$V$ 上的非退化双线性对称型 $(\ ,\ )$，以及 Lagrange（= 极大各向同性）子空间 $L_{+},L_{-}\subset V$。若 $X$ 与 $\mathcal{Q}$ 如 4.2.1，而 $\mathcal{Q}'\subset\mathcal{Q}$ 是满足 $H^{0}(X,\mathcal{Q}')=0$ 且 $S:=\operatorname{Supp}(\mathcal{Q}/\mathcal{Q}')$ 有限的子层，则对 $(\mathcal{Q},\mathcal{Q}')$ 如下对应一个 Lagrange 三元组 $(V;L_{+},L_{-})$（参 [Mu]）：
\begin{enumerate}[label=(\arabic*),leftmargin=2.2em]
\item $V:=H^{0}(X,\mathcal{Q}''/\mathcal{Q}')$，其中 $\mathcal{Q}'':=\operatorname{Hom}(\mathcal{Q}',\omega_{X})\supset\mathcal{Q}$；
\item $L_{+}:=H^{0}(X,\mathcal{Q}/\mathcal{Q}')\subset V$；
\item $L_{-}:=H^{0}(X,\mathcal{Q}'')\subset V$；
\item $V$ 上的双线性型由下述两者诱导：自然配对 $\mathcal{Q}''/\mathcal{Q}'\otimes\mathcal{Q}''/\mathcal{Q}'\to(j_{*}\omega_{X\setminus S})/\omega_{X}$，以及「留数之和」映射 $H^{0}(X,(j_{*}\omega_{X\setminus S})/\omega_{X})\to\mathbb{C}$。这里 $j$ 是嵌入 $X\setminus S\to X$。在此情形下，可以把 $R\Gamma(X,\mathcal{Q})$ 与复形
\begin{equation}\tag{164}
0\longrightarrow L_{-}\longrightarrow V/L_{+}\longrightarrow 0
\end{equation}
等同起来，该复形集中位于次数 0 与 1。特别地 $H^{0}(X,\mathcal{Q})=L_{+}\cap L_{-}$，$H^{1}(X,\mathcal{Q})=V/(L_{+}+L_{-})$。而 $H^{0}(X,\mathcal{Q})=L_{+}\cap L_{-}$ 与 $H^{1}(X,\mathcal{Q})=V/(L_{+}+L_{-})$ 之间的 Serre 配对由 $V$ 上的双线性型诱导。
\end{enumerate}

\bdp{4.2.3} 对 Lagrange 三元组 $(V;L_{+},L_{-})$，置
\begin{equation}\tag{165}
\det(V;L_{+},L_{-}):=\det L_{+}\otimes\det L_{-}\otimes(\det V)^{*}.
\end{equation}

\SourcePage{135}
$\det(V;L_{+},L_{-})$ 正是复形 (164) 的行列式。公式 (165) 在 Lagrange 三元组的栈上定义了一个线丛 $\det$。在 4.2.4 与 4.2.8 中，我们将在此栈上构造一个 $(\mathbb{Z}/2\mathbb{Z})$-分次线丛 $\operatorname{Pf}$，以及一个典范同构 $\operatorname{Pf}\otimes\operatorname{Pf}^{-}\xrightarrow{\ \sim\ }\det$。这里 $\operatorname{Pf}^{-}(V;L_{+},L_{-}):=\operatorname{Pf}(V^{-};L_{+},L_{-})$，而 $V^{-}$ 表示 $V$ 配备型 $-(\ ,\ )$。朴素的「定义」会是 $\operatorname{Pf}^{?}(V;L_{+},L_{-}):=\det(L_{+}\cap L_{-})$，或 $\operatorname{Pf}^{?}(V;L_{+},L_{-})^{*}:=\det((L_{+}\cap L_{-})^{*})=\det(V/(L_{+}+L_{-}))$（参 4.2.1 的注 (i)）。

\bdp{4.2.4} 对 Lagrange 三元组 $(V;L_{+},L_{-})$，如下定义 $\operatorname{Pf}(V;L_{+},L_{-})$。记 $\operatorname{Cl}(V)$ 为带典范 $(\mathbb{Z}/2\mathbb{Z})$-分次的 Clifford 代数（$V\subset\operatorname{Cl}(V)$ 为奇）。设 $M$ 是不可约 $(\mathbb{Z}/2\mathbb{Z})$-分次 $\operatorname{Cl}(V)$-模（实际上，即使不计分次，$M$ 也不可约）。$M$ 在张量一维 $(\mathbb{Z}/2\mathbb{Z})$-分次向量空间的意义下唯一。置 $M_{L_{-}}=M/L_{-}M$，$M^{L_{+}}:=\{m\in M\mid L_{+}m=0\}$。则 $M^{L_{+}}$ 与 $M_{L_{-}}$ 是一维 $(\mathbb{Z}/2\mathbb{Z})$-分次空间。我们置
\begin{equation}\tag{166}
\operatorname{Pf}(V;L_{+},L_{-}):=M^{L_{+}}\otimes(M_{L_{-}})^{*}.
\end{equation}
特别地，可以取 $M=\operatorname{Cl}(V)/\operatorname{Cl}(V)L_{+}$。此时 $M^{L_{+}}=\mathbb{C}$，故
\begin{equation}\tag{167}
\operatorname{Pf}(V;L_{+},L_{-})^{*}=\operatorname{Cl}(V)/(L_{-}\cdot\operatorname{Cl}(V)+\operatorname{Cl}(V)\cdot L_{+}).
\end{equation}
显然 (166) 或 (167) 把 $\operatorname{Pf}$ 定义为 Lagrange 三元组栈上的 $(\mathbb{Z}/2\mathbb{Z})$-分次线丛。\footnote{换言之，从个别 Lagrange 三元组过渡到族是显然的。此原则贯穿我们对 Pfaffian 的全部讨论（只有在 4.2.14 的无穷维情形中，我们才显式考虑族，因为这确实需要一些小心）。} 该分次对应于 $-1\in\operatorname{Aut}(V;L_{+},L_{-})$ 在 $\operatorname{Pf}(V;L_{+},L_{-})$ 上的作用。

若 $V$ 是 $V_{1}$ 与 $V_{2}$ 的正交直和，则 $\operatorname{Cl}(V)$ 是超代数 $\operatorname{Cl}(V_{1})$ 与 $\operatorname{Cl}(V_{2})$ 的张量积。因此，若 $(V^{1};L^{1}_{+},L^{1}_{-})$ 与 $(V^{2};L^{2}_{+},L^{2}_{-})$ 是 Lagrange 三元组，则有典范同构
\begin{equation}\tag{168}
\begin{aligned}
&\operatorname{Pf}(V^{1}\oplus V^{2};L^{1}_{+}\oplus L^{2}_{+},L^{1}_{-}\oplus L^{2}_{-})\\
&\qquad=\operatorname{Pf}(V^{1};L^{1}_{+},L^{1}_{-})\otimes\operatorname{Pf}(V^{2};L^{2}_{+},L^{2}_{-}),
\end{aligned}
\end{equation}
其中 $\oplus$ 表示正交直和。

\SourcePage{136}
$\operatorname{Pf}(V;L_{+},L_{-})$ 为偶，当且仅当 $\dim(L_{+}\cap L_{-})$ 为偶。这由 (168) 与下面引理的陈述 (i) 推出。

\bdp{4.2.5} \emph{引理.}
\begin{enumerate}[label=(\roman*),leftmargin=2.2em]
\item 任意 Lagrange 三元组 $(V;L_{+},L_{-})$ 都可表为 Lagrange 三元组 $(V^{1};L^{1}_{+},L^{1}_{-})$ 与 $(V^{2};L^{2}_{+},L^{2}_{-})$ 的正交直和，使得 $L^{1}_{+}\cap L^{1}_{-}=0$，$L^{2}_{+}=L^{2}_{-}$。
\item 此外，若固定子空间 $\Lambda\subset L_{+}$ 使 $L_{+}=\Lambda\oplus(L_{+}\cap L_{-})$，则可选取上述分解 $(V;L_{+},L_{-})=(V^{1};L^{1}_{+},L^{1}_{-})\oplus(V^{2};L^{2}_{+},L^{2}_{-})$，使得 $L^{1}_{+}=\Lambda$。
\end{enumerate}

\emph{证明.}
\begin{enumerate}[label=(\roman*),leftmargin=2.2em]
\item 选取子空间 $P\subset V$ 使 $V=(L_{+}+L_{-})\oplus P$。然后置 $V^{2}:=(L_{1}\cap L_{2})\oplus P$，$V^{1}:=(V^{2})^{\perp}$。
% [Tutor Note: 原书 p.136 证明 (i) 中印作 $L_1\cap L_2$，按上下文应为 $L_+\cap L_-$；底本如此，保留原样。]
\item 选取子空间 $P\subset\Lambda^{\perp}$ 使 $\Lambda^{\perp}=L_{+}\oplus P$（这意味着 $V=(L_{+}+L_{-})\oplus P$，因为 $\Lambda^{\perp}/L_{+}\to V/(L_{+}+L_{-})$ 是同构）。然后如上进行。\hfill$\Box$
\end{enumerate}

\bdp{4.2.6} 本小节（读者可以跳过）我们在 $\operatorname{Pf}(V;L_{+},L_{-})$ 与 4.2.3 的朴素 $\operatorname{Pf}^{?}(V;L_{+},L_{-})$ 之间构造一个典范同构。回忆 $\operatorname{Pf}^{?}(V;L_{+},L_{-}):=\det(L_{+}\cap L_{-})$，故 $\operatorname{Pf}^{?}(V;L_{+},L_{-})^{*}=\det((L_{+}\cap L_{-})^{*})=\det(V/(L_{+}+L_{-}))$。这里应理解为：配对 $\det W\otimes\det W^{*}\to\mathbb{C}$（$W:=L_{+}\cap L_{-}$）由 $(e_{1}\wedge\dots\wedge e_{k})\otimes(e^{k}\wedge\dots\wedge e^{1})\mapsto 1$ 定义，其中 $e_{1},\dots,e_{k}$ 是 $W$ 的基，$e^{1},\dots,e^{k}$ 是 $W^{*}$ 的对偶基（从「超」的观点看，这个配对是合理的；例如它与分解 $W=W_{1}\oplus W_{2}$ 相容）。

为定义同构 $\operatorname{Pf}(V;L_{+},L_{-})\xrightarrow{\ \sim\ }\operatorname{Pf}^{?}(V;L_{+},L_{-})$，我们使用 $\operatorname{Cl}(V)$ 上由
\begin{equation}\tag{169}
\operatorname{Cl}_{0}(V)=\mathbb{C},\qquad \operatorname{Cl}_{k+1}(V)=\operatorname{Cl}_{k}(V)+V\cdot\operatorname{Cl}_{k}(V)
\end{equation}
定义的典范滤链。我们有 $\operatorname{Cl}_{k}(V)/\operatorname{Cl}_{k-1}(V)=\bigwedge^{k}V$。置 $r:=\dim(L_{+}\cap L_{-})$。有典范满射
% 下式起于 p.136 末、止于 p.137 首，跨页处只拆行不拆义。
\SourcePage{137}
\[
\begin{aligned}
\varphi:\operatorname{Cl}_{r}(V)&\longrightarrow\bigwedge\nolimits^{r}V\longrightarrow\bigwedge\nolimits^{r}(V/(L_{+}+L_{-}))\\
&=\det(V/(L_{+}+L_{-}))=\operatorname{Pf}^{?}(V;L_{+},L_{-})^{*}.
\end{aligned}
\]
由 4.2.5(i) 容易推出：典范映射 $\operatorname{Cl}_{r}(V)\to\operatorname{Cl}(V)/(L_{-}\cdot\operatorname{Cl}(V)+\operatorname{Cl}(V)\cdot L_{+})=\operatorname{Pf}(V;L_{+},L_{-})^{*}$ 经过 $\varphi$ 分解。由此诱导的映射 $f:\operatorname{Pf}^{?}(V;L_{+},L_{-})^{*}\to\operatorname{Pf}(V;L_{+},L_{-})^{*}$ 是同构。$f^{*}$ 就是所求的同构 $\operatorname{Pf}(V;L_{+},L_{-})\xrightarrow{\ \sim\ }\operatorname{Pf}^{?}(V;L_{+},L_{-})$。

这里有一个等价定义。设 $M$ 是不可约 $(\mathbb{Z}/2\mathbb{Z})$-分次 $\operatorname{Cl}(V)$-模。典范嵌入 $\det(L_{+}\cap L_{-})\subset\bigwedge^{*}(L_{+}\cap L_{-})=\operatorname{Cl}(L_{+}\cap L_{-})\subset\operatorname{Cl}(V)$ 诱导映射 $\det(L_{+}\cap L_{-})\otimes M_{L_{+}\cap L_{-}}\to M^{L_{+}\cap L_{-}}$，它实际上是同构。容易由 4.2.5(i) 推出，复合
\[
\begin{aligned}
M^{L_{+}}\longrightarrow M^{L_{+}\cap L_{-}}&\xrightarrow{\ \sim\ }\det(L_{+}\cap L_{-})\otimes M_{L_{+}\cap L_{-}}\\
&\longrightarrow\det(L_{+}\cap L_{-})\otimes M_{L_{-}}
\end{aligned}
\]
是同构。它诱导同构 $\operatorname{Pf}(V;L_{+},L_{-}):=M^{L_{+}}\otimes(M_{L_{-}})^{\otimes-1}\to\det(L_{+}\cap L_{-})=\operatorname{Pf}^{?}(V;L_{+},L_{-})$，这实际上与上面构造的同构互逆。

\bdp{4.2.7} 在构造同构 $\operatorname{Pf}\otimes\operatorname{Pf}^{-}\xrightarrow{\ \sim\ }\det$ 之前，我们先构造一个典范同构
\begin{equation}\tag{170}
\begin{aligned}
&\operatorname{Pf}(V\oplus V^{*};L_{+}\oplus L_{+}^{\perp},L_{-}\oplus L_{-}^{\perp})\\
&\qquad\xrightarrow{\ \sim\ }\det(V;L_{+},L_{-}),
\end{aligned}
\end{equation}
其中 $V$ 是不带任何双线性型的有限维向量空间，$L_{\pm}\subset V$ 是任意子空间，而 $V\oplus V^{*}$ 配备显然的双线性型（(170) 的左端有意义，因为 $L_{\pm}\oplus L_{\pm}^{\perp}$ 是 Lagrange 的；(170) 的右端由 (165) 定义）。置
\begin{equation}\tag{171}
M=\bigwedge V\otimes(\det L_{+})^{*},\qquad \bigwedge V:=\bigoplus_{i}\bigwedge\nolimits^{i}V.
\end{equation}
$M$ 是不可约 $\operatorname{Cl}(V\oplus V^{*})$-模，且 $M^{L_{+}\oplus L_{+}^{\perp}}=\mathbb{C}$，故按 (166)，$\operatorname{Pf}(V\oplus V^{*};L_{+}\oplus L_{+}^{\perp},L_{-}\oplus L_{-}^{\perp})=(M_{L_{-}\oplus L_{-}^{\perp}})^{*}$。显然 $M_{L_{-}}=\bigwedge(V/L_{-})\otimes(\det L_{+})^{*}$，且 $M_{L_{-}\oplus L_{-}^{\perp}}=\det(V/L_{-})\otimes(\det L_{+})^{*}=\det(V;L_{+},L_{-})^{*}$（见 (165)）。于是我们构造了同构 (170)。

\SourcePage{138}
\bdp{4.2.8} 现在设 $(V;L_{+},L_{-})$ 是 Lagrange 三元组。我们将构造一个典范同构
\begin{equation}\tag{172}
\begin{aligned}
&\operatorname{Pf}(V;L_{+},L_{-})\otimes\operatorname{Pf}(V^{-};L_{+},L_{-})\\
&\qquad\xrightarrow{\ \sim\ }\det(V;L_{+},L_{-}),
\end{aligned}
\end{equation}
其中 $V^{-}$ 表示 $V$ 配备双线性型 $-(\ ,\ )$。若 $W$ 是配备非退化对称双线性型的有限维向量空间，则 $(V\otimes W;L_{+}\otimes W,L_{-}\otimes W)$ 是 Lagrange 三元组。(170) 可以改写为典范同构
\begin{equation}\tag{173}
\det(V;L_{+},L_{-})\xrightarrow{\ \sim\ }\operatorname{Pf}(V\otimes H;L_{+}\otimes H,L_{-}\otimes H),
\end{equation}
其中 $H$ 表示 $\mathbb{C}^{2}$ 配备双线性型 $\begin{pmatrix}0&1\\1&0\end{pmatrix}$。另一方面，(168) 给出同构
\begin{equation}\tag{174}
\begin{aligned}
&\operatorname{Pf}(V;L_{+},L_{-})\otimes\operatorname{Pf}(V^{-};L_{+},L_{-})\\
&\qquad\xrightarrow{\ \sim\ }\operatorname{Pf}(V\otimes H';L_{+}\otimes H',L_{-}\otimes H'),
\end{aligned}
\end{equation}
其中 $H'$ 表示 $\mathbb{C}^{2}$ 配备双线性型 $\begin{pmatrix}1&0\\0&-1\end{pmatrix}$。于是同构 $\varphi:H'\xrightarrow{\ \sim\ }H$ 诱导同构
\[
\varphi_{*}:\operatorname{Pf}(V;L_{+},L_{-})\otimes\operatorname{Pf}(V^{-};L_{+},L_{-})\xrightarrow{\ \sim\ }\det(V;L_{+},L_{-}).
\]
\emph{引理.} 若 $\psi\in\operatorname{Aut}H'$，则
\begin{equation}\tag{175}
(\varphi\psi)_{*}=(\det\psi)^{n}\varphi_{*},\qquad n=\dim(L_{+}\cap L_{-}).
\end{equation}
\emph{证明.} $\operatorname{Aut}H'$ 通过某个特征标 $\chi:\operatorname{Aut}H'\to\mathbb{C}^{*}$ 作用在 (174) 的右端。$\operatorname{Aut}H'$ 的任意特征标都具有形式 $\psi\mapsto(\det\psi)^{m}$，$m\in\mathbb{Z}/2\mathbb{Z}$。$\chi\begin{pmatrix}1&0\\0&-1\end{pmatrix}=(-1)^{n}$，$n:=\dim(L_{+}\cap L_{-})$，因为 $-1\in\operatorname{Aut}(V;L_{+},L_{-})$ 在 $\operatorname{Pf}(V;L_{+},L_{-})$ 上的作用是 $(-1)^{n}$（见 4.2.4）。故 $m=n\bmod 2$。\hfill$\Box$

我们如下定义 (172)：取任意满足 $\det\varphi=1$ 的同构 $\varphi:H'\xrightarrow{\ \sim\ }H$，令 (172) 为 $\varphi_{*}$。

\emph{注.}
\begin{enumerate}[label=(\roman*),leftmargin=2.2em]
\item (172) 与 $(V;L_{+},L_{-})$ 到正交直和的分解相容；也就是说，若有一个这样的分解 $(V;L_{+},L_{-})=(V^{1};L^{1}_{+},L^{1}_{-})\oplus(V^{2};L^{2}_{+},L^{2}_{-})$，则同构\SourcePage{139}
(172) 对 $(V;L_{+},L_{-})$、$(V^{1};L^{1}_{+},L^{1}_{-})$ 与 $(V^{2};L^{2}_{+},L^{2}_{-})$ 这三个三元组，都与 (168) 及典范同构 $\det(V;L_{+},L_{-})=\det(V^{1};L^{1}_{+},L^{1}_{-})\otimes\det(V^{2};L^{2}_{+},L^{2}_{-})$ 相容。
\item (170) 与 $(V;L_{+},L_{-})$ 到直和的分解相容。
\end{enumerate}

\bdp{4.2.9} 本小节（读者可以跳过）给出 (172) 的一个等价构造。我们将使用超代数反同构 $*:\operatorname{Cl}(V^{-})\xrightarrow{\ \sim\ }\operatorname{Cl}(V)$，它在 $V$ 上恒等（对任意 $v_{1},\dots,v_{k}\in V$，有 $(v_{1}\cdots v_{k})^{*}=(-1)^{k(k-1)/2}v_{k}\cdots v_{1}$）。我们还使用典范映射 $\operatorname{sTr}:\operatorname{Cl}(V)=\operatorname{Cl}_{n}(V)\to\operatorname{Cl}_{n}(V)/\operatorname{Cl}_{n-1}(V)=\det V$，其中 $n=\dim V$，而 $\operatorname{Cl}_{k}(V)$ 由 (169) 定义。它具有「超迹性质」
\begin{equation}\tag{176}
\operatorname{sTr}(ab)=(-1)^{p(a)p(b)}\operatorname{sTr}(ba),
\end{equation}
其中 $a,b\in\operatorname{Cl}(V)$ 是次数为 $p(a),p(b)\in\mathbb{Z}/2\mathbb{Z}$ 的齐次元素。事实上，只需在 $a\in V$、$p(ab)=n\bmod 2$ 的情形证明 (176)。此时 $b\in\operatorname{Cl}_{n-1}(V)$，(176) 显然。或者可以验证：$\operatorname{sTr}(a)$ 与算子 $a:M\to M$ 的超迹至多相差一个符号，其中 $M$ 是不可约 $\operatorname{Cl}(V)$-模。

现在考虑映射
\begin{equation}\tag{177}
\begin{aligned}
&\det L_{-}\otimes\operatorname{Pf}(V;L_{+},L_{-})^{*}\otimes\det L_{+}\otimes\operatorname{Pf}(V^{-};L_{+},L_{-})^{*}\\
&\qquad\longrightarrow\det V,
\end{aligned}
\end{equation}
它由 $a_{-}\otimes x\otimes a_{+}\otimes y\mapsto\operatorname{sTr}(a_{-}xa_{+}y^{*})$ 定义。这里 $a_{\pm}\in\det L_{\pm}\subset\bigwedge^{*}(L_{\pm})=\operatorname{Cl}(L_{\pm})\subset\operatorname{Cl}(V)$，$x\in\operatorname{Pf}(V;L_{+},L_{-})^{*}=\operatorname{Cl}(V)/(L_{-}\cdot\operatorname{Cl}(V)+\operatorname{Cl}(V)\cdot L_{+})$，$y^{*}\in\operatorname{Cl}(V)/(L_{+}\cdot\operatorname{Cl}(V)+\operatorname{Cl}(V)\cdot L_{-})$，故 (177) 是良定义的。容易看出（例如由 4.2.5(i)），(177) 是同构。它诱导同构
\[
\operatorname{Pf}(V;L_{+},L_{-})\otimes\operatorname{Pf}(V^{-};L_{+},L_{-})\xrightarrow{\ \sim\ }\det L_{+}\otimes\det L_{-}\otimes(\det V)^{*}=\det(V;L_{+},L_{-}).
\]
可以证明，这个同构等于 (172)。

\SourcePage{140}
\bdp{4.2.10} 设 $X$ 与 $\mathcal{Q}$ 如 4.2.1，$\mathcal{Q}'\subset\mathcal{Q}$ 如 4.2.2。对这些数据，我们已对应一个 Lagrange 三元组 $(V;L_{+},L_{-})$，使得 $\det(V;L_{+},L_{-})=\det R\Gamma(X,\mathcal{Q})$（见 4.2.2）。置 $\operatorname{Pf}_{\mathcal{Q}'}(\mathcal{Q}):=\operatorname{Pf}(V;L_{+},L_{-})$。按 4.2.9，有典范同构 $\operatorname{Pf}_{\mathcal{Q}'}(\mathcal{Q})\otimes\operatorname{Pf}_{\mathcal{Q}'}(\mathcal{Q}^{-})\xrightarrow{\ \sim\ }\det R\Gamma(X,\mathcal{Q})$。为定义 $\operatorname{Pf}(\mathcal{Q})$，只需对所有满足 $\mathcal{Q}'\subset\tilde{\mathcal{Q}}'$ 的对 $(\mathcal{Q}',\tilde{\mathcal{Q}}')$，定义相容的同构系统 $\operatorname{Pf}_{\mathcal{Q}'}(\mathcal{Q})\xrightarrow{\ \sim\ }\operatorname{Pf}_{\tilde{\mathcal{Q}}'}(\mathcal{Q})$。为定义 $\operatorname{Pf}(\mathcal{Q})\otimes\operatorname{Pf}(\mathcal{Q}^{-})\xrightarrow{\ \sim\ }\det R\Gamma(X,\mathcal{Q})$，只需证明下述图表的交换性：
% [Tutor Note: 原书 p.140 此交换图排印错乱：底行右端的对象缺失，且两支之间的纵向箭头被印成横箭头（落在顶行张量积号下方）。译文把该箭头排为纵向，底行右端照原样留空。]
\[
\begin{tikzcd}[column sep=3.4em, row sep=2.6em]
  \operatorname{Pf}_{\mathcal{Q}'}(\mathcal{Q})\otimes\operatorname{Pf}_{\mathcal{Q}'}(\mathcal{Q}^{-})
    \arrow[r, "\sim"]
    \arrow[d]
  & \det R\Gamma(X,\mathcal{Q}) \\
  \operatorname{Pf}_{\tilde{\mathcal{Q}}'}(\mathcal{Q})\otimes\operatorname{Pf}_{\tilde{\mathcal{Q}}'}(\mathcal{Q}^{-})
    \arrow[r, "\sim"]
  & \phantom{\det R\Gamma(X,\mathcal{Q})}
\end{tikzcd}
\]
对应 $\tilde{\mathcal{Q}}'$ 的 Lagrange 三元组 $(\tilde{V};\tilde{L}_{+},\tilde{L}_{-})$ 与对应 $\mathcal{Q}'$ 的三元组 $(V;L_{+},L_{-})$ 如下相关：若 $\Lambda=H^{0}(X,\tilde{\mathcal{Q}}'/\mathcal{Q}')\subset H^{0}(X,\mathcal{Q}/\mathcal{Q}')=L_{+}$，则
\begin{equation}\tag{178}
\tilde{V}=\Lambda^{\perp}/\Lambda,\qquad \tilde{L}_{+}=L_{+}/\Lambda\subset\tilde{V},\qquad \tilde{L}_{-}=L_{-}\cap\Lambda^{\perp}\hookrightarrow\tilde{V}
\end{equation}
（注意 $\Lambda\cap L_{-}=H^{0}(X,\tilde{\mathcal{Q}}')=0$）。于是剩下的是一些线性代数（见 4.2.11）。容易验证，我们对 $\operatorname{Pf}(\mathcal{Q})$ 以及 $\operatorname{Pf}(\mathcal{Q})\otimes\operatorname{Pf}(\mathcal{Q}^{-})\xrightarrow{\ \sim\ }\det R\Gamma(X,\mathcal{Q})$ 的定义对 $\mathcal{Q}$ 的族有意义。

\bdp{4.2.11} 设 $(V;L_{+},L_{-})$ 是 Lagrange 三元组，$\Lambda\subset L_{+}$ 是满足 $\Lambda\cap L_{-}=0$ 的子空间。则由 (178) 定义的 $(\tilde{V};\tilde{L}_{+},\tilde{L}_{-})$ 是 Lagrange 三元组。在此情形下，我们说 $(\tilde{V};\tilde{L}_{+},\tilde{L}_{-})$ 是 $(V;L_{+},L_{-})$ 的 \emph{子商}。容易证明，子商的子商仍是子商。于是可以考虑以 Lagrange 三元组为对象的范畴 $T$，其中从 $(V;L_{+},L_{-})$ 到 $(V';L'_{+},L'_{-})$ 的态射定义为 $(V;L_{+},L_{-})$ 与 $(V';L'_{+},L'_{-})$ 的某个子商之间的同构。再考虑范畴 $C$，其对象是有限维向量空间的有限复形，态射是拟同构。记 $\mathbb{I}$ 为如下范畴：对象是 $(\mathbb{Z}/2\mathbb{Z})$-分次一维向量空间，态射是保持\SourcePage{141}%
分次的同构。

复形 (164) 作为 $C$ 的对象，函子性地依赖于 $(V;L_{+},L_{-})\in T$：若 $(\tilde{V};\tilde{L}_{+},\tilde{L}_{-})$ 是对应于 $\Lambda\subset L_{+}$ 的 $(V;L_{+},L_{-})$ 的子商，则有拟同构
% [Tutor Note: 原书 p.141 此图把两支纵向的嵌入箭头印成横向钩箭头（分置两行之间偏左、偏右处），未与上下两行对齐；译文按图意排为纵向 hook 箭头：$\tilde{L}_-\hookrightarrow L_-$ 与 $\Lambda^\perp/L_+\hookrightarrow V/L_+$。]
\[
\begin{tikzcd}[column sep=3.6em, row sep=2.4em]
  L_{-} \arrow[r] & V/L_{+} \\
  \tilde{L}_{-} \arrow[r] \arrow[u, hook] & \tilde{V}/\tilde{L}_{+}=\Lambda^{\perp}/L_{+} \arrow[u, hook]
\end{tikzcd}
\]
应用 [Kn-Mu] 的函子 $\det:C\to\mathbb{I}$，我们看到 $\det(V;L_{+},L_{-})\in\mathbb{I}$ 函子性地依赖于 $(V;L_{+},L_{-})\in T$。若 $(\tilde{V};\tilde{L}_{+},\tilde{L}_{-})$ 是对应于 $\Lambda\subset L_{+}$ 的 $(V;L_{+},L_{-})$ 的子商，则 $\det(V;L_{+},L_{-})=(\det L_{+})\otimes(\det L_{-})\otimes(\det V)^{*}$ 与 $\det(\tilde{V};\tilde{L}_{+},\tilde{L}_{-})=(\det\tilde{L}_{+})\otimes(\det\tilde{L}_{-})\otimes(\det\tilde{V})^{*}$ 之间的同构来自自然同构 $\det L_{+}=\det\Lambda\otimes\det\tilde{L}_{+}$、$\det L_{-}=\det\tilde{L}_{-}\otimes\det(V/\Lambda^{\perp})$ 与 $\det V=\det\Lambda\otimes\det\tilde{V}\otimes\det(V/\Lambda^{\perp})$。

正如 4.2.10 所解释的，我们必须把 $\operatorname{Pf}$ 定义为函子 $T\to\mathbb{I}$，并证明 4.2.8 中的同构 $\operatorname{Pf}(V;L_{+},L_{-})\otimes\operatorname{Pf}^{-}(V;L_{+},L_{-})\xrightarrow{\ \sim\ }\det(V;L_{+},L_{-})$ 具有函子性。

若 $(\tilde{V};\tilde{L}_{+},\tilde{L}_{-})$ 是对应于 $\Lambda\subset L_{+}$ 的 $(V;L_{+},L_{-})$ 的子商，则
\[
\begin{aligned}
\operatorname{Pf}(V;L_{+},L_{-})^{*}&=\operatorname{Cl}(V)/(L_{-}\cdot\operatorname{Cl}(V)+\operatorname{Cl}(V)\cdot L_{+}),\\
\operatorname{Pf}(\tilde{V};\tilde{L}_{+},\tilde{L}_{-})^{*}
&=\operatorname{Cl}(\Lambda^{\perp})/\bigl((L_{-}\cap\Lambda^{\perp})\cdot\operatorname{Cl}(\Lambda^{\perp})\\
&\qquad+\operatorname{Cl}(\Lambda^{\perp})\cdot L_{+}\bigr).
\end{aligned}
\]
于是嵌入 $\operatorname{Cl}(\Lambda^{\perp})\to\operatorname{Cl}(V)$ 诱导映射
\begin{equation}\tag{179}
\operatorname{Pf}(\tilde{V};\tilde{L}_{+},\tilde{L}_{-})^{*}\to\operatorname{Pf}(V;L_{+},L_{-})^{*}.
\end{equation}
这定义了 $\operatorname{Pf}^{*}$ 为函子 $T\to\{(\mathbb{Z}/2\mathbb{Z})\text{-分次一维空间}\}$（容易看出，复合对应复合）。还需证明
\begin{enumerate}[label=\alph*),leftmargin=2.2em]
\item (179) 是同构，
\end{enumerate}
% ===================== end 源码/units/ch04_u02.tex =====================

% ===================== begin 源码/units/ch04_u03.tex =====================
% ch04_u03 —— 原书 p.142–p.151（4.2.11 末、4.2.12–4.2.16、4.3.1–4.3.2）
% 首句承接 ch04_u02 末的 a) 项（p.141 末），故本节自 b) 项起；用 start=2 接续原书编号。
\SourcePage{142}
\begin{enumerate}[label=\alph*),leftmargin=2.2em,start=2]
\item (179) 与来自 4.2.8 的配对
$\operatorname{Pf}(V;L_{+},L_{-})^{*}\otimes\operatorname{Pf}(V^{-};L_{+},L_{-})^{*}\xrightarrow{\ \sim\ }\det(V;L_{+},L_{-})^{*}$
及
$\operatorname{Pf}(\tilde{V};\tilde{L}_{+},\tilde{L}_{-})^{*}\otimes\operatorname{Pf}(\tilde{V}^{-};L_{+},L_{-})^{*}\xrightarrow{\ \sim\ }\det(\tilde{V};\tilde{L}_{+},\tilde{L}_{-})^{*}$
相容。
% [Tutor Note: 原书 p.142 第二个配对中，Pfaffian 自变量里的 L 缺波浪号（第一个配对中是带波浪号的 L），底本如此，保留原样。]
\end{enumerate}

b) 可直接验证，a) 由 b) 推出。也可以把 $(\tilde{V};\tilde{L}_{+},\tilde{L}_{-})$ 归到它是\emph{极大子商}的情形（即 $\Lambda\oplus(L_{+}\cap L_{-})=L_{+}$），再用 4.2.5(ii) 证明 a)。

\bdp{4.2.12} 设 $E$ 是 $X$ 上的向量丛。则 $E\oplus(E^{*}\otimes\omega_{X})$ 具有显然的 $\omega$-正交丛结构。我们将构造典范同构
\begin{equation}\tag{180}
\operatorname{Pf}(E\oplus(E^{*}\otimes\omega_{X}))\xrightarrow{\ \sim\ }\det R\Gamma(X,E).
\end{equation}
选取子层 $E'\subset E$ 与局部自由层 $E''\supset E$，使得 $H^{0}(X,E')=0$、$H^{1}(X,E'')=0$，且 $E''/E'$ 具有限支撑。置
\[
V:=H^{0}(X,E''/E'),\qquad L_{+}:=H^{0}(X,E'/E)\subset V,\qquad L_{-}:=H^{0}(X,E'')\subset V.
\]
则 $R\Gamma(X,E)$ 可等同于复形 $0\to L_{-}\to V/L_{+}\to 0$，而 $\det R\Gamma(X,E)$ 可等同于 $\det(V;L_{+},L_{-})$。另一方面，$\mathcal{Q}:=E\oplus(E^{*}\otimes\omega_{X})$ 的 Pfaffian 可用子层 $\mathcal{Q}':=E'\oplus((E'')^{*}\otimes\omega_{X})\subset\mathcal{Q}$ 来计算。于是 $\operatorname{Pf}_{\mathcal{Q}'}(\mathcal{Q})$ 等于 (170) 的左端。故 (170) 给出同构 (180)。可以验证 (180) 不依赖 $E'$ 与 $E''$。

\bdp{4.2.13} Lagrange 三元组的概念有一个有用的无穷维推广。先回忆一些基本定义。

\emph{定义.} \emph{Tate 空间}是指完备拓扑向量空间，其 $0$ 的邻域基由一族可公度的向量子空间构成（即对基中任意 $U_{1},U_{2}$，有 $\dim U_{1}/(U_{1}\cap U_{2})<\infty$）。

\emph{注.} Tate 空间隐含在 Tate 的卓越著作 [T] 之中。事实上，[T] 中发展的曲线上的留数理论，最自然的解释正是通过 Tate 空间的自同态代数\SourcePage{143}的典范中心扩张；这一点也隐含在 [T] 中。Tate 扩张的构造可参看 7.13.18。

设 $V$ 是 Tate 空间。若对每个开子空间 $U\subset V$ 都存在有限集 $\{v_{1},\dots,v_{n}\}\subset V$，使得 $P\subset U+\mathbb{C}v_{1}+\dots+\mathbb{C}v_{n}$，则称向量子空间 $P\subset V$ \emph{有界}。$V$ 的\emph{拓扑对偶}是 $V$ 上连续线性泛函构成的空间 $V^{*}$，其上的（线性）拓扑使得 $V$ 的有界子空间的正交补构成 $0\in V^{*}$ 的邻域基。显然 $V^{*}$ 是 Tate 空间，且典范态射 $V\to(V^{*})^{*}$ 是同构。

\emph{例}（坐标 Tate 空间）。设 $I$ 是集合。若 $A\setminus(A\cap B)$ 与 $B\setminus(B\cap A)$ 都有限，则称 $A,B\subset I$ \emph{可公度}。可公度性是一个等价关系。设 $\mathcal{A}$ 是子集 $A\subset I$ 的族的一个等价类。$\mathcal{A}$ 的元素称为\emph{半无限子集}。记 $\mathbb{C}((I,\mathcal{A}))$ 为形式线性组合 $\sum_{i}c_{i}e_{i}$（$c_{i}\in\mathbb{C}$）构成的空间，其中要求 $c_{i}$ 在 $i\notin A$ 时为零，而 $A$ 是某个半无限子集。它是 Tate 向量空间（拓扑由子空间 $\mathbb{C}[[A]]:=\{\sum_{i\in A}c_{i}e_{i}\}$ 定义，其中 $A$ 半无限）。$\mathbb{C}((I,\mathcal{A}))$ 的对偶空间是 $\mathbb{C}((I,\mathcal{A}'))$，其中 $\mathcal{A}'$ 由 $\mathcal{A}$ 中子集的补集构成。任何 Tate 向量空间都同构于适当选取的 $I,\mathcal{A}$ 的 $\mathbb{C}((I,\mathcal{A}))$；这样的同构由相应的子集 $\{e_{i}\}\subset V$ 给出，它称为 $V$ 的\emph{拓扑基}。

$V$ 中的 \emph{c-格}是开有界子空间。$V$ 中的 \emph{d-格}\footnote{c 与 d 分别是 「compact」 与 「discrete」 的首字母。}是离散子空间 $\Gamma\subset V$，使得对某个 c-格 $P\subset V$ 有 $\Gamma+P=V$。若 $W\subset V$ 是 d-格（相应地，c-格），则存在 c-格（相应地，d-格）$W'\subset V$，使得 $V=W\oplus W'$。若 $W\subset V$ 是 d-格（相应地，c-格），则 $W^{\perp}\subset V^{*}$ 也是 d-格（相应地，c-格），且 $(W^{\perp})^{\perp}=W$。

Tate 空间 $V$ 上的（连续）双线性型称为\emph{非退化}的，是指它诱导出拓扑同构 $V\to V^{*}$。设 $V$ 是配备非退化对称双线性型的 Tate 空间。若 $L^{\perp}=L$，则称子空间 $L\subset V$ 为 \emph{Lagrange} 的。

\SourcePage{144}
\emph{定义.} \emph{Tate Lagrange 三元组}由以下数据构成：配备非退化对称双线性型的 Tate 空间 $V$，$V$ 中的 Lagrange c-格 $L_{+}\subset V$，以及 $V$ 中的 Lagrange d-格 $L_{-}\subset V$。

\emph{例.} 设 $\mathcal{Q}$ 是 $X$ 上的 $\omega$-正交丛。对 $x\in X$，记 $\mathcal{Q}\otimes\mathcal{O}_{x}$（相应地，$\mathcal{Q}\otimes K_{x}$）为 $\mathcal{Q}$ 拉回到 $\operatorname{Spec}\mathcal{O}_{x}$（相应地，$\operatorname{Spec}K_{x}$）上之后的整体截面空间。$\mathcal{Q}\otimes K_{x}$ 是 Tate 空间，带非退化对称双线性型 $\operatorname{Res}(\ ,\ )$。对每个非空有限子集 $S\subset X$，有 Tate Lagrange 三元组
\begin{equation}\tag{181}
V:=\oplus_{x\in S}(\mathcal{Q}\otimes K_{x}),\qquad L_{+}:=\oplus_{x\in S}(\mathcal{Q}\otimes\mathcal{O}_{x}),\qquad L_{-}:=\Gamma(X\setminus S,\mathcal{Q}).
\end{equation}

设 $(V;L_{+},L_{-})$ 是 Tate Lagrange 三元组。则对任意 c-格 $\Lambda\subset L_{+}$，只要 $\Lambda\cap L_{-}=0$，就有 (178) 所定义的有限维 Lagrange 三元组 $(\tilde{V};\tilde{L}_{+},\tilde{L}_{-})$。如 4.2.11 所述，$\operatorname{Pf}(\tilde{V};\tilde{L}_{+},\tilde{L}_{-})$ 与 $\det(\tilde{V};\tilde{L}_{+},\tilde{L}_{-})$ 不依赖 $\Lambda$。置
\[
\operatorname{Pf}(V;L_{+},L_{-}):=\operatorname{Pf}(\tilde{V};\tilde{L}_{+},\tilde{L}_{-}),\qquad
\det(V;L_{+},L_{-}):=\det(\tilde{V};\tilde{L}_{+},\tilde{L}_{-}).
\]
等价地，可以把 $\det(V;L_{+},L_{-})$ 定义为复形 (164) 的行列式，而把 $\operatorname{Pf}(V;L_{+},L_{-})$ 用 (166) 或 (167) 来定义（此时应设 (166) 中的 $\operatorname{Cl}(V)$-模 $M$ 是\emph{离散}的，即对每个 $m\in M$，$\{v\in V\mid vm=0\}$ 是开集）。

\emph{例.} 若 $(V;L_{+},L_{-})$ 由 (181) 定义，则 $\operatorname{Pf}(V;L_{+},L_{-})=\operatorname{Pf}(\mathcal{Q})$，$\det(V;L_{+},L_{-})=\det R\Gamma(X,\mathcal{Q})$。

4.2.7 与 4.2.8 中的构造在 Tate 情形下仍有意义，只需作如下显然的改动：a) 在 4.2.7 中应设 $L_{+}$ 是 c-格，$L_{-}$ 是 d-格；b) (171) 应换成如下公式：
\begin{equation}\tag{182}
M=\varinjlim_{U}\bigwedge(V/U)\otimes\det(L_{+}/U)^{*}
\end{equation}
其中 $U$ 取遍 $L_{+}$ 中的 c-格。(182) 的右端是\emph{费米 Fock 空间}，即 $V$ 的\emph{半无限外幂}的直和（参看 [KR] 第 4 讲及其中的文献）。

\SourcePage{145}
\emph{注.} 用三元组 (181) 表述 $\operatorname{Pf}(\mathcal{Q})$ 的式子可以改写如下。对 $x\in X$，考虑交换李超代数 $\mathfrak{a}_{\mathcal{O}_{x}}\subset\mathfrak{a}_{K_{x}}$，使 $\mathfrak{a}_{\mathcal{O}_{x}}$（相应地，$\mathfrak{a}_{K_{x}}$）的奇部为 $\mathcal{Q}\otimes\mathcal{O}_{x}$（相应地，$\mathcal{Q}\otimes K_{x}$），偶部为 $0$。$\mathcal{Q}\otimes K_{x}$ 上的对称双线性型定义了中心扩张 $0\to\mathbb{C}\to\tilde{\mathfrak{a}}_{K_{x}}\to\mathfrak{a}_{K_{x}}\to 0$，它在 $\mathfrak{a}_{\mathcal{O}_{x}}$ 上有典范分裂。Clifford 代数 $\operatorname{Cl}(\mathcal{Q}\otimes K_{x})$ 就是扭泛包络代数 $\operatorname{U}'\mathfrak{a}_{K_{x}}$，而
\[
M_{x}:=\operatorname{Cl}(\mathcal{Q}\otimes K_{x})/\operatorname{Cl}(\mathcal{Q}\otimes K_{x})\cdot(\mathcal{Q}\otimes\mathcal{O}_{x})
\]
是 $\operatorname{U}'\mathfrak{a}_{K_{x}}$ 上的真空模。按 (167)，$\operatorname{Pf}(\mathcal{Q})^{*}$ 是 $\Gamma(X\setminus S,\mathcal{Q})$ 在 $\otimes_{x\in S}M_{x}$ 上作用的余不变量空间。

\bdp{4.2.14} 本小节讨论 Tate Lagrange 三元组的\emph{族}。设 $R$ 是交换环。\emph{Tate $R$-模}是指同构于 $P\oplus Q^{*}$ 的拓扑 $R$-模，其中 $P,Q$ 是有限生成投射 $R$-模的（无限）直和（$0\in P\oplus Q^{*}$ 的邻域基由所有有限生成子模 $M\subset Q$ 对应的 $M^{\perp}\subset Q^{*}$ 构成）。这个\emph{坏}\footnote{有限秩投射 $R((t))$-模不必是上述意义的 Tate 模。我们的 Tate $R$-模概念对 $\operatorname{Spec}R$ 不是局部的。它还有其它缺点。}定义对我们的目的够用。事实上，我们主要处理同构于 $V_{0}\hat{\otimes}R$ 的 Tate $R$-模，其中 $V_{0}$ 是 Tate 空间。

若如下定义 c-格与 d-格，则 4.2.13 中 Tate 线性代数的讨论对 Tate $R$-模仍成立。

\emph{定义.} Tate $R$-模 $V$ 中的 \emph{c-格}是开有界子模 $P\subset V$，使 $V/P$ 投射。$V$ 中的 \emph{d-格}是子模 $\Gamma\subset V$，使得对 $V$ 中的某个 c-格 $P$ 有 $\Gamma\cap P=0$，且 $V/(\Gamma+P)$ 是有限型投射模。\footnote{于是这对所有 c-格 $P'\subset P$ 也成立。}

现若 $1/2\in R$，则可完全像 4.2.13 那样定义 Tate Lagrange 三元组（当然，若 $1/2\notin R$，则应改用二次型代替双线性型，这很容易）。Tate Lagrange 三元组的 Pfaffian\SourcePage{146}$(V;L_{+},L_{-})$ 在 $R$ 上按 4.2.13 中同样的方式定义，只作如下小的修改：为过渡到 (178) 所定义的有限维 Lagrange 三元组 $(\tilde{V};\tilde{L}_{+},\tilde{L}_{-})$，须设 $\Lambda\subset L_{+}$ 是 c-格，使 $\Lambda\cap L_{-}=0$ 且 $V/(\Lambda+L_{-})$ 投射（这两个性质等价于 $\Lambda^{\perp}+L_{-}=V$）。

\emph{例.} 设 $D\subset X\otimes R$ 是 $\operatorname{Spec}R$ 上的有限闭子概形，可局部由一个方程定义（即 $D$ 是有效相对 Cartier 除子）。设 $\mathcal{Q}$ 是 $X\otimes R$ 上的向量丛。设态射 $D\to\operatorname{Spec}R$ 满射。则
\[
V:=\varprojlim_{m}\varinjlim_{n}H^{0}(X\otimes R,\mathcal{Q}(nD)/\mathcal{Q}(-mD))
\]
是 Tate $R$-模\footnote{事实上 $V$ 同构于某个 Tate 空间 $V_{0}$ 的 $V_{0}\hat{\otimes}R$。确实，可以设 $R$ 在 $\mathbb{C}$ 上有限生成，再用 7.12.11。我们需要 7.12.11 的下述情形：$R$ 在 $\mathbb{C}$ 上有限生成，且 7.12.11 中的投射模是有限生成模的直和；此时 7.12.11 由 Serre 定理（[Se] 定理 1；另见 [Ba68] 第 4 章 §2）与 Eilenberg 引理 [Ba63] 推出。}，
\[
L_{+}:=\varprojlim_{m}H^{0}(X\otimes R,\mathcal{Q}/\mathcal{Q}(-mD))\subset V
\]
是 c-格，且
\[
L_{-}:=H^{0}((X\otimes R)\setminus D,\mathcal{Q})\subset V
\]
是 d-格。若 $\mathcal{Q}$ 是 $\omega$-正交丛，则 $(V;L_{+},L_{-})$ 是 Lagrange 三元组，且 $\operatorname{Pf}(\mathcal{Q})=\operatorname{Pf}(V;L_{+},L_{-})$（参看 4.2.13）。

\bdp{4.2.15} 记 $\mathcal{B}$ 为 $\mathbb{C}$ 上配备非退化对称双线性型的有限维向量空间构成的群胚。本小节（读者可以跳过）构造典范同构
\begin{equation}\tag{183}
\begin{aligned}
&\operatorname{Pf}(V\otimes W;L_{+}\otimes W,L_{-}\otimes W)\\
&\qquad\xrightarrow{\ \sim\ }\operatorname{Pf}(V;L_{+},L_{-})^{\otimes\dim W}\otimes|\det W|^{\otimes p(V;L_{+},L_{-})},
\end{aligned}
\end{equation}
% [Tutor Note: 原书 p.146 把编号 (183) 单独印在公式上一行的左端，公式本身未带编号；底本如此，此处按 equation 排版。]
\begin{equation}\tag{184}
\operatorname{Pf}(\mathcal{Q}\otimes W)\xrightarrow{\ \sim\ }\operatorname{Pf}(\mathcal{Q})^{\otimes\dim W}\otimes|\det W|^{\otimes p(\mathcal{Q})}
\end{equation}

\SourcePage{147}
其中 $W\in\mathcal{B}$，$(V;L_{+},L_{-})$ 是（Tate）Lagrange 三元组，$\mathcal{Q}$ 是 $X$ 上的 $\omega$-正交丛，$|\det W|$ 是把 $W$ 看作空间（不是超空间！）时的行列式，而 $p(V;L_{+},L_{-}),p(\mathcal{Q})\in\mathbb{Z}/2\mathbb{Z}$ 分别是 $\operatorname{Pf}(V;L_{+},L_{-})$ 与 $\operatorname{Pf}(\mathcal{Q})$ 的奇偶性。$|\det W|^{\otimes n}$ 对 $n\in\mathbb{Z}/2\mathbb{Z}$ 有意义，因为有典范同构 $|\det W|^{\otimes 2}\xrightarrow{\ \sim\ }\mathbb{C}$，$(w_{1}\wedge\dots\wedge w_{r})^{\otimes 2}\mapsto\det(w_{i},w_{j})$。

为定义 (183) 与 (184)，注意 $\mathcal{B}$ 是以 $\oplus$ 为张量「积」的张量范畴。(183) 与 (184) 的两端都是从 $\mathcal{B}$ 到一维超空间范畴的张量函子（为把 (184) 的右端看作张量函子，把它改写为 $|\operatorname{Pf}(\mathcal{Q})|^{\otimes\dim W}\otimes(\det W)^{\otimes p(\mathcal{Q})}$，其中 $|\operatorname{Pf}(\mathcal{Q})|$ 由 $\operatorname{Pf}(\mathcal{Q})$ 改变 $(\mathbb{Z}/2\mathbb{Z})$-分次使其为偶而得到，而 $\det W$ 是把 $W$ 看作\emph{超空间}时的行列式）。

断言：把 (183) 与 (184) 定义为张量函子的同构，存在唯一的方式使得对 $W=(\mathbb{C},1)$，(183) 与 (184) 等于 $\operatorname{id}$。这里 $1$ 表示双线性型 $(x,y)\mapsto xy$，$x,y\in\mathbb{C}$。

为证明，对由 (183) 或 (184) 的左端除以其右端得到的张量函子 $F$，应用下述引理。

\emph{引理.} 每个张量函子 $F:\mathcal{B}\to\{1\text{ 维向量空间}\}$ 都同构于张量函子 $F_{1}$，后者由 $F_{1}(W)=L^{\otimes\dim W}$ 定义，其中 $L:=F(\mathbb{C},1)$。存在唯一同构 $F\xrightarrow{\ \sim\ }F_{1}$，它诱导恒等映射 $F(\mathbb{C},1)\to F_{1}(\mathbb{C},1)$。

\emph{证明.} 对每个 $W\in\mathcal{B}$，函子 $F$ 诱导同态 $f_{W}:\operatorname{Aut}W\to\mathbb{C}^{*}$。因 $\operatorname{Aut}W$ 是正交群，故 $f_{W}(g)=(\det g)^{n(W)}$，其中 $n(W)\in\mathbb{Z}/2\mathbb{Z}$。显然 $n(W)=n$ 与 $W$ 无关。置 $W_{1}:=(\mathbb{C},1)$。$F$ 把交换性同构 $\left(\begin{smallmatrix}0&1\\1&0\end{smallmatrix}\right):W_{1}\oplus W_{1}\to W_{1}\oplus W_{1}$ 映为 $\operatorname{id}$。故 $n=0$，即 $f_{W}$ 对每个 $W$ 平凡。其余是显然的，因为 $\mathcal{B}$ 的对象同构类的半群 $|\mathcal{B}|$ 就是 $\mathbb{Z}_{+}$。\hfill$\Box$

\emph{注.}
\begin{enumerate}[label=(\roman*),leftmargin=2.2em]
\item (183) 在 4.2.8 中已被隐含地使用。
\item 我们将在 4.2.16 中使用 (183)。
\end{enumerate}

\SourcePage{148}
\bdp{4.2.16} 本小节（读者当然可以跳过）解释把 $\mathbb{C}$ 换成特征 $2$ 的域 $k$ 时会发生什么。此时必须区分二次型（见 [Bourb59] §3, n°4）与对称双线性型。在 Lagrange 三元组的定义中，$V$ 应带非退化\emph{二次}型。于是 $\omega$-正交丛的定义中，$\mathcal{Q}$ 应带非退化\emph{二次}型 $\mathcal{Q}\to\omega_{X}$（因 $k$ 的特征为 $2$，非退化性蕴含 $\mathcal{Q}$ 的秩为偶）。4.2.8 中 $\operatorname{Pf}\otimes\operatorname{Pf}^{-}\xrightarrow{\ \sim\ }\det$ 的构造须修改。若 $(V;L_{+},L_{-})$ 是 Lagrange 三元组，而 $W$ 带非退化对称\emph{双线性}型，则 $(V\otimes W;L_{+}\otimes W,L_{-}\otimes W)$ 是 Lagrange 三元组。双线性型 $\left(\begin{smallmatrix}1&0\\0&-1\end{smallmatrix}\right)$ 与 $\left(\begin{smallmatrix}0&1\\1&0\end{smallmatrix}\right)$ 在特征 $2$ 下不等价，但可以对 $W=H$ 与 $W=H'$ 使用 (183) 来构造 $\operatorname{Pf}\otimes\operatorname{Pf}^{-}\xrightarrow{\ \sim\ }\det$。最后，须在特征 $2$ 下构造 (183) 与 (184)。为简单起见设 $k$ 完全。则 (183) 与 (184) 的同构的特征性质\footnote{在非完全情形表述此性质时，应把 $\mathcal{B}$ 看作栈而非群胚。}的表述与 4.2.15 相同，但存在性与唯一性的证明须修改。半群 $|\mathcal{B}|$（见 4.2.15 引理证明的末尾）不再是 $\mathbb{Z}_{+}$；它有生成元 $a,b$ 与定义关系 $a+b=3a$（$a$ 对应 $1$ 阶矩阵 $(1)$，$b$ 对应 $\left(\begin{smallmatrix}0&1\\1&0\end{smallmatrix}\right)$）。故 $\mathcal{B}$ 对应的群是 $\mathbb{Z}$，这已够用。

\bdsec{4.3. Pfaffian（二）}

\bdp{4.3.1} 固定 $\mathbb{C}$ 上 $n$ 维向量空间 $W$ 与其上的非退化对称双线性型 $(\ )$。为简化记号，把 $O(W)$ 与 $SO(W)$ 写作 $O_{n}$ 与 $SO_{n}$。

设 $\mathcal{F}$ 是 $X$ 上的 $SO_{n}$-torsor。相应的秩 $n$ 向量丛 $W_{\mathcal{F}}$ 带双线性型 $(\ )_{\mathcal{F}}$，且有典范同构 $\det W_{\mathcal{F}}=\mathcal{O}_{X}\otimes\det W$。设 $\mathcal{L}\in\omega^{1/2}(X)$，即 $\mathcal{L}$ 是 $\omega_{X}$ 的平方根。

\SourcePage{149}
则 $W_{\mathcal{F}}\otimes\mathcal{L}$ 是 $\omega$-正交丛，故 $\operatorname{Pf}(W_{\mathcal{F}}\otimes\mathcal{L})$ 有意义（见 4.2）。考虑「归一化」Pfaffian
\begin{equation}\tag{185}
\operatorname{Pf}_{\mathcal{L},\mathcal{F}}:=\operatorname{Pf}(W_{\mathcal{F}}\otimes\mathcal{L})\otimes\operatorname{Pf}(W\otimes\mathcal{L})^{\otimes-1}
\end{equation}
与「归一化」行列式
\begin{equation}\tag{186}
\nu(\mathcal{F}):=\det R\Gamma(X,W_{\mathcal{F}})\otimes\det R\Gamma(X,\mathcal{O}_{X}\otimes W)^{\otimes-1}.
\end{equation}
如 4.2.1 所述，存在典范同构 $c_{\pm i}:\operatorname{Pf}^{\otimes 2}\xrightarrow{\ \sim\ }\det R\Gamma$。使用例如 $c_{i}$ 得到同构\footnote{故同构 (187)=(189) 依赖 $-1$ 的平方根的选取。若把 (187) 乘以 $i^{\pm p(\mathcal{F})^{2}}$，则此依赖消失，其中 $p$ 是典范映射 $\operatorname{Bun}_{SO_{n}}\to\pi_{0}(\operatorname{Bun}_{SO_{n}})=\pi_{1}(SO_{n})=\mathbb{Z}/2\mathbb{Z}$，而 $p(\mathcal{F})^{2}\in\mathbb{Z}/4\mathbb{Z}$。由于 4.2.1 的注 (iii) 所述的理由，我们宁愿不这样做。}
\begin{equation}\tag{187}
\operatorname{Pf}_{\mathcal{L},\mathcal{F}}^{\otimes 2}\xrightarrow{\ \sim\ }\nu_{\mathcal{L}}(\mathcal{F})
\end{equation}
其中
\begin{equation}\tag{188}
\nu_{\mathcal{L}}(\mathcal{F}):=\det R\Gamma(X,W_{\mathcal{F}}\otimes\mathcal{L})\otimes\det R\Gamma(X,W\otimes\mathcal{L})^{\otimes-1}.
\end{equation}
[Del87] 中的构造 7.2 给出典范同构
\[
\nu_{\mathcal{L}}(\mathcal{F})=\nu(\mathcal{F})\otimes\langle\det W_{\mathcal{F}}\otimes(\det W)^{\otimes-1},\mathcal{L}\rangle
\]
因 $\det W_{\mathcal{F}}=\mathcal{O}_{X}\otimes\det W$，故 $\nu_{\mathcal{L}}(\mathcal{F})=\nu(\mathcal{F})$，且
\begin{equation}\tag{189}
\operatorname{Pf}_{\mathcal{L},\mathcal{F}}^{\otimes 2}=\nu(\mathcal{F}).
\end{equation}
$\mathcal{F}$ 变化时，$\operatorname{Pf}_{\mathcal{L},\mathcal{F}}$ 与 $\nu(\mathcal{F})$ 是 $\operatorname{Bun}_{SO_{n}}$ 上两个线丛的纤维，二者分别记作 $\operatorname{Pf}_{\mathcal{L}}$ 与 $\nu$。

记 $\nu^{1/2}(\operatorname{Bun}_{SO_{n}})$ 为 $\nu$ 的平方根构成的范畴。有函子
\begin{equation}\tag{190}
\operatorname{Pf}:\omega^{1/2}(X)\to\nu^{1/2}(\operatorname{Bun}_{SO_{n}})
\end{equation}
它由 $\mathcal{L}\mapsto\operatorname{Pf}_{\mathcal{L}}$ 定义。

\SourcePage{150}
$\omega^{1/2}(X)$ 与 $\nu^{1/2}(\operatorname{Bun}_{SO_{n}})$ 是 Picard 范畴 $\mu_{2}\operatorname{tors}(X)$ 与 $\mu_{2}\operatorname{tors}(\operatorname{Bun}_{SO_{n}})$ 上的 torsor。有 Picard 函子 $\ell^{\operatorname{Spin}}:\mu_{2}\operatorname{tors}(X)\to\mu_{2}\operatorname{tors}(\operatorname{Bun}_{SO_{n}})$；这就是 4.1 中的函子 $\ell=\ell^{\tilde{G}}$ 在 $G=SO_{n}$、$\tilde{G}=\operatorname{Spin}_{n}$、$\Pi=\mathbb{Z}/2\mathbb{Z}$ 时的特例。在 4.3.8–4.3.15 中我们将证明，函子 $\operatorname{Pf}:\omega^{1/2}(X)\to\nu^{1/2}(\operatorname{Bun}_{SO_{n}})$ 具有典范的 $\ell^{\operatorname{Spin}}$-仿射函子结构。在此之先，4.3.2–4.3.7 将证明：对有限子集 $S\subset X$，4.1.7 所定义的 $SO_{n}(K_{S})$ 在 $\operatorname{Bun}_{SO_{n},S}$ 上的作用，提升为 $SO_{n}(K_{S})$ 的某个中心扩张在 $\operatorname{Pf}_{\mathcal{L}}$ 到 $\operatorname{Bun}_{SO_{n},S}$ 的拉回上的作用。一旦引入此作用，就很容易用 $SO_{n}(K_{S})$-不变性来刻画函子 $\operatorname{Pf}$ 上的 $\ell^{\operatorname{Spin}}$-仿射结构（见 4.3.8–4.3.10）。

\bdp{4.3.2} 设 $V$ 是带\emph{偶型}非退化对称双线性型的 Tate 空间，即存在 Lagrange c-格 $L\subset V$（见 4.2.13）；若 $\dim V<\infty$，这表示 $\dim V$ 为偶。记 $O(V)$ 为 $V$ 的保持该型的拓扑自同构群。回忆典范中心扩张的熟知构造
\begin{equation}\tag{191}
0\to\mathbb{C}^{*}\to\tilde{O}(V)\to O(V)\to 0.
\end{equation}
设 $M$ 是 Clifford 代数 $\operatorname{Cl}(V)$ 上的不可约 $(\mathbb{Z}/2\mathbb{Z})$-分次离散模（离散是指：对每个 $m\in M$，$\{v\in V\mid vm=0\}$ 是开集）。则 $M$ 在张量一维 $(\mathbb{Z}/2\mathbb{Z})$-分次空间的意义下唯一。故 $O(V)$ 在 $M$ 中有自然的射影表示。(191) 就是与此表示对应的扩张，即
\[
\tilde{O}(V):=\{(g,\varphi)\mid g\in O(V),\ \varphi\in\operatorname{Aut}_{\mathbb{C}}M,\ \varphi(vm)=g(v)\cdot\varphi(m)\ \text{对}\ m\in M\}.
\]
显然 $\tilde{O}(V)$ 不依赖 $M$ 的选取（事实上 $\operatorname{Aut}_{\mathbb{C}}M$ 是 $\operatorname{Cl}(V)$ 的自然完备化的可逆元群）。若 $(g,\varphi)\in\tilde{O}(V)$，则 $\varphi$ 或为偶或为奇。记 $\chi(g)\in\mathbb{Z}/2\mathbb{Z}$ 为 $\varphi$ 的奇偶性。则 $\chi:O(V)\to\mathbb{Z}/2\mathbb{Z}$ 是同态。

\SourcePage{151}
$-1\in O(V)$ 在 $\tilde{O}(V)$ 中的原像不是中心的。事实上，若 $\varphi:M\to M$ 满足：$m$ 为偶时 $\varphi(m)=m$，$m$ 为奇时 $\varphi(m)=-m$，则 $[-1]:=(-1,\varphi)\in\tilde{O}(V)$，且
\begin{equation}\tag{192}
[-1]\cdot\tilde{g}=(-1)^{\chi(g)}\cdot\tilde{g}\cdot[-1],\qquad g\in O(V)
\end{equation}
其中 $\tilde{g}$ 表示 $g$ 在 $\tilde{O}(V)$ 中的一个原像。

$O(V)$ 与 $\operatorname{Aut}_{\mathbb{C}}M$ 有自然的群归纳概形结构。更确切地说，把 $\mathbb{C}$-代数 $R$ 对应到集合 $O(V\hat{\otimes}R)$ 与 $\operatorname{Aut}_{\mathbb{C}}(M\hat{\otimes}R)$ 的函子是归纳概形（若 $\dim V=\infty$，则它们可表示为一族\emph{不可数}的滤过闭子概形之并）。故 $\tilde{O}(V)$ 是群归纳概形。

记 $\operatorname{Lagr}(V)$ 为 $V$ 中 Lagrange c-格的集合。它有自然的归纳概形结构：$\operatorname{Lagr}(V)=\varinjlim_{\Lambda}\operatorname{Lagr}(\Lambda^{\perp}/\Lambda)$，其中 $\Lambda$ 取遍 $V$ 中的迷向 c-格（故 $\operatorname{Lagr}(V)$ 的 $R$-点就是 4.2.14 意义下 $V\hat{\otimes}R$ 中的 Lagrange c-格）。记 $\mathcal{P}=\mathcal{P}_{M}$ 为 $\operatorname{Lagr}(V)$ 上的线丛，它在 $L\in\operatorname{Lagr}(V)$ 上的纤维等于 $M^{L}:=\{m\in M\mid Lm=0\}$。$O(V)$ 在 $\operatorname{Lagr}(V)$ 上的作用典范地提升为 $\tilde{O}(V)$ 在 $\mathcal{P}$ 上的作用。

$\operatorname{Lagr}(V)$ 有两个连通分支，由一维 $(\mathbb{Z}/2\mathbb{Z})$-分次空间 $M^{L}$（$L\in\operatorname{Lagr}(V)$）的奇偶性区分。此断言的证明容易化归到 $\dim V$ 有限（且为偶）的情形。同样的论证表明：$L_{1},L_{2}\in\operatorname{Lagr}(V)$ 属于同一分支，当且仅当 $\dim(L_{1}/(L_{1}\cap L_{2}))$ 为偶。显然，$\operatorname{Lagr}(V)$ 的连通分支对 $g\in O(V)$ 不变，当且仅当 $\chi(g)=0$。因此 $\chi:O(V)\to\mathbb{Z}/2\mathbb{Z}$ 是群归纳概形的态射。

我们来证明 (191) 来自群归纳概形的正合列
\begin{equation}\tag{193}
0\to\mathbb{G}_{m}\to\tilde{O}(V)\to O(V)\to 0.
\end{equation}
只需说明态射 $\tilde{O}(V)\to O(V)$ 是 $\mathbb{G}_{m}$-torsor。为此固定 $L\in\operatorname{Lagr}(V)$，并置 $M=\operatorname{Cl}(V)/\operatorname{Cl}(V)L$，于是 $\mathcal{P}=\mathcal{P}_{M}$ 在 $L$ 上的纤维等于 $\mathbb{C}$。定义 $f:O(V)\to\operatorname{Lagr}(V)$ 为 $f(g)=gL$。
% ===================== end 源码/units/ch04_u03.tex =====================

% ===================== begin 源码/units/ch04_u04.tex =====================
% ch04_u04 —— 原书 p.152–p.161（4.3.3–4.3.16）
% 承接 p.151 末 4.3.2 的证明（$f(g)=gL$）。
\SourcePage{152}

置 $\mathcal{P}':=\mathcal{P}\setminus\{\text{零截面}\}$；这是 $\operatorname{Lagr}(V)$ 上的 $\mathbb{G}_{m}$-torsor。
易证自然态射 $\tilde{O}(V)\to f^{*}\mathcal{P}'$ 是同构，故 $\tilde{O}(V)$ 是 $O(V)$ 上的 $\mathbb{G}_{m}$-torsor。

\emph{注.} 设 $L\in\operatorname{Lagr}(V)$。则 (193) 在 $O(V)$ 中 $L$ 的稳定化子上典范分裂：若 $g\in O(V)$、$gL=L$，则 $g$ 在 $\tilde{O}(V)$ 中有唯一原像，它在 $M^{L}$ 上作用为恒等。

\bdp{4.3.3} 置 $O:=\mathbb{C}[[t]]$，$K:=\mathbb{C}((t))$。以 $\omega_{O}$ 表 $O$ 的（完备化）微分模。固定 $\omega_{O}$ 的一个平方根，即一维自由 $O$-模 $\omega_{O}^{1/2}$，连同同构 $\omega_{O}^{1/2}\otimes\omega_{O}^{1/2}\xrightarrow{\ \sim\ }\omega_{O}$。$W$ 的含义同 4.3.1。我们要构造 $O_{n}(K):=O(W\otimes K)$ 的中心扩张，其中 $O_{n}(K)$ 看作 $\mathbb{C}$ 上的群归纳概形。

置 $\omega_{K}^{1/2}:=\omega_{O}^{1/2}\otimes_{O}K$，$\omega_{K}:=\omega_{O}\otimes_{O}K$。考虑 Tate 空间 $V:=\omega_{K}^{1/2}\otimes W$。$W$ 上的双线性型诱导 $K$-双线性型 $V\times V\to\omega_{K}$；与 $\operatorname{Res}:\omega_{K}\to\mathbb{C}$ 复合，得非退化对称双线性型 $V\times V\to\mathbb{C}$，其型为偶。把扩张 (193) 限制到 $O_{n}(K)\hookrightarrow O(V)$，得中心扩张
\begin{equation}\tag{194}
0\to\mathbb{G}_{m}\to\tilde{O}_{n}(K)\to O_{n}(K)\to 0.
\end{equation}
它在 $O_{n}(O)\subset O_{n}(K)$ 上典范分裂（用 4.3.2 末的注，取 $L=\omega_{O}^{1/2}\otimes W\subset V$）。群 $\operatorname{Aut}\omega_{O}^{1/2}=\mu_{2}$ 作用在扩张 (194) 上，并保持 $O_{n}(O)$ 上的分裂。

\bdp{4.3.4} \emph{引理.} 由 $-1\in\operatorname{Aut}\omega_{O}^{1/2}$ 诱导的 $\tilde{O}_{n}(K)$ 自同构，把 $\tilde{g}\in\tilde{O}_{n}(K)$ 映到 $(-1)^{\theta(g)}\tilde{g}$，其中 $g$ 是 $\tilde{g}$ 在 $O_{n}(K)$ 中的像，$\theta:O_{n}(K)\to K^{*}/(K^{*})^{2}=\mathbb{Z}/2\mathbb{Z}$ 是旋量范数。

\emph{证明.} 由 (192)，只需证 $g\in O_{n}(K)\subset O(V)$ 时 $\chi(g)=\theta(g)$。按 $\theta$ 的定义（见 [D71] 第 II 章 §7），只需证：若 $g$ 是关于某非迷向 $x\in K^{n}$ 的正交补的反射，则 $\chi(g)$ 等于 $(x,x)\in K^{*}$ 在 $K^{*}/(K^{*})^{2}=\mathbb{Z}/2\mathbb{Z}$ 中的像。可设 $x\in O^{n}$，$x\notin tO^{n}$。

\SourcePage{153}
$L:=\omega_{O}^{1/2}\otimes W$ 是 $V$ 中的 Lagrange c-格，故 $\chi(g)$ 是 $\dim L/(L\cap gL)=\dim O/(x,x)O$ 的奇偶性。

\emph{注记.}
\begin{enumerate}[label=(\roman*),leftmargin=2.2em]
\item 不用反射，也可计算 $\chi$ 在 $SO_{n}(K)$ 的某分裂嘉当子群上的限制，并注意 $g\in O_{n}(\mathbb{C})$ 时 $\chi(g)=0$。
\item $\theta$ 在 $SO_{n}(K)$ 上的限制是边缘态射
\begin{equation}\tag{195}
SO_{n}(K)\to H^{1}(K,\mu_{2})=\mathbb{Z}/2\mathbb{Z},
\end{equation}
它来自正合列 $0\to\mu_{2}\to\operatorname{Spin}_{n}\to SO_{n}\to 0$。
\item 若 $g\in O_{n}(K)=O(W\otimes K)$，则 $\dim(W\otimes O)/((W\otimes O)\cap g(W\otimes O))$ 为偶，当且仅当 $\theta(g)=0$。这由引理 4.3.4 的证明推出。
\end{enumerate}

\bdp{4.3.5} 考虑 (194) 在 $SO_{n}(K)$ 上的限制：
\begin{equation}\tag{196}
0\to\mathbb{G}_{m}\to\tilde{SO}_{n}(K)\to SO_{n}(K)\to 0.
\end{equation}
它在 $SO_{n}(O)$ 上典范分裂。扩张 (196) 依赖 $\omega_{O}^{1/2}$ 的选取，故更应写成 $\tilde{SO}_{n}(K)_{\mathcal{C}}$，其中 $\mathcal{C}$ 是 $\omega_{O}$ 的平方根。设 $\mathcal{C}'$ 是 $\omega_{O}$ 的另一平方根，则 $\mathcal{C}'=\mathcal{C}\otimes\mathcal{A}$，其中 $\mathcal{A}$ 是 $\operatorname{Spec} O$ 上的 $\mu_{2}$-torsor（或 $\operatorname{Spec}\mathbb{C}$ 上的，二者一样）。考虑 $\mathbb{Z}/2\mathbb{Z}$ 被 $\mathbb{G}_{m}$ 的（平凡）扩张，使 $\mathcal{A}$ 成为其分裂构成的 $\mu_{2}$-torsor。它经 (195) 的拉回是（平凡）扩张
\begin{equation}\tag{197}
0\to\mathbb{G}_{m}\to\tilde{SO}_{n}(K)_{\mathcal{A}}\to SO_{n}(K)\to 0,
\end{equation}
并带 $SO_{n}(O)$ 上的分裂（4.1.8 已在更一般的情形引入此扩张）。

引理 4.3.4 给出 $\tilde{SO}_{n}(K)_{\mathcal{C}'}$ 与扩张 $\tilde{SO}_{n}(K)_{\mathcal{C}}$、$\tilde{SO}_{n}(K)_{\mathcal{A}}$ 之和之间的典范同构。它与 $SO_{n}(O)$ 上的分裂相容。

\SourcePage{154}
\bdp{4.3.6} 设 $S$、$O_{S}$、$K_{S}$ 的含义同 4.1.7。固定 $\mathcal{L}\in\omega^{1/2}(X)$，记 $\omega_{K_{S}}^{1/2}$ 为 $\mathcal{L}$ 到 $\operatorname{Spec} K_{S}$ 的拉回的截面空间。然后照 4.3.3 的办法：置 $V:=\omega_{K_{S}}^{1/2}\otimes W$，用「留数之和」映射 $\omega_{K_{S}}\to\mathbb{C}$ 定义 $V$ 上的标量积，把 $SO_{n}(K_{S})$ 嵌入 $O(V)$，最终得中心扩张
\begin{equation}\tag{198}
0\to\mathbb{G}_{m}\to\tilde{SO}_{n}(K_{S})_{\mathcal{L}}\to SO_{n}(K_{S})\to 0,
\end{equation}
并带 $SO_{n}(O_{S})$ 上的典范分裂。

\emph{注.} (198) 是各扩张 (196)（$K=K_{x}$，$x\in S$）的「超和」。说明如下：设 $G_{i}$（$i\in I$）是带态射 $\theta_{i}:G_{i}\to\mathbb{Z}/2\mathbb{Z}$ 的群，$\tilde{G}_{i}$ 是 $G_{i}$ 被 $\mathbb{G}_{m}$ 的中心扩张，则这些扩张的超和是 $\bigoplus_{i\in I}G_{i}$ 被 $\mathbb{G}_{m}$ 的扩张：在通常的和上，添加上标准扩张
\[
0\to\mathbb{G}_{m}\to A\to\bigoplus_{i\in I}(\mathbb{Z}/2\mathbb{Z})\to 0
\]
的拉回。其中 $A$ 由 $\mathbb{G}_{m}$ 与元素 $e_{i}$（$i\in I$）生成，定义关系为 $e_{i}^{2}=1$、$ce_{i}=e_{i}c$（$c\in\mathbb{G}_{m}$）、$e_{i}e_{j}=(-1)\cdot e_{j}e_{i}$（$i\neq j$）。在我们的情形，$\theta_{x}:SO_{n}(K_{x})\to\mathbb{Z}/2\mathbb{Z}$ 是旋量范数。

若 $\mathcal{L},\mathcal{L}'\in\omega^{1/2}(X)$，则 $\mathcal{L}'=\mathcal{L}\otimes\mathcal{E}$，其中 $\mathcal{E}$ 是 $\mu_{2}$-torsor。由 4.3.5，$\tilde{SO}_{n}(K_{S})_{\mathcal{L}'}$ 与扩张 $\tilde{SO}_{n}(K_{S})_{\mathcal{L}}$、$\tilde{SO}_{n}(K_{S})_{\mathcal{E}}$ 之和之间有典范同构（$\tilde{SO}_{n}(K_{S})_{\mathcal{E}}$ 的定义见 4.1.8）。

\bdp{4.3.7} 4.3.1 在 $\operatorname{Bun}_{SO_{n}}$ 上定义了线丛 $\operatorname{Pf}_{\mathcal{L}}$（$\mathcal{L}\in\omega^{1/2}(X)$）。记 $\operatorname{Pf}_{\mathcal{L}}^{S}$ 为 $\operatorname{Pf}_{\mathcal{L}}$ 到 4.1.7 所定义概形 $\operatorname{Bun}_{SO_{n},S}$ 的拉回。有 $SO_{n}(O_{S})\times\mathbb{G}_{m}$ 在 $\operatorname{Pf}_{\mathcal{L}}^{S}$ 上的显然作用（$\lambda\in\mathbb{G}_{m}$ 乘以 $\lambda$）。我们要把它延拓为 $\tilde{SO}_{n}(K_{S})_{\mathcal{L}}$ 在 $\operatorname{Pf}_{\mathcal{L}}^{S}$ 上的作用，与 $SO_{n}(K_{S})$ 在 $\operatorname{Bun}_{SO_{n},S}$ 上的作用相容。

设 $u\in\operatorname{Bun}_{SO_{n},S}$，$\tilde{g}\in\tilde{SO}_{n}(K_{S})_{\mathcal{L}}$。记 $\mathcal{F}$、$\mathcal{F}'$ 为对应 $u$ 与 $gu$ 的 $SO(W)$-丛，其中 $g\in SO_{n}(K_{S})$ 是 $\tilde{g}$ 的像。须定义同构 $\operatorname{Pf}_{\mathcal{L},\mathcal{F}}\cong\operatorname{Pf}_{\mathcal{L},\mathcal{F}'}$，即同构

\SourcePage{155}
\[
\operatorname{Pf}(W_{\mathcal{F}}\otimes\mathcal{L})\xrightarrow{\ \sim\ }\operatorname{Pf}(W_{\mathcal{F}'}\otimes\mathcal{L}).
\]
由 4.2.13，只需构造同构 $\operatorname{Pf}(V;L_{+},L_{-})\xrightarrow{\ \sim\ }\operatorname{Pf}(V;L_{+},L'_{-})$，其中 $V$ 是 4.3.6 的 Tate 空间，$L_{+}=\omega_{O_{S}}^{1/2}\otimes W\subset V$，而 $L_{-},L'_{-}\subset V$ 是离散 Lagrange 子空间，满足 $L'_{-}=gL_{-}$。由 (166)，这等价于构造同构 $f:(M_{L_{-}})^{\otimes-1}\xrightarrow{\ \sim\ }(M_{gL_{-}})^{\otimes-1}$。定义 $f$ 为 $\tilde{g}\in\tilde{O}(V)$ 在 $M$ 上作用所诱导。

注意：$\lambda\in\mathbb{G}_{m}\subset\tilde{SO}_{n}(K_{S})_{\mathcal{L}}$ 在 $\operatorname{Pf}_{\mathcal{L}}^{S}$ 上乘以 $\lambda^{-1}$。

\bdp{4.3.8} 如 4.3.1 所述，目标是给函子 (190) 定义一个典范的 $\ell^{\operatorname{Spin}}$-仿射结构。也就是说，对 $\mathcal{L}\in\omega^{1/2}(X)$ 与 $X$ 上的 $\mu_{2}$-torsor $\mathcal{E}$，须定义同构
\begin{equation}\tag{199}
\operatorname{Pf}_{\mathcal{L}}\otimes\ell^{\operatorname{Spin}}_{\mathcal{E}}\xrightarrow{\ \sim\ }\operatorname{Pf}_{\mathcal{L}'},\qquad \mathcal{L}':=\mathcal{L}\otimes\mathcal{E}.
\end{equation}
还须验证同构 (199) 的某些相容性质。

为简化记号，把 $\ell^{\operatorname{Spin}}_{\mathcal{E}}$ 写成 $\ell_{\mathcal{E}}$。设 $S\subset X$ 有限。4.1.7–4.1.8 构造了中心扩张 $\tilde{SO}_{n}(K_{S})_{\mathcal{E}}$ 在 $\ell^{S}_{\mathcal{E}}:=\ell_{\mathcal{E}}$ 到 $\operatorname{Bun}_{SO_{n},S}$ 的拉回上的作用。故由 4.3.6–4.3.7，$\tilde{SO}_{n}(K_{S})_{\mathcal{L}'}$ 同时作用在 $\operatorname{Pf}_{\mathcal{L}}^{S}\otimes\ell^{S}_{\mathcal{E}}$ 与 $\operatorname{Pf}_{\mathcal{L}'}^{S}$ 上。回忆 (199) 两端在平凡 $SO_{n}$-丛上的纤维都等于 $\mathbb{C}$。

\bdp{4.3.9} \emph{定理.} 存在唯一的同构 (199)，使对每个 $S$，相应同构 $\operatorname{Pf}_{\mathcal{L}}^{S}\otimes\ell^{S}_{\mathcal{E}}\xrightarrow{\ \sim\ }\operatorname{Pf}_{\mathcal{L}'}^{S}$ 是 $\tilde{SO}_{n}(K_{S})_{\mathcal{L}'}$-等变的，且由 (199) 诱导的平凡 $SO_{n}$-丛上纤维之间的同构是恒等。

证明见 4.3.11–4.3.13。较弱陈述的简短证明见 [BLaSo] 的 §5.2。

\bdp{4.3.10} \emph{命题.} 同构 (199) 给出函子 $\operatorname{Pf}:\omega^{1/2}(X)\to\nu^{1/2}(\operatorname{Bun}_{SO_{n}})$ 上的 $\ell^{\operatorname{Spin}}$-仿射结构。

证明见 4.3.15。

\SourcePage{156}
\bdp{4.3.11} 开始证定理 4.3.9。$n>2$ 时 (199) 的唯一性显然：此时 $SO_{n}$ 半单，故有 $G=SO_{n}$、$S\neq\emptyset$ 的同构 (155)。若 $n=2$，$SO_{n}(K_{S})$ 在 $\operatorname{Bun}_{SO_{n},S}$ 上的作用不可迁；但 adele 上的 $SO_{n}$ 在 $\varprojlim_{S}\operatorname{Bun}_{SO_{n},S}(\mathbb{C})$ 上可迁，这足以给出唯一性。

证存在性时设 $n>2$。$n=2$ 的情形可用嵌入 $SO_{2}\hookrightarrow SO_{3}$ 与相应态射 $\operatorname{Bun}_{SO_{2}}\to\operatorname{Bun}_{SO_{3}}$ 处理，或用 4.3.14 末的注处理。

考虑 $\operatorname{Bun}_{SO_{n},S}$ 上 $SO_{n}(K_{S})$-等变的线丛 $\mathcal{C}_{S}:=\operatorname{Pf}_{\mathcal{L}}^{S}\otimes\ell^{S}_{\mathcal{E}}\otimes(\operatorname{Pf}_{\mathcal{L}'}^{S})^{*}$。$\operatorname{Bun}_{SO_{n},S}$ 中对应平凡 $SO_{n}$-丛（连同 $S$ 上的显然平凡化）的点的稳定化子等于 $SO_{n}(A_{S})$，其中 $A_{S}:=H^{0}(X\setminus S,\mathcal{O}_{X})$。于是 $SO_{n}(K_{S})$ 在 $\mathcal{C}_{S}$ 上的作用诱导态射 $f_{S}:SO_{n}(A_{S})\to\mathbb{G}_{m}$。只需证所有 $S$ 的 $f_{S}$ 都平凡（则 $S\neq\emptyset$ 时可用 (155) 得 $\mathcal{C}_{S}$ 的 $SO_{n}(K_{S})$-等变平凡化；这些平凡化当然彼此相容）。

记 $\Sigma$ 为 $X$ 的有限子概形构成的概形（即 $X$ 各对称幂的不交并）。对不必约化的\footnote{这一点在 $S$ 变动时很重要。对固定的 $S$，环 $A_{S}$、$O_{S}$、$K_{S}$ 只依赖 $S_{\mathrm{red}}$。} $S\in\Sigma$，$A_{S}$、$O_{S}$、$K_{S}$ 有意义（例如 $O_{S}$ 是 $X$ 沿 $S$ 的完备化上的函数环），且环 $A_{S}$、$O_{S}$、$K_{S}$ 自然地组织成族（即显然可定义 $\Sigma$ 上三个环归纳概形，它们在 $S\in\Sigma$ 上的纤维分别为 $A_{S}$、$O_{S}$、$K_{S}$）。

易证这些态射 $f_{S}$ 构成族（即它们来自 $\Sigma$ 上群归纳概形之间的态射）。显然 $S\subset S'$ 时 $f_{S'}$ 到 $SO_{n}(A_{S})$ 的限制等于 $f_{S}$。4.3.12–4.3.13 将由此两点推出 $f_{S}=1$。

\SourcePage{157}
\bdp{4.3.12} 设 $Y$ 是 $\mathbb{C}$ 上有限型的分离概形，$R$ 是 $\mathbb{C}$-代数。置
\[
Y_{rat}(R)=\varinjlim_{U}\operatorname{Mor}(U,Y),
\]
其中极限取遍所有开 $U\subset X\otimes R$，使 $U$ 在 $\operatorname{Spec} R$ 的任一点上的纤维非空。换言之，$Y_{rat}(R)$ 的元素是由 $\operatorname{Spec} R$ 参数化的有理映射族 $X\to Y$。函子 $Y_{rat}$ 称为有理映射空间 $X\to Y$。易证 $Y_{rat}$ 是 fppf 拓扑的层，即 [LMB93] 意义下的「空间」。

有空间 $Y(A_{S})$（$S\in\Sigma$），它们构成族（即 $\Sigma$ 上有自然空间，其每个 $S$ 上的纤维等于 $Y(A_{S})$）。故 $Y_{rat}$ 上的正则函数定义 $Y(A_{S})$ 上的正则函数族 $f_{S}$（$S\in\Sigma$），使 $S\subset S'$ 时 $f_{S'}$ 到 $Y(A_{S})$ 的拉回等于 $f_{S}$。易见 $Y_{rat}$ 上的函数等同于具有此性质的函数族 $f_{S}$。

\bdp{4.3.13} \emph{命题.} 设 $G$ 是连通代数群。
\begin{enumerate}[label=(\roman*),leftmargin=2.2em]
\item $G_{rat}$ 上每个正则函数是常数。特别地，每个群态射 $G_{rat}\to\mathbb{G}_{m}$ 平凡。
\item 此外，对每个 $\mathbb{C}$-代数 $R$，$G_{rat}\otimes R$ 上每个正则函数是常数（即 $R$ 的元素）。
\end{enumerate}

\emph{证明.} 把 $G$ 表为 $\bigcup_{i}U_{i}$，其中 $U_{i}$ 是开集，同构于 $(\mathbb{A}^{1}\setminus\{0\})^{r}\times\mathbb{A}^{s}$（例如取 $U\subset G$ 为关于某 Borel 子群的大胞腔，则 $G$ 被有限多个形如 $gU$（$g\in G$）的集合覆盖）。有开覆盖 $G_{rat}=\bigcup_{i}(U_{i})_{rat}$，且 $(U_{i})_{rat}\cap(U_{j})_{rat}\neq\emptyset$。故只需对 $G=(\mathbb{G}_{m})^{r}\times(\mathbb{G}_{a})^{s}$ 证明命题。此外，只需对 $\mathbb{G}_{a}$ 与 $\mathbb{G}_{m}$ 证明 (ii)。

例如考虑 $\mathbb{G}_{m}$ 的情形。取 $X$ 上的丰沛线丛 $\mathcal{A}$，置 $V_{n}:=H^{0}(X,\mathcal{A}^{\otimes n})$，$V'_{n}:=V_{n}\setminus\{0\}$。定义 $\pi_{n}:V'_{n}\times V'_{n}\to(\mathbb{G}_{m})_{rat}$ 为 $(f,g)\mapsto f/g$。$(\mathbb{G}_{m})_{rat}\otimes R$ 上的正则函数 $\varphi$ 定义 $(V'_{n}\times V'_{n})\otimes R$ 上的正则函数 $\pi_{n}^{*}\varphi$，它关于 $\mathbb{G}_{m}$ 在 $V'_{n}\times V'_{n}$ 上的显然作用不变。$n$ 充分大时 $\dim V_{n}>1$，故 $\pi_{n}^{*}\varphi$ 延拓为

\SourcePage{158}
$(V_{n}\times V_{n})\otimes R$ 上 $\mathbb{G}_{m}$-不变的正则函数，它必为常数。故 $\varphi$ 为常数。

\bdp{4.3.14} 本小节后文不用（除 5.3.10 所需的 $\operatorname{GRAS}_{G}$ 的定义外）。

设 $G$ 是连通代数群。下述接近 $\operatorname{Bun}_{G}$ 的方式似乎自然。

记 $\operatorname{GRAS}_{G}$ 为 $X$ 上带有理截面的 $G$-torsor 构成的空间。其精确定义与 4.3.12 中 $Y_{rat}$ 的定义颇相似。我们愿称 $\operatorname{GRAS}_{G}$ 为对应 $G$ 与 $X$ 的\emph{大 Grassmann 流形}：因为对固定的有限 $S\subset X$，在 $X\setminus S$ 上平凡化的 $X$ 上 $G$-丛的空间可等同于归纳概形 $G(K_{S})/G(O_{S})=\prod_{x\in X}G(K_{x})/G(O_{x})$（见 5.3.10），而 $G(K_{x})/G(O_{x})$ 称为仿射 Grassmann 流形或环路 Grassmann 流形（见 4.5 或 [MV]）。

态射 $\pi:\operatorname{GRAS}_{G}\to\operatorname{Bun}_{G}$ 是光滑拓扑下的 $G_{rat}$-torsor（若 $G$ 的约化部分等于 $GL_{n}$、$SL_{n}$ 或 $Sp_{n}$，则存在光滑满射 $S\to\operatorname{Bun}_{G}$ 及截面 $S\to\operatorname{GRAS}_{G}$，这是显然的；一般 $G$ 可用 [DSim]）。

考虑函子
\begin{equation}\tag{200}
\pi^{*}:\operatorname{Vect}(\operatorname{Bun}_{G})\to\operatorname{Vect}(\operatorname{GRAS}_{G}),
\end{equation}
其中 $\operatorname{Vect}$ 表向量丛范畴。由 4.3.13，(200) 全忠实。可证对任意概形 $T$，$G_{rat}\times T$ 上每个向量丛都来自 $T$。这说明 (200) 是等价。

\emph{注.} (199) 的构造可解释如下：先构造 (199) 左右两端到 $\operatorname{GRAS}_{SO_{n}}$ 的拉回之间的同构，再用 (200) 的全忠实性。其实不必用同构 (155)。故 (199) 的构造在 $SO_{2}$ 的情形也有效。

\SourcePage{159}
\bdp{4.3.15} 证明命题 4.3.10。各同构 (199) 彼此相容（用 4.3.9 的唯一性）。剩下要证 (199) 的张量平方等于复合
\begin{equation}\tag{201}
\operatorname{Pf}_{\mathcal{L}}^{\otimes 2}\xrightarrow{\ \sim\ }\nu_{\mathcal{L}}\xrightarrow{\ \sim\ }\nu\xrightarrow{\ \sim\ }\nu_{\mathcal{L}'}\xrightarrow{\ \sim\ }\operatorname{Pf}_{\mathcal{L}'}^{\otimes 2},
\end{equation}
其中 $\nu_{\mathcal{L}}$ 由 (188) 定义。

固定 $X$ 上的 $SO_{n}$-torsor $\mathcal{F}$ 及其在某非空有限 $S\subset X$ 上的 $X\setminus S$ 平凡化。用该平凡化计算 (199) 与 (201) 诱导的同构 $\operatorname{Pf}_{\mathcal{L},\mathcal{F}}^{\otimes 2}\xrightarrow{\ \sim\ }\operatorname{Pf}_{\mathcal{L}',\mathcal{F}}^{\otimes 2}$。

回忆 $\operatorname{Pf}_{\mathcal{L},\mathcal{F}}:=\operatorname{Pf}(W_{\mathcal{F}}\otimes\mathcal{L})\otimes\operatorname{Pf}(W\otimes\mathcal{L})^{\otimes-1}$。由 4.2.13，
\[
\operatorname{Pf}(W_{\mathcal{F}}\otimes\mathcal{L})=\operatorname{Pf}(V;L_{+},L_{-}),\qquad
\operatorname{Pf}(W\otimes\mathcal{L})=\operatorname{Pf}(V;L_{+}^{0},L_{-}),
\]
其中 $V=\mathcal{L}_{K_{S}}\otimes W$，$L_{-}=\Gamma(X\setminus S,\mathcal{L}\otimes W)$，$L_{+}^{0}=\mathcal{L}_{O_{S}}\otimes W$，而 $L_{+}$ 是 $W_{\mathcal{F}}\otimes\mathcal{L}$ 到 $\operatorname{Spec} O_{S}$ 的拉回的截面空间（用 4.3.6 的记号）。用 (166) 得
\begin{equation}\tag{202}
\operatorname{Pf}_{\mathcal{L},\mathcal{F}}=M^{L_{+}}\otimes(M^{L_{+}^{0}})^{*},
\end{equation}
其中 $M$ 是 $\operatorname{Cl}(V)$ 上不可约的 $\mathbb{Z}/2\mathbb{Z}$-分次离散模。$\operatorname{Pf}_{\mathcal{L}',\mathcal{F}}$ 有类似的描述，用 $V'$、$L'_{+}$、$(L_{+}^{0})'$、$L'_{-}$，其中 $V'=\mathcal{L}'_{K_{S}}\otimes W$，等等。固定 4.3.8 的 $\mu_{2}$-torsor $\mathcal{E}$ 在 $S$ 上的平凡化。它给出 $\mathcal{E}$ 在 $\operatorname{Spec} O_{S}$ 上的平凡化，从而给出等同
\begin{equation}\tag{203}
(V',L'_{+},(L_{+}^{0})')=(V,L_{+},L_{+}^{0}).
\end{equation}
因 (202) 不含 $L_{-}$，得同构 $\operatorname{Pf}_{\mathcal{L},\mathcal{F}}\xrightarrow{\ \sim\ }\operatorname{Pf}_{\mathcal{L}',\mathcal{F}}$。易证它与 (199) 诱导者一致（注意 $\mathcal{F}$ 在 $X\setminus S$ 上的平凡化与 $\mathcal{E}$ 在 $S$ 上的平凡化给出 $\ell^{\operatorname{Spin}}_{\mathcal{E}}$ 在 $\mathcal{F}$ 上的平凡化，因为 (150) 左端有特异元素）。

现在须证 (201) 诱导的同构 $\operatorname{Pf}_{\mathcal{L},\mathcal{F}}^{\otimes 2}\xrightarrow{\ \sim\ }\operatorname{Pf}_{\mathcal{L}',\mathcal{F}}^{\otimes 2}$，在把 $\operatorname{Pf}_{\mathcal{L},\mathcal{F}}$、$\operatorname{Pf}_{\mathcal{L}',\mathcal{F}}$ 与 (202) 右端等同后是恒等。

\SourcePage{160}
$\mathcal{F}$ 在 $X\setminus S$ 上的平凡化给出同构 $\nu_{\mathcal{L}}(\mathcal{F})\xrightarrow{\ \sim\ }d(L_{+}^{0},L_{+})$，其中 $d(L_{+}^{0},L_{+})$ 是相对行列式，即对任意 c-格 $U\subset L\cap L_{+}^{0}$ 有 $d(L_{+}^{0},L_{+})=\det(L_{+}/U)\otimes\det(L_{+}^{0}/U)^{\otimes-1}$。类似有 $\nu_{\mathcal{L}'}(\mathcal{F})=d((L_{+}^{0})',L'_{+})$。(201) 中的同构 $\nu_{\mathcal{L}}(\mathcal{F})\xrightarrow{\ \sim\ }\nu_{\mathcal{L}'}(\mathcal{F})$ 在 [Del87] 中如下定义：任取 $\mathcal{L}$ 与 $\mathcal{L}'$ 到 $\operatorname{Spec} O_{S}$ 的拉回之间的同构 $f$。$f$ 给出同构 $f_{*}:(V,L_{+},L_{+}^{0})\xrightarrow{\ \sim\ }(V',L'_{+},(L_{+}^{0})')$，从而给出 $d(L_{+}^{0},L_{+})\xrightarrow{\ \sim\ }d(L'_{+},(L_{+}^{0})')$，它其实不依赖 $f$ 的选取。用 $\mu_{2}$-torsor $\mathcal{E}=\mathcal{L}'\otimes\mathcal{L}^{\otimes-1}$ 在 $\operatorname{Spec} O_{S}$ 上的上述平凡化来定义 $f$ 较方便。此时 $f_{*}$ 与 (203) 一致。

于是我们把 $\nu_{\mathcal{L}}(\mathcal{F})$ 与 $\nu_{\mathcal{L}'}(\mathcal{F})$ 都等同于 $d(L_{+}^{0},L_{+})$，使 (201) 中的同构成为恒等。又把 $\operatorname{Pf}_{\mathcal{L},\mathcal{F}}$ 与 $\operatorname{Pf}_{\mathcal{L}',\mathcal{F}}$ 都等同于 (202) 右端。剩下要证 (187) 的同构与 $\mathcal{L}'$ 的类似物诱导同一同构
\begin{equation}\tag{204}
(M^{L_{+}}\otimes(M^{L_{+}^{0}})^{*})^{\otimes 2}\xrightarrow{\ \sim\ }d(L_{+}^{0},L_{+}).
\end{equation}

据 4.2.8 与 4.2.13，(187) 诱导的 (204) 可描述如下。有典范同构
\begin{equation}\tag{205}
N^{L_{+}\otimes H}\otimes(N^{L_{+}^{0}\otimes H})^{*}\xrightarrow{\ \sim\ }d(L_{+}^{0},L),
\end{equation}
其中 $N$ 是 Clifford 代数 $\operatorname{Cl}(V\oplus V^{*})=\operatorname{Cl}(V\oplus V)=\operatorname{Cl}(V\otimes H)$ 上不可约的 $(\mathbb{Z}/2\mathbb{Z})$-分次离散模，$H$ 表 $\mathbb{C}^{2}$ 配双线性型 $\bigl(\begin{smallmatrix}0&1\\1&0\end{smallmatrix}\bigr)$（构造 (205) 时取 $N$ 为 (182) 的右端）。
% [Tutor Note: 原书 p.160 (205) 右端作 d(L_+^0,L)，而下文的 (204) 作 d(L_+^0,L_+)；此处照录原书。]
另一方面，$P:=M\otimes M$ 是 $\operatorname{Cl}(V)\otimes\operatorname{Cl}(V)=\operatorname{Cl}(V\otimes H'')$ 上不可约的 $(\mathbb{Z}/2\mathbb{Z})$-分次离散模，其中 $H''$ 表 $\mathbb{C}^{2}$ 配双线性型 $\bigl(\begin{smallmatrix}1&0\\0&1\end{smallmatrix}\bigr)$。把 (204) 左端改写为 $P^{L_{+}\otimes H''}\otimes(P^{L_{+}^{0}\otimes H''})^{*}$。于是正交同构 $\psi:H''\xrightarrow{\ \sim\ }H$ 诱导 (204)。要使 (204) 是 (187) 诱导者，须用 $\det\psi=i$ 规范化 $\psi$（还是 $-i$？？我们应当查一下！）。
% [Tutor Note: 原书 p.160 此句末尾的括注「还是 -i（原文用减号）？？我们应当查一下！」是作者在底本中的自注，照录。]

\SourcePage{161}
因上述描述不含 $L_{-}$，$\mathcal{L}'$ 的 (187) 类似物诱导同一同构 (204)，证毕。

\bdp{4.3.16} 本小节与 4.3.17 将用于 4.4.14（水平性定理 2.7.3 证明的末段）与定理 5.4.5（本工作的主要结果）的证明。读者暂时可以跳过。

照例置 $O:=\mathbb{C}[[t]]$，$K:=\mathbb{C}((t))$。固定 $\mathcal{L}\in\omega^{1/2}(X)$，即 $\mathcal{L}$ 是 $\omega_{X}$ 的平方根。再固定 $\omega_{O}$ 的一个平方根，记为 $\omega_{O}^{1/2}$。于是可定义 2.6.5 中概形 $X^{\wedge}$ 的二重覆叠 $X_{2}^{\wedge}$。回忆 $X^{\wedge}$ 的 $R$-点是一个 $R$-态射 $\alpha:\operatorname{Spec}(R\hat{\otimes}O)\to X\otimes R$，其微分在 $\operatorname{Spec} R$ 的任一点都不为零。记 $\mathcal{L}_{R}$ 为 $\mathcal{L}$ 到 $X\otimes R$ 的拉回。按定义，$X_{2}^{\wedge}(R)$ 在 $\gamma\in X^{\wedge}(R)$ 上的纤维，是 $R\hat{\otimes}\omega_{O}$ 的平方根群胚中同构 $H^{0}(\operatorname{Spec} R\hat{\otimes}O,\alpha^{*}\mathcal{L}_{R})\xrightarrow{\ \sim\ }R\hat{\otimes}\omega_{O}^{1/2})$ 的集合。
% [Tutor Note: 原书 p.161 此行末尾 $\omega_O^{1/2}$ 之后多一个右括号，括号不配对，属排版笔误，照录。]

3.5.2 引入的群归纳概形 $\operatorname{Aut}_{2}O:=\operatorname{Aut}(O,\omega_{O}^{1/2})$ 以结构搬运的方式作用在 $X_{2}^{\wedge}$ 上。

设 $M$ 是 2.8.1 中 $G=SO(W)=SO_{n}$ 这一特殊情形的概形。记 $M_{2}^{\wedge}$ 为 $M$ 与 $X_{2}^{\wedge}$ 在 $X$ 上的纤维积（故 $M_{2}^{\wedge}$ 是 2.8.3 中概形 $M^{\wedge}$ 的二重覆叠）。则半直积 $\operatorname{Aut}_{2}O\ltimes SO_{n}(K)$ 作用在 $M_{2}^{\wedge}$ 上。事实上 $M_{2}^{\wedge}$ 是 $M^{\wedge}$ 与 $X_{2}^{\wedge}$ 在 $X^{\wedge}$ 上的纤维积，而 $\operatorname{Aut}_{2}O\ltimes SO_{n}(K)$ 作用在图表
\[
\begin{tikzcd}
M^{\wedge} \arrow[d]\\
X_{2}^{\wedge} \arrow[r] & X^{\wedge}
\end{tikzcd}
\]
上（$\operatorname{Aut}O\ltimes SO_{n}(K)$ 在 $M^{\wedge}$ 上的作用在 2.8.4 中定义；$\operatorname{Aut}_{2}O\ltimes SO_{n}(K)$ 经其商 $\operatorname{Aut}_{2}O$ 与 $\operatorname{Aut}O$ 作用在 $X_{2}^{\wedge}$ 与 $X^{\wedge}$ 上）。

记 $\operatorname{Pf}_{\mathcal{L}}^{\wedge}$ 为 4.3.1 定义的 $\operatorname{Bun}_{SO_{n}}$ 上线丛 $\operatorname{Pf}_{\mathcal{L}}$ 到 $M_{2}^{\wedge}$ 的拉回。要把 $\operatorname{Aut}_{2}O\ltimes SO_{n}(K)$ 在 $M_{2}^{\wedge}$ 上的作用提升为 $\operatorname{Aut}_{2}O\ltimes\tilde{SO}_{n}(K)$ 在 $\operatorname{Pf}_{\mathcal{L}}^{\wedge}$ 上的作用，其中 $\tilde{SO}_{n}(K)$ 是中心扩张
% 原书 p.161 末句未完，续 p.162（下一单元）。
% ===================== end 源码/units/ch04_u04.tex =====================

% ===================== begin 源码/units/ch04_u05.tex =====================
% ch04_u05 —— 原书 p.162–p.171（4.3.17–4.4.16）
% 承接 p.161 末句：……其中 $\tilde{SO}_n(K)$ 是中心扩张
\SourcePage{162}

(196)（对应于 $\omega_{O}^{1/2}$）。$\operatorname{Aut}_{2}O$ 在 $\operatorname{Pf}_{\mathcal{L}}^{\wedge}$ 上的作用是清楚的：$\operatorname{Aut}_{2}O$ 作用在概形 $M_{2}^{\wedge}$ 上，而 $M_{2}^{\wedge}$ 看作 $\operatorname{Bun}_{SO_{n}}$ 上的概形。另一方面，$\tilde{SO}_{n}(K)$ 作用在 $\operatorname{Pf}_{\mathcal{L},\hat x}$ 上，后者定义为 $\operatorname{Pf}_{\mathcal{L}}^{\wedge}$ 在 $M_{2}^{\wedge}$ 于 $\hat x\in X_{2}^{\wedge}$ 处纤维上的限制。事实上，该纤维等于 $\operatorname{Bun}_{SO_{n},\underline{x}}$，其中 $x$ 是 $\hat x$ 在 $X$ 中的像；由 4.3.7，中心扩张 $\tilde{SO}_{n}(K_{x})_{\mathcal{L}}$ 作用在 $\operatorname{Pf}_{\mathcal{L}}$ 到 $\operatorname{Bun}_{SO_{n},\underline{x}}$ 的拉回上。该扩张只依赖 $\mathcal{L}_{x}$，即 $\mathcal{L}$ 到 $\operatorname{Spec}O_{x}$ 的拉回。因 $\hat x$ 给出 $(O,\omega_{O}^{1/2})$ 与 $(O_{x},H^{0}(\operatorname{Spec}O_{x},\mathcal{L}_{x}))$ 之间的同构，我们得到同构 $\tilde{SO}_{n}(K_{x})_{\mathcal{L}}\xrightarrow{\ \sim\ }\tilde{SO}_{n}(K)$，从而得到所要的 $\tilde{SO}_{n}(K)$ 作用。

\bdp{4.3.17} \emph{命题.}
\begin{enumerate}[label=(\roman*),leftmargin=2.2em]
\item $\tilde{SO}_{n}(K)$ 在 $\operatorname{Pf}_{\mathcal{L},\hat x}$（$\hat x\in X_{2}^{\wedge}$）上的作用，来自 $\tilde{SO}_{n}(K)$ 在 $\operatorname{Pf}_{\mathcal{L}}^{\wedge}$ 上的（显然唯一的）作用。
\item $\operatorname{Aut}_{2}O$ 与 $\tilde{SO}_{n}(K)$ 在 $\operatorname{Pf}_{\mathcal{L}}^{\wedge}$ 上的作用，定义了 $\operatorname{Aut}_{2}O\ltimes\tilde{SO}_{n}(K)$ 的一个作用。
\end{enumerate}

\emph{注.} 陈述 (ii) 可按 2.8.2 的精神来理解：$\operatorname{Aut}_{2}O$ 的作用给出 $\pi^{*}\operatorname{Pf}_{\mathcal{L}}$ 上沿 $X$ 的联络，其中 $\pi$ 是态射 $M\to\operatorname{Bun}_{G}$。而 $\operatorname{Aut}_{2}O$ 的作用与 $\tilde{SO}_{n}(K)$ 的作用相容，意思是某个中心扩张 $\tilde{J}^{\mathrm{mer}}(SO_{n})_{\mathcal{L}}$ 在 $\pi^{*}\operatorname{Pf}_{\mathcal{L}}$ 上的作用是水平的。

\emph{证明.} 要定义 $\operatorname{Aut}_{2}O\ltimes\tilde{SO}_{n}(K)$ 在 $\operatorname{Pf}_{\mathcal{L}}^{\wedge}$ 上具有所需性质的作用，我们按 4.3.7 的办法进行。设 $R$ 是 $\mathbb{C}$-代数。取 $M_{2}^{\wedge}$ 的一个 $R$-点 $u$，以及 $\operatorname{Aut}_{2}O\ltimes\tilde{SO}_{n}(K)$ 的一个 $R$-点 $\tilde{g}$。回忆 $SO_{n}$ 是 $SO(W)$ 的缩写。记 $\mathcal{F}$ 与 $\mathcal{F}'$ 为 $X\otimes R$ 上对应 $u$ 与 $gu$ 的 $SO(W)$-torsor，其中 $g$ 是 $\tilde{g}$ 在 $\operatorname{Aut}_{2}O\ltimes SO_{n}(K)$ 中的像。我们须定义一个同构
\begin{equation}\tag{206}
\operatorname{Pf}(W_{\mathcal{F}}\otimes\mathcal{L}_{R})\xrightarrow{\ \sim\ }\operatorname{Pf}(W_{\mathcal{F}'}\otimes\mathcal{L}_{R})
\end{equation}
其中 $\mathcal{L}_{R}$ 是 $\mathcal{L}$ 到 $X\otimes R$ 的拉回。

\SourcePage{163}

置 $V:=\omega_{O}^{1/2}\otimes_{O}K\otimes W$。这是 $\mathbb{C}$ 上的 Tate 空间，带有非退化对称双线性型（见 4.3.3）。由 4.2.14，
\begin{equation}\tag{207}
\operatorname{Pf}(W_{\mathcal{F}}\otimes\mathcal{L}_{R})=\operatorname{Pf}(V\hat\otimes R;L_{+}\hat\otimes R,L^{u}_{-})
\end{equation}
其中 $L_{+}:=\omega_{O}^{1/2}\otimes W\subset V$（故 $L_{+}$ 是 $V$ 中的 Lagrange c-格），而 Lagrange d-格 $L^{u}_{-}\subset V\hat\otimes R$ 定义如下。点 $u\in M_{2}^{\wedge}(R)$ 是四元组 $(\alpha,\mathcal{F},\gamma,f)$，其中 $\alpha$、$\mathcal{F}$、$\gamma$ 的含义同 2.8.4（在我们这个特殊情形 $G=SO(W)$ 中），而 $f$ 是 $R\hat\otimes\omega_{O}$ 的平方根构成的群胚中、$H^{0}(\operatorname{Spec}R\hat\otimes K,\alpha^{*}\mathcal{L}_{R})$ 与 $R\hat\otimes\omega_{O}^{1/2}$ 之间的一个同构。记 $\Gamma_{\alpha}$ 的含义同 2.8.4。则
\[
L^{u}_{-}:=H^{0}((X\otimes R)\setminus\Gamma_{\alpha},W_{\mathcal{F}}\otimes\mathcal{L}_{R})\subset H^{0}(\operatorname{Spec}R\hat\otimes K,\alpha^{*}W_{\mathcal{F}}\otimes\alpha^{*}\mathcal{L}_{R})\xrightarrow[\ \sim\ ]{\varphi}V\hat\otimes R
\]
（同构 $\varphi$ 由 $\gamma$ 与 $f$ 诱导）。

计入 (207)，可见构造 (206) 等价于定义一个同构
\begin{equation}\tag{208}
\operatorname{Pf}(V\hat\otimes R;L_{+}\hat\otimes R,L^{u}_{-})\xrightarrow{\ \sim\ }\operatorname{Pf}(V\hat\otimes R;L_{+}\hat\otimes R,L^{gu}_{-}).
\end{equation}
群归纳概形 $\operatorname{Aut}_{2}O\ltimes SO(W\otimes K)$ 以显然的方式作用在 $V$ 上，且容易看出 $L^{gu}_{-}=gL^{u}_{-}$。由 (166)，(208) 的左端等于 $(M\otimes R)_{L_{-}}$ 的逆，其中 $M$ 是 Clifford 模 $Cl(V)/Cl(V)L_{+}$，$L_{-}:=L^{u}_{-}$。于是只剩构造同构 $(M\otimes R)_{L_{-}}\xrightarrow{\ \sim\ }(M\otimes R)_{gL_{-}}$。我们把它定义为 $\tilde{g}$ 在 $M\otimes R$ 上的作用\footnote{回忆 $g$ 是 $\operatorname{Aut}_{2}O\ltimes\tilde{SO}_{n}(K)=\operatorname{Aut}_{2}O\ltimes\tilde{SO}(W\otimes K)$ 的一个 $R$-点。按 $\tilde{SO}_{n}(K)$ 的定义，它作用在 $M$ 上。群归纳概形 $\operatorname{Aut}_{2}O$ 作用在 $(V,L_{+})$ 上，从而作用在 $M$ 上。}所诱导的同构。\hfill$\Box$

\bdsec{4.4. $\operatorname{Bun}_G$ 上的半形式}

\bdp{4.4.1} 设 $G$ 半单。固定 $\mathfrak{g}$ 上一个 $G$-不变的非退化对称双线性型。置 $n:=\dim\mathfrak{g}$，并把 $SO(\mathfrak{g})$ 写成 $SO_{n}$。伴随表示 $G\to SO(\mathfrak{g})$ 诱导态射 $f:\operatorname{Bun}_{G}\to\operatorname{Bun}_{SO_{n}}$。对 $\mathcal{L}\in\omega^{1/2}(X)$ 置 $\lambda'_{\mathcal{L}}:=f^{*}\operatorname{Pf}_{\mathcal{L}}$，其中 $\operatorname{Pf}_{\mathcal{L}}$ 是 4.3.1 中的线丛；

\SourcePage{164}

故 $\lambda'_{\mathcal{L}}$ 在 $\mathfrak{F}\in\operatorname{Bun}_{G}$ 处的纤维等于 $\operatorname{Pf}(\mathfrak{g}_{\mathfrak{F}}\otimes\mathcal{L})\otimes\operatorname{Pf}(\mathfrak{g}\otimes\mathcal{L})^{\otimes-1}$。同构 (189) 诱导一个同构
\begin{equation}\tag{209}
(\lambda'_{\mathcal{L}})^{\otimes 2}=\omega^{\sharp}_{\operatorname{Bun}_{G}}
\end{equation}
这里 $\omega^{\sharp}_{\operatorname{Bun}_{G}}$ 是正规化典范丛 (146)；按 2.1.1，它在 $\mathfrak{F}\in\operatorname{Bun}_{G}$ 处的纤维等于 $\det R\Gamma(X,\mathfrak{g}_{\mathfrak{F}})\otimes(\det R\Gamma(X,\mathfrak{g}\otimes O_{X}))^{\otimes-1}$。

\bdp{4.4.2} 考虑函子
\begin{equation}\tag{210}
\lambda':\omega^{1/2}(X)\to(\omega^{\sharp})^{1/2}(\operatorname{Bun}_{G}),
\end{equation}
$\mathcal{L}\mapsto\lambda'_{\mathcal{L}}$。由 4.3.10，$\lambda'$ 关于 Picard 函子 $\tilde\ell':\mu_{2}\operatorname{tors}(X)\to\mu_{2}\operatorname{tors}(\operatorname{Bun}_{G})$ 是仿射的；该函子把 $X$ 上的 $\mu_{2}$-torsor $\mathcal{E}$ 送到 $\tilde\ell'_{\mathcal{E}}$，即 $\operatorname{Bun}_{SO_{n}}$ 上的 torsor $\ell^{\operatorname{Spin}}_{\mathcal{E}}$ 到 $\operatorname{Bun}_{G}$ 的拉回。

\bdp{4.4.3} \emph{命题.} $\tilde\ell'=\ell'$，其中 $\ell'$ 是下述两个函子的复合：由 (56) 诱导的函子 $\mu_{2}\operatorname{tors}(X)\to Z\operatorname{tors}(X)$，以及 4.1.1–4.1.4 中构造的函子 $\ell:Z\operatorname{tors}(X)\to\mu_{\infty}\operatorname{tors}(\operatorname{Bun}_{G})$。这里 $Z=\pi_{1}(G)^{\vee}=$ ${}^{L}G$ 的中心（见 4.1.1 的注）。

假定该命题，我们定义典范的 $\ell$-仿射函子
\begin{equation}\tag{211}
\lambda:Z\operatorname{tors}_{\theta}(X)\to\mu_{\infty}\operatorname{tors}_{\theta}(\operatorname{Bun}_{G})
\end{equation}
为 $\mathcal{E}\cdot\mathcal{L}\mapsto\lambda_{\mathcal{E}\cdot\mathcal{L}}:=\ell_{\mathcal{E}}\cdot\lambda'_{\mathcal{L}}$，其中 $\mathcal{E}\in Z\operatorname{tors}(X)$，$\mathcal{L}\in\omega^{1/2}(X)$。（注意：正规化问题！！！？？？）
% [Tutor Note: 原书 p.164 此括注「Attention: normalization problem!!!???」是作者本人的自注，连用感叹号与问号，底本如此，保留原样。]

为证明命题 4.4.3，注意 $\tilde\ell'$ 是对应于 $G$ 被 $\mu_{2}$ 的扩张的函子 (152)，该扩张由 $SO(\mathfrak{g})$ 的旋量扩张诱导。因此 $\tilde\ell'$ 是 $\ell:Z\operatorname{tors}(X)\to\mu_{\infty}\operatorname{tors}(\operatorname{Bun}_{G})$ 与函子 $\mu_{2}\operatorname{tors}(X)\to Z\operatorname{tors}(X)$ 的复合，后者由态射 $\mu_{2}\to Z=\pi_{1}(G)^{\vee}$ 诱导，而该态射对偶于 $\pi_{1}(G)\to\pi_{1}(SO(\mathfrak{g}))=\mathbb{Z}/2\mathbb{Z}$。于是只需证明下面的事实。

\SourcePage{165}

\bdp{4.4.4} \emph{引理.} 态射 $\pi_{1}(G)\to\pi_{1}(SO(\mathfrak{g}))=\mathbb{Z}/2\mathbb{Z}$ 对偶于群 ${}^{L}G$ 的态射 (56)。

\emph{证明.} 有典范同构 $f:P/P_{G}\xrightarrow{\ \sim\ }\operatorname{Hom}(\pi_{1}(G)(1),\mu_{\infty})$，其中 $P_{G}$ 是 $G$ 的权群，$P$ 是其泛覆叠 $\tilde{G}$ 的权群；权 $\lambda\in P$ 是嘉当子群 $\tilde{H}\subset\tilde{G}$ 的特征标，而 $f(\lambda)$ 是它在 $\pi_{1}(G)(1)\subset\tilde{H}$ 上的限制。设 $M$ 是 $so(\mathfrak{g})$ 的旋量表示。则 $\tilde{G}$ 作用在 $M$ 上，且 $\pi_{1}(G)(1)\subset\tilde{G}$ 按某个特征标 $\chi\in\operatorname{Hom}(\pi_{1}(G)(1),\mu_{\infty})$ 作用。按 (56) 的定义（另见 3.4.1 中 $\lambda^{\#}$ 的定义），引理只是说 $\chi=f(\rho)$，其中 $\rho\in P$ 是基本权之和。

设 $\mathfrak{b}\subset\mathfrak{g}$ 是 Borel 子代数。选取 $\mathfrak{b}$-不变的旗 $0\subset V_{1}\subset\ldots\subset V_{n}=\mathfrak{g}$，使 $\dim V_{k}=k$、$V_{k}^{\perp}=V_{n-k}$，且 $\mathfrak{b}$ 是其中某个 $V_{k}$。设 $\mathfrak{b}'$ 是该旗在 $so(\mathfrak{g})$ 中的稳定化子。这是 $so(\mathfrak{g})$ 中含 $\mathfrak{b}$ 的 Borel 子代数。设 $m\in M$ 是关于 $\mathfrak{b}'$ 的最高向量。则 $\mathbb{C}m$ 是 $\mathfrak{b}$-不变的，且相应的 $\mathfrak{b}$ 的特征标等于正根之和的一半，即 $\rho$。故 $\chi=f(\rho)$。\hfill$\Box$

\emph{注.} 按 Kostant（参看 [Ko61] 中引理 5.9 的证明），$\mathfrak{g}$-模 $M$ 同构于 $2^{[r/2]}$ 个最高权为 $\rho$ 的不可约 $\mathfrak{g}$-模之和（其中 $r$ 是 $\mathfrak{g}$ 的秩）。

\bdp{4.4.5} 我们对 (211) 的构造略依赖 $\mathfrak{g}$ 上标量积的选取（见 4.4.1）。由于 $\mathfrak{g}$ 上有若干「典范」标量积，读者或许更喜欢 (211) 的如下版本。

为简化记号，设 $G$ 单。则 $\mathfrak{g}$ 上不变对称双线性型的空间是一维的。记之为 $\beta$。取 $\beta$ 的一个平方根，即一维向量空间 $\beta^{1/2}$，连同一个同构 $\beta^{1/2}\otimes\beta^{1/2}\xrightarrow{\ \sim\ }\beta$。于是 $\mathfrak{g}\otimes\beta^{1/2}$ 带有典范双线性型。考虑表示 $G\to SO(\mathfrak{g}\otimes\beta^{1/2})$，然后按 4.4.1–4.4.3 进行（例如，此时 $\lambda'_{\mathcal{L}}$ 在 $\mathcal{F}\in\operatorname{Bun}_{G}$ 处的纤维等于 $\operatorname{Pf}(\mathfrak{g}_{\mathcal{F}}\otimes\mathcal{L}\otimes\beta^{1/2})\otimes\operatorname{Pf}(\mathfrak{g}\otimes\mathcal{L}\otimes\beta^{1/2})^{\otimes-1}$）。这样得到的函子 (211)

\SourcePage{166}

略依赖 $\beta^{1/2}$ 的选取。更确切地说，$-1\in\operatorname{Aut}\beta^{1/2}$ 作用在 $\lambda'_{\mathcal{L}}$ 上，从而作用在 $\lambda_{\mathcal{M}}$（$\mathcal{M}\in Z\operatorname{tors}_{\theta}(X)$）上，其方式为乘以 $(-1)^{p}$，其中 $p:\operatorname{Bun}_{G}\to\mathbb{Z}/2\mathbb{Z}$ 是复合
\[
\operatorname{Bun}_{G}\to\pi_{0}(\operatorname{Bun}_{G})=\pi_{1}(G)\to\pi_{1}(SO(\mathfrak{g}))=\mathbb{Z}/2\mathbb{Z}.
\]
我们是否该把 $\lambda_{\mathcal{M}}$ 看作一个 SUPER-层？？？！
% [Tutor Note: 原书 p.166「Do we want to consider $\lambda_{\mathcal{M}}$ as a SUPER-sheaf???!」是作者行文中的自问，连用问号与感叹号，底本如此，保留原样。]

\bdp{4.4.6} 我们已把 $\mathcal{L}\in Z\operatorname{tors}_{\theta}(X)$ 对应到 $\operatorname{Bun}_{G}$ 上的线丛 $\lambda_{\mathcal{L}}$（见 4.4.1–4.4.3）。对 $x\in X$，记 $\lambda_{\mathcal{L},\underline{x}}$ 为 $\lambda_{\mathcal{L}}$ 到 $\operatorname{Bun}_{G,\underline{x}}$ 的拉回。在 4.4.7–4.4.10 中我们将定义 $G(K_{x})$ 的中心扩张 $\tilde{G}(K_{x})_{\mathcal{L}}$，它作用在 $\lambda_{\mathcal{L},\underline{x}}$ 上。在 4.4.11–4.4.13 中我们考虑 $\tilde{G}(K_{x})_{\mathcal{L}}$ 的李代数。

\bdp{4.4.7} 设 $O$、$K$ 与 $\omega_{O}$ 的含义同 4.3.3。固定 $\omega_{O}$ 的一个平方根 $\mathcal{L}$。然后我们构造群归纳概形的中心扩张
\begin{equation}\tag{212}
0\to\mathbb{G}_{m}\to\tilde{G}(K)_{\mathcal{L}}\to G(K)\to 0
\end{equation}
如下。$\mathcal{L}$ 定义了中心扩张 (196)。固定 $\mathfrak{g}$ 上一个非退化不变对称双线性型\footnote{不固定 $\mathfrak{g}$ 上的型，读者也可按 4.4.5 处理。}，并把 $SO(\mathfrak{g})$ 写成 $SO_{n}$，$n:=\dim\mathfrak{g}$。我们把 (212) 定义为 $G(K)$ 的中心扩张，它与由 (196) 经伴随表示 $G\to SO(\mathfrak{g})=SO_{n}$ 诱导的扩张相反。扩张 (212) 在 $G(O)$ 上分裂。

\emph{注.} 在 $G=SO_{r}$ 的情形，我们的记号有歧义：$\tilde{G}(K)\ne\tilde{SO}_{r}(K)$。但愿这一歧义无害。

\bdp{4.4.8} 设 $\mathcal{L}\in\omega^{1/2}(X)$，$x\in X$。按 4.4.7，$\mathcal{L}$ 到 $\operatorname{Spec}O_{x}$ 的限制定义了 $G(K_{x})$ 的中心扩张，记之为 $\tilde{G}(K_{x})_{\mathcal{L}}$。记 $\lambda'_{\mathcal{L},\underline{x}}$ 为线丛 $\lambda'_{\mathcal{L}}$（来自 4.4.1）到 $\operatorname{Bun}_{G,\underline{x}}$ 的拉回。由 4.3.7，$G(K_{x})$ 在 $\operatorname{Bun}_{G,\underline{x}}$ 上的作用提升为 $\tilde{G}(K_{x})_{\mathcal{L}}$ 在 $\lambda'_{\mathcal{L}}$ 上的典范作用。
% [Tutor Note: 原书 p.166 末三处均印作 $\lambda'_{\mathcal{L}}$（无逗号与下划线 x 的下标），而同段前文定义的是 $\lambda'_{\mathcal{L},\underline{x}}$；底本如此，保留原样。]

子群 $\mathbb{G}_{m}\subset\tilde{G}(K_{x})_{\mathcal{L}}$ 以显然的方式作用在 $\lambda'_{\mathcal{L}}$ 上（见 4.4.7 中 $\tilde{G}(K_{x})_{\mathcal{L}}$ 的定义，以及 4.3.7 的最后

\SourcePage{167}

一句）。$G(O_{x})\subset\tilde{G}(K_{x})_{\mathcal{L}}$ 在 $\lambda'_{\mathcal{L},\underline{x}}$ 上的作用是显然的那个。

\bdp{4.4.9} 在 4.4.7 中我们定义了一个函子
\begin{equation}\tag{213}
\omega^{1/2}(O)\to\{G(K)\text{ 被 }\mathbb{G}_{m}\text{ 的中心扩张}\}
\end{equation}
其中 $\omega^{1/2}(O)$ 是 $\omega_{O}$ 的平方根构成的群胚。(213) 的左端是 3.4.4 意义下的 $\mu_{2}$-范畴。(213) 的右端是 $Z$-范畴，$Z:=\pi_{1}(G)^{\vee}=\operatorname{Hom}(\pi_{1}(G),\mathbb{G}_{m})$。事实上，边缘态射\footnote{先验地，(214) 是抽象群的态射；但按 4.1.7 的注，它实际上是群归纳概形的态射。另见 4.5.4。}
\begin{equation}\tag{214}
G(K)\to H^{1}(K,\pi_{1}^{\mathrm{et}}(G))=\pi_{1}(G)=Z^{\vee}
\end{equation}
诱导一个态射\footnote{事实上是同构（见 4.5.4）。}
\begin{equation}\tag{215}
Z\to\operatorname{Hom}(G(K),\mathbb{G}_{m}),
\end{equation}
即 (213) 右端上的一个 $Z$-结构。用 (56) 定义的态射 $\mu_{2}\to Z$，我们把 (213) 的右端看作 $\mu_{2}$-范畴。于是 (213) 是 $\mu_{2}$-函子（用 4.3.4、4.3.4 的注 (ii) 以及 4.4.4）。故由 3.4.4，函子 (213) 给出一个 $Z$-函子
\begin{equation}\tag{216}
Z\operatorname{tors}_{\theta}(O)\to\{G(K)\text{ 被 }\mathbb{G}_{m}\text{ 的中心扩张}\}.
\end{equation}
由 (213) 给出的、对应于 $\mathcal{L}\in Z\operatorname{tors}_{\theta}(O)$ 的 $G(K)$ 的中心扩张将记作 $\tilde{G}(K)_{\mathcal{L}}$。扩张
\begin{equation}\tag{217}
0\to\mathbb{G}_{m}\to\tilde{G}(K)_{\mathcal{L}}\to G(K)\to 0
\end{equation}
在 $G(O)$ 上分裂。

\emph{注记.}
\begin{enumerate}[label=(\roman*),leftmargin=2.2em]
\item 按 3.4.7 (i)，(213) 右端上的 $Z$-结构给出一个 Picard 函子
\begin{equation}\tag{218}
Z\operatorname{tors}(O)=Z\operatorname{tors}\to\{G(K)\text{ 被 }\mathbb{G}_{m}\text{ 的中心扩张}\}.
\end{equation}
\end{enumerate}

\SourcePage{168}

显式地，(218) 是如下两个函子的复合。其一是典范等价
\begin{equation}\tag{219}
\{\mathbb{G}_{m}\text{ 对 }Z^{\vee}\text{ 的平凡扩张}\}=Z\operatorname{tors},\qquad
\text{扩张}\mapsto\text{其分裂构成的 }Z\text{-torsor}
\end{equation}
其二是从 (219) 左端到 (218) 右端的函子，它由 (214) 诱导。换言之，(218) 就是 4.1.8 中的函子 $\mathcal{E}\mapsto\tilde{G}(K)_{\mathcal{E}}$。

(ii) 由 3.4.7 (iv)，函子 (216) 关于 Picard 函子 (218) 是仿射的。

\bdp{4.4.10} 设 $\mathcal{L}\in Z\operatorname{tors}_{\theta}(X)$。按 4.4.9，$\mathcal{L}$ 在 $Z\operatorname{tors}_{\theta}(O_{x})$ 中的像定义了 $G(K_{x})$ 的中心扩张，记之为 $\tilde{G}(K_{x})_{\mathcal{L}}$。记 $\lambda_{\mathcal{L},\underline{x}}$ 为 $\lambda_{\mathcal{L}}$ 到 $\operatorname{Bun}_{G,\underline{x}}$ 的拉回。$G(K_{x})$ 在 $\operatorname{Bun}_{G,\underline{x}}$ 上的作用提升为 $\tilde{G}(K_{x})_{\mathcal{L}}$ 在 $\lambda_{\mathcal{L},\underline{x}}$ 上的典范作用（用 4.3.7–4.3.9、4.1.8 以及 4.4.9 的注记）。$G(O_{x})\times\mathbb{G}_{m}\subset\tilde{G}(K_{x})_{\mathcal{L}}$ 以显然的方式作用在 $\lambda_{\mathcal{L},\underline{x}}$ 上。

\bdp{4.4.11} \emph{命题.} 对应于 (217) 的李代数扩张是 2.5.1 中的扩张
\[
0\to\mathbb{C}\to\widetilde{\mathfrak{g}}\otimes K\to\mathfrak{g}\otimes K\to 0.
\]

\emph{证明.} 对应于 (217) 的李代数扩张不依赖 $\mathcal{L}\in Z\operatorname{tors}_{\theta}(O)$，故可用 (212) 代替 (217)，最终用 (194)。于是只需用 Kac–Peterson–Frenkel 定理，它说李代数扩张
\begin{equation}\tag{220}
0\to\mathbb{C}\to\tilde{o}_{n}(K)\to o_{n}(K)\to 0
\end{equation}
对应于 (194)，由上闭链 $(u,v)\mapsto\frac{1}{2}\operatorname{Res}\operatorname{Tr}(du,v)$ 定义，$u,v\in o_{n}(K)$。事实上，要用 [KP] 或 [Fr81] 中的命题 I.3.11，须用到 $\tilde{o}_{n}(K)$ 的如下刻画（它不涉及群 $\tilde{O}_{n}(K)$）：设 $V$ 的含义同 4.3.3，$M$ 是 $Cl(V)$ 上的不可约离散模，则有

\SourcePage{169}

$\tilde{o}_{n}(K)$ 在 $M$ 中的表示，它与 $\tilde{o}_{n}(K)$ 在 $Cl(V)$ 上的作用相容，且使 $1\in\mathbb{C}\subset\tilde{o}_{n}(K)$ 恒等地作用在 $M$ 上。\hfill$\Box$

\bdp{4.4.12} 设 $\lambda_{\mathcal{L}}$ 与 $\lambda_{\mathcal{L},\underline{x}}$ 的含义同 4.4.10。按 4.4.10 与 4.4.11，$\mathfrak{g}\otimes K_{x}$ 在 $\operatorname{Bun}_{G,\underline{x}}$ 上的作用提升为 $\widetilde{\mathfrak{g}\otimes K_{x}}$ 在 $\lambda_{\mathcal{L},\underline{x}}$ 上的典范作用；它在 $\mathbb{C}\times(\mathfrak{g}\otimes O_{x})\subset\widetilde{\mathfrak{g}\otimes K_{x}}$ 上的限制是显然的那个。特别地，$1\in\mathbb{C}\subset\widetilde{\mathfrak{g}\otimes K_{x}}$ 作用为乘以 1。

$\lambda_{\mathcal{L}}$ 带有一个同构 $\lambda_{\mathcal{L}}^{\otimes 2n}\xrightarrow{\ \sim\ }(\omega^{\sharp}_{\operatorname{Bun}_{G}})^{\otimes n}$（对某个 $n\ne 0$），故作用在 $\lambda_{\mathcal{L}}$ 上的微分算子层是 $D'$。因此由 1.2.5，$\widetilde{\mathfrak{g}\otimes K_{x}}$ 在 $\lambda_{\mathcal{L},\underline{x}}$ 上的作用诱导一个典范态射
\[
h_{x}:\mathfrak{z}_{x}\to\Gamma(\operatorname{Bun}_{G},D').
\]
显然 $h_{x}$ 不依赖 $\mathcal{L}\in Z\operatorname{tors}_{\theta}(X)$。

\bdp{4.4.13} 本小节我们证明 4.4.12 中的 $h_{x}$ 与 2.5.4 中的 $h_{x}$ 一致。读者可跳过此证明，干脆忘掉 $h_{x}$ 的旧定义（当初引入它，只是为了在 §2 中回避 $\omega_{\operatorname{Bun}_{G}}$ 的平方根的讨论）。

要证 $h_{x}$ 的两个定义等价，只需说明：若 $\mathcal{L}$ 是 $\omega_{X}$ 的平方根，则同构 $\lambda_{\mathcal{L}}^{\otimes 2}\xrightarrow{\ \sim\ }\omega^{\sharp}_{\operatorname{Bun}_{G}}$ 在它们到 $\operatorname{Bun}_{G,\underline{x}}$ 的拉回之间诱导一个 $\mathfrak{g}\otimes K_{x}$-等变同构。这可以直接证明，但事实上不可能有别的情形。事实上，$\widetilde{\mathfrak{g}\otimes K_{x}}$-等变性的障碍是一个上闭链 $\widetilde{\mathfrak{g}\otimes K_{x}}\to H^{0}(\operatorname{Bun}_{G,\underline{x}},\mathcal{O})$。因 $\operatorname{Hom}(\mathfrak{g}\otimes K_{x},\mathbb{C})=0$，只需说明 $\operatorname{Bun}_{G,\underline{x}}$ 上每个正则函数 $f$ 局部常值。按 2.3.1，$\operatorname{Bun}_{G,\underline{x}}$ 是 $\operatorname{Bun}_{G,nx}$（$n\in\mathbb{N}$）的逆向极限。显然 $f$ 来自某个 $\operatorname{Bun}_{G,nx}$ 上的正则函数。故只需证明下述引理。

\emph{引理.} $\operatorname{Bun}_{G,nx}$ 上每个正则函数局部常值。

\emph{证明.} 取 $y\in X\setminus\{x\}$，考虑概形 $M$，它参数化 $X$ 上在 $nx$ 与 $y$ 的形式邻域上都平凡化的 $G$-丛（这里除子 $nx$ 看作子概形）。$G(K_{y})$ 作用在 $M$ 上，而正则

\SourcePage{170}

函数 $f$ 是 $H^{0}(M,\mathcal{O}_{M})$ 的 $G(O_{y})$-不变元。显然 $H^{0}(M,\mathcal{O}_{M})$ 是可积离散 $\mathfrak{g}\otimes K_{y}$-模。众所周知且极易证明：这样的模中 $(\mathfrak{g}\otimes O_{y})$-不变元必是 $(\mathfrak{g}\otimes K_{y})$-不变元。故 $f$ 是 $(\mathfrak{g}\otimes K_{y})$-不变的。因 $\mathfrak{g}\otimes K_{y}$ 在 $M$ 上的作用是（形式地）传递的，$f$ 局部常值。\hfill$\Box$

\emph{注.} 上述引理是众所周知的。标准的证法是把 $\operatorname{Bun}_{G,nx}$ 表示成 $\Gamma\backslash G(K_{y})/G(O_{y})$（对某个 $\Gamma\subset G(K_{y})$；$n=0$ 的情形见 [La-So]），然后利用 $G(K_{y})/G(O_{y})$ 上正则函数局部常值这一事实。

\bdp{4.4.14} 现在我们来完成水平性定理 2.7.3 的证明（证明的开头见 2.8.3–2.8.5）。

设 $M$ 是 $X$ 上的概形，它在 $x\in X$ 处的纤维是 $\operatorname{Bun}_{G,\underline{x}}$。固定 $\mathcal{L}\in\omega^{1/2}(X)$ 与 $\mathcal{L}^{\mathrm{loc}}\in\omega^{1/2}(O)$（即 $\mathcal{L}$ 是 $\omega_{X}$ 的平方根，$\mathcal{L}^{\mathrm{loc}}$ 是 $\omega_{O}$ 的平方根）。则有 4.3.16 中定义的概形 $X_{2}^{\wedge}$。记 $M_{2}^{\wedge}$ 为 $M$ 与 $X_{2}^{\wedge}$ 在 $X$ 上的纤维积。半直积 $\operatorname{Aut}_{2}O\ltimes G(K)$ 作用在 $M_{2}^{\wedge}$ 上（参看 4.3.16）。

取其中心扩张 $\operatorname{Aut}_{2}O\ltimes\widetilde{G(K)}$，其中 $\widetilde{G(K)}$ 是对应于 $\mathcal{L}^{\mathrm{loc}}$ 的中心扩张 (212)，而 $\operatorname{Aut}_{2}O=\operatorname{Aut}(O,\mathcal{L}^{\mathrm{loc}})$ 按结构的搬运作用在 $\widetilde{G(K)}=\widetilde{G(K)}_{\mathcal{L}^{\mathrm{loc}}}$ 上。记 $\lambda_{\mathcal{L}}^{\wedge}$ 为 Pfaffian 线丛 $\lambda'_{\mathcal{L}}$（来自 4.4.1）到 $M_{2}^{\wedge}$ 的拉回。因 $\operatorname{Aut}_{2}O$ 作用在 $M_{2}^{\wedge}$ 上，而 $M_{2}^{\wedge}$ 看作 $\operatorname{Bun}_{G}$ 上的概形，我们得到 $\operatorname{Aut}_{2}O$ 在 $\lambda_{\mathcal{L}}^{\wedge}$ 上的作用。另一方面，$\widetilde{G(K)}$ 作用在 $\lambda_{\mathcal{L},\hat x}^{\wedge}$ 上，后者定义为 $\lambda_{\mathcal{L}}^{\wedge}$ 在 $M_{2}^{\wedge}$ 于 $\hat x\in X_{2}^{\wedge}$ 处纤维上的限制。事实上，该纤维等于 $\operatorname{Bun}_{G,\underline{x}}$，其中 $x$ 是 $\hat x$ 在 $X$ 中的像；由 4.4.8，中心扩张 $\tilde{G}(K_{x})_{\mathcal{L}}$ 作用在 $\lambda'_{\mathcal{L},\underline{x}}=\lambda_{\mathcal{L},\hat x}^{\wedge}$ 上。该扩张只依赖 $\mathcal{L}_{x}$，即 $\mathcal{L}$ 到 $\operatorname{Spec}O_{x}$ 的拉回。因 $\hat x$ 给出同构 $(O_{x},\mathcal{L}_{x})\xrightarrow{\ \sim\ }(O,\mathcal{L}^{\mathrm{loc}})$，我们得到同构 $\tilde{G}(K_{x})_{\mathcal{L}}\xrightarrow{\ \sim\ }\widetilde{G(K)}$，从而得到 $\widetilde{G(K)}$ 在 $\lambda_{\mathcal{L},\hat x}^{\wedge}$ 上的作用。如 2.8.5 所述，为完成 2.7.3 的证明，只需说明

\SourcePage{171}

(i) 对应于各种 $\hat x\in X_{2}^{\wedge}$ 的 $\widetilde{G(K)}$ 在 $\lambda_{\mathcal{L},\hat x}^{\wedge}$ 上的作用，来自 $\widetilde{G(K)}$ 在 $\lambda_{\mathcal{L}}^{\wedge}$ 上的（显然唯一的）作用。

(ii) 该作用与 $\operatorname{Aut}_{2}O$ 的作用相容（即实际上我们有 $\operatorname{Aut}_{2}O\ltimes\widetilde{G(K)}$ 在 $\lambda_{\mathcal{L}}^{\wedge}$ 上的作用）。

这立即由 4.3.17 推出。

\bdp{4.4.15} 本小节与下一小节中，我们陈述并证明 4.4.14 的陈述 i) 与 ii) 的一个推广，它将在本书主要结果（定理 5.4.5）的证明中用到。推广是显然的（把 $\omega^{1/2}(X)$ 换成 $Z\operatorname{tors}_{\theta}(X)$，等等），读者目前完全可以跳过这两小节。

固定 $\mathcal{L}\in Z\operatorname{tors}_{\theta}(X)$ 与 $\mathcal{L}^{\mathrm{loc}}\in Z\operatorname{tors}_{\theta}(O)$。记 $X_{Z}^{\wedge}$ 为 $X^{\wedge}$ 的平展 $Z$-覆叠，使得 $X_{Z}^{\wedge}(R)$ 中对应于态射 $\alpha:\operatorname{Spec}(R\hat\otimes O)\to X$ 的点的原像，是群胚 $Z\operatorname{tors}_{\theta}(R\hat\otimes O)$ 中同构 $\mathcal{L}_{R}^{\mathrm{loc}}\xrightarrow{\ \sim\ }\alpha^{*}\mathcal{L}$ 的集合\footnote{这里为方便起见，采用 3.4.5 中 $Z\operatorname{tors}_{\theta}$ 的定义。}，其中 $\mathcal{L}_{R}^{\mathrm{loc}}$ 是 $\mathcal{L}^{\mathrm{loc}}$ 到 $\operatorname{Spec}R\hat\otimes O$ 的拉回。4.6.6 中的群归纳概形 $\operatorname{Aut}_{Z}O=\operatorname{Aut}(O,\mathcal{L}^{\mathrm{loc}})$ 按结构的搬运作用在 $X_{Z}^{\wedge}$ 上。记 $M_{Z}^{\wedge}$ 为 $M$ 与 $X_{Z}^{\wedge}$ 在 $X$ 上的纤维积。记 $\lambda_{\mathcal{L}}^{\wedge}$ 为 4.4.3 中定义的线丛 $\lambda_{\mathcal{L}}$ 到 $M_{Z}^{\wedge}$ 的拉回。半直积 $\operatorname{Aut}_{Z}O\ltimes G(K)$ 作用在 $M_{Z}^{\wedge}$ 上。取其中心扩张 $\operatorname{Aut}_{Z}O\ltimes\widetilde{G(K)}$，其中 $\widetilde{G(K)}$ 是对应于 $\mathcal{L}^{\mathrm{loc}}$ 的中心扩张 (217)，而 $\operatorname{Aut}_{Z}O=\operatorname{Aut}(O,\mathcal{L}^{\mathrm{loc}})$ 按结构的搬运作用在 $\widetilde{G(K)}=\widetilde{G(K)}_{\mathcal{L}^{\mathrm{loc}}}$ 上。我们把 $\operatorname{Aut}_{Z}O\ltimes G(K)$ 在 $M_{Z}^{\wedge}$ 上的作用提升为 $\operatorname{Aut}_{Z}O\ltimes\widetilde{G(K)}$ 在 $\lambda_{\mathcal{L}}^{\wedge}$ 上的作用。

正如 4.4.14，可定义 $\operatorname{Aut}_{Z}O$ 在 $\lambda_{\mathcal{L}}^{\wedge}$ 上的作用，以及 $\widetilde{G(K)}$ 在 $\lambda_{\mathcal{L},\hat x}^{\wedge}:=$ $\lambda_{\mathcal{L}}^{\wedge}$ 在 $M_{Z}^{\wedge}$ 于 $\hat x\in\hat X_{Z}$ 处纤维上的限制上的作用。
% [Tutor Note: 原书 p.171 同段前文印作 $X_{Z}^{\wedge}$，此处末句印作 $\hat X_{Z}$；底本如此，保留原样。]

\bdp{4.4.16} \emph{命题.}

(i) 对应于各种 $\hat x\in X_{Z}^{\wedge}$ 的 $\widetilde{G(K)}$ 在 $\lambda_{\mathcal{L},\hat x}^{\wedge}$ 上的作用，来自 $\widetilde{G(K)}$ 在 $\lambda_{\mathcal{L}}^{\wedge}$ 上的（显然唯一的）作用。
% ===================== end 源码/units/ch04_u05.tex =====================

% ===================== begin 源码/units/ch04_u06.tex =====================
% ch04_u06 —— 原书 p.172–p.181（4.5 全部：仿射 Grassmann 流形）
% 承接 p.171 末的命题陈述之 (ii)（本单元不写 \bdchap）。
\SourcePage{172}

(ii) $\operatorname{Aut}_Z O$ 与 $\widetilde{G}(K)$ 在 $\lambda_{\mathcal{L}}^{\wedge}$ 上的作用，定义了 $\operatorname{Aut}_Z O\ltimes\widetilde{G}(K)$ 的一个作用。

\emph{证明.} 把 $\mathcal{L}\in Z\operatorname{tors}_{\theta}(X)$ 写成 $\mathcal{L}=\mathcal{E}\cdot\mathcal{L}_0$，其中 $\mathcal{E}\in Z\operatorname{tors}(X)$，$\mathcal{L}_0\in\omega^{1/2}(X)$。按定义，$\lambda_{\mathcal{L}}=l_{\mathcal{E}}\otimes\lambda'_{\mathcal{L}_0}$（$\operatorname{Bun}_G$ 上 $\mu_\infty$-torsor $l_{\mathcal{E}}$ 的定义见 4.1.4 或 4.1.6）。

把 $\mathcal{L}^{\mathrm{loc}}$ 看作 $\omega^{1/2}(O)$ 的对象（这是可以的：$Z\operatorname{tors}_{\theta}(O)$ 与 $\omega^{1/2}(O)$ 都恰有一个同构类）。用 $\mathcal{L}_0$ 与 $\mathcal{L}^{\mathrm{loc}}$ 构造 $X^{\wedge}_2$、$M^{\wedge}_2$ 与 $\lambda^{\wedge}_{\mathcal{L}_0}$（见 4.4.14）。

把 $\mathcal{E}$ 看作 $Z$-覆叠 $\mathcal{E}\to X$。置 $X^{\wedge}_{\mathcal{E}}:=\mathcal{E}\times_X X^{\wedge}$，$M^{\wedge}_{\mathcal{E}}:=\mathcal{E}\times_X M^{\wedge}$，其中 $X^{\wedge}$ 与 $M^{\wedge}$ 的含义同 2.6.5 与 2.8.3。记 $l^{\wedge}_{\mathcal{E}}$ 为 $l_{\mathcal{E}}$ 到 $M^{\wedge}_{\mathcal{E}}$ 的拉回。

置 $M^{\wedge}_{\mathcal{E},2}:=\mathcal{E}\times_X M^{\wedge}_2$。有平展覆叠 $M^{\wedge}_{\mathcal{E},2}\to M^{\wedge}_2$、$M^{\wedge}_{\mathcal{E},2}\to M^{\wedge}_{\mathcal{E}}$，以及 $p:M^{\wedge}_{\mathcal{E},2}\to M^{\wedge}_Z$。显然 $p^*\lambda^{\wedge}_{\mathcal{L}}$ 是 $l^{\wedge}_{\mathcal{E}}$ 与 $\lambda^{\wedge}_{\mathcal{L}_0}$ 到 $M^{\wedge}_{\mathcal{E},2}$ 的拉回的张量积。下面分别考察 $l^{\wedge}_{\mathcal{E}}$ 与 $\lambda^{\wedge}_{\mathcal{L}_0}$。

半直积 $\operatorname{Aut} O\ltimes G(K)$ 作用在 $M^{\wedge}_{\mathcal{E}}$ 上；$\operatorname{Aut} O$ 在 $M^{\wedge}_{\mathcal{E}}$ 上的作用，典范地提升为它在 $l^{\wedge}_{\mathcal{E}}$ 上的作用（参看 4.4.14 或 2.8.5）。$G(K)$ 的作用对象是 $l^{\wedge}_{\mathcal{E}}$ 于 $X^{\wedge}_{\mathcal{E}}$ 各点纤维上的限制（见 4.1.7）。容易看出，这些作用来自 $\operatorname{Aut} O\ltimes G(K)$ 在 $l^{\wedge}_{\mathcal{E}}$ 上的一个作用。

另一方面，由 4.4.14，有 $\operatorname{Aut}_2 O\ltimes\widetilde{G}(K)$ 在 $\lambda^{\wedge}_{\mathcal{L}_0}$ 上的典范作用。

于是得到 $\operatorname{Aut}_2 O\ltimes\widetilde{G}(K)$ 在 $p^*\lambda^{\wedge}_{\mathcal{L}}$ 上的作用。它与 $\operatorname{Aut}_2 O$ 在 $\lambda^{\wedge}_{\mathcal{L}}$ 上的作用相容，也与 $\widetilde{G}(K)$ 在 $\lambda^{\wedge}_{\mathcal{L},\hat x}$（$\hat x\in X^{\wedge}_Z$）上的作用相容。

$p$ 是平展满射，故 $\operatorname{Aut}_2 O\ltimes\widetilde{G}(K)$ 在 $p^*\lambda^{\wedge}_{\mathcal{L}}$ 上的作用下降为它在 $\lambda^{\wedge}_{\mathcal{L}}$ 上的作用。$\operatorname{Aut}_Z O$ 由 $\operatorname{Aut}_2 O$ 与 $Z$ 生成，故只需证 $Z\subset\operatorname{Aut}_Z O$ 在 $\lambda^{\wedge}_{\mathcal{L}}$ 上的作用与 $\widetilde{G}(K)$ 的作用相容。这一点显然：对每个 $\hat x\in X^{\wedge}_Z$，$Z$ 与 $\widetilde{G}(K)$ 在 $\lambda^{\wedge}_{\mathcal{L},\hat x}$ 上的作用都相容。\hfill$\Box$

\bdsec{4.5. 仿射 Grassmann 流形}

仿射 Grassmann 流形 $\mathcal{GR}$ 是 fpqc-商 $G(K)/G(O)$，其中 $O=\mathbb{C}[[t]]$，$K=\mathbb{C}((t))$。\SourcePage{173}本节回顾 $\mathcal{GR}$ 的一些基本性质。4.6 节将构造并研究\emph{局部 Pfaffian 丛}；它是 $\mathcal{GR}$ 上的线丛。

仿射 Grassmann 流形在主要定理 5.2.6 的证明中起关键作用。不过读者暂时可以跳过本节。

在 4.5.1 – ? 中，$G$ 表示任一连通仿射代数群。连通性是无害的假设，因为 $G(K)/G(O)=G^0(K)/G^0(O)$，其中 $G^0$ 是 $G$ 的连通分支。
% [Tutor Note: 原书 p.173「In 4.5.1 – ?」中的交叉引用只印出一个问号、无编号，底本如此，保留原样。]

\bdp{4.5.1} \emph{定理.}
\begin{enumerate}[label=(\roman*),leftmargin=2.2em]
\item fpqc-商 $G(K)/G(O)$ 是归纳有限型的归纳概形。
\item $G(K)/G(O)$ 形式光滑。\footnote{形式光滑的定义见 7.11.1。}
\item 投影 $p:G(K)\to G(K)/G(O)$ 在 Zariski 拓扑下局部有截面。
\item $G(K)/G(O)$ 是归纳真的，当且仅当 $G$ 约化。
\item $G(K)$（等价地，$G(K)/G(O)$）既约，当且仅当 $\operatorname{Hom}(G,\mathbb{G}_m)=0$。
\end{enumerate}

\emph{注.} 这个定理是众所周知的。下面给出的证明，其实质部分只是引用 Faltings、Beauville、Laszlo 与 Sorger 的工作。

\emph{证明.} 对 $G=GL_n$，(i) 与 (iv) 成立。事实上，存在归纳真的归纳概形 $Gr(K^n)$，它参数化 $K^n$ 中的 c-格（细节见 7.11.2(iii)）。$GL_n(K)/GL_n(O)$ 等同于 $Gr(K^n)$ 中参数化 $O$-不变 c-格的闭子归纳概形。为对任意 $G$ 证明 (i) 与 (iv)，需要下面的引理。

\emph{引理.} 设 $G_1\subset G_2$ 是仿射代数群，且商 $U:=G_1\backslash G_2$ 拟仿射，即 $U$ 是仿射概形\SourcePage{174}$Z$ 的开子概形。设 fpqc-商 $G_2(K)/G_2(O)$ 是归纳有限型的归纳概形。则 $G_1(K)/G_1(O)$ 亦然，且态射
\begin{equation}\tag{221}
G_1(K)/G_1(O)\to G_2(K)/G_2(O)
\end{equation}
是局部闭嵌入。若 $U$ 仿射，则 (221) 是闭嵌入。

读者可用 4.5.2 中 $G(K)/G(O)$ 的整体解释轻易证明引理。我们更愿意给出局部证明。

\emph{证明.} 考虑态射 $f:G_1(K)\to Z(K)$。显然 $Z(O)$ 是 $Z(K)$ 的闭子概形，$U(O)$ 是 $Z(O)$ 的开子概形。故 $Y:=f^{-1}(U(O))$ 是 $G_2(K)$ 的局部闭子归纳概形；当 $U$ 仿射时它是闭的。显然 $Y\cdot G_2(O)=Y$，故 $Y$ 是某个局部闭子归纳概形 $Y'\subset G_2(K)/G_2(O)$ 的原像；当 $U$ 仿射时 $Y'$ 是闭的。由 $G_1(K)\subset Y$，有自然态射
\begin{equation}\tag{222}
G_1(K)\to Y'.
\end{equation}

我们断言 (222) 是 $G_1(O)$-torsor（$G_1(O)$ 经右平移作用在 $G_1(K)$ 上），从而 $G_1(K)/G_1(O)=Y'$。为看出 (222) 是 $G_1(O)$-torsor，注意态射 $Y\to Y'$ 是 $G_2(O)$-torsor，态射 $\varphi:Y\to U(O)=G_1(O)\backslash G_2(O)$ 是 $G_2(O)$-等变的，且 $G_1(K)=\varphi^{-1}(\bar e)$，其中 $\bar e\in G_1(O)\backslash G_2(O)$ 是 $e\in G_2(O)$ 的像。\hfill$\Box$

下面证明对任意 $G$ 的 (i) 与 (iv)。取嵌入 $G\hookrightarrow GL_n$。若 $G$ 约化，则 $GL_n/G$ 仿射，故由引理知 $G(K)/G(O)$ 是归纳真的归纳概形。对任意 $G$，我们将构造嵌入 $i:G\hookrightarrow G':=GL_n\times\mathbb{G}_m$，使 $G'/i(G)$ 拟仿射；由此推出 (i)。为构造 $i$，取 $GL_n$-模 $V$，使 $G\subset GL_n$ 是某个一维子空间 $l\subset V$ 的稳定化子。$G$ 在 $l$ 中的作用由某个 $\chi:G\to\mathbb{G}_m$ 给出。定义 $i:G\hookrightarrow G':=GL_n\times\mathbb{G}_m$，$i(g)=(g,\chi(g)^{-1})$。为证 $G'/i(G)$ 拟仿射，把 $V$ 看作 $G'$-模（$\lambda\in\mathbb{G}_m$ 作用为乘 $\lambda$），并注意非零 $v\in l$ 在 $G'$ 中的稳定化子等于 $i(G)$。故 $G'/i(G)\simeq G'v$，而 $G'v$ 拟仿射。

\SourcePage{175}

下面完成 (iv) 的证明。若 $G(K)/G(O)$ 是归纳真的，且 $G'$ 是 $G$ 的正规子群，则按引理 $G'(K)/G'(O)$ 也是归纳真的。显然 $\mathbb{G}_a(K)/\mathbb{G}_a(O)$ 不是归纳真的。因此只有当 $G$ 约化时，$G(K)/G(O)$ 才归纳真。

为证 (iii)，只需证 $p:G(K)\to G(K)/G(O)$ 在任一 $\mathbb{C}$-点 $x\in G(K)/G(O)$ 的邻域上有截面（这里用到由 (i) 知 $\mathbb{C}$-点在 $G(K)/G(O)$ 中稠密）。由 $p$ 的 $G(K)$-等变性，可化到 $x$ 是 $e\in G(K)$ 的像的情形。于是须构造含 $e$ 的子归纳概形 $\Gamma\subset G(K)$，使态射
\begin{equation}\tag{223}
\Gamma\times G(O)\to G(K),\qquad (\gamma,g)\mapsto\gamma g
\end{equation}
是开浸入。按 Faltings [Fal94, p.350–351]，当 $\Gamma$ 的 $R$-点集由
\[
\Gamma(R)=\operatorname{Ker}\bigl(G(R[t^{-1}])\xrightarrow{f}G(R)\bigr)\subset G(R((t)))=G(R\hat\otimes K)
\]
定义时，(223) 是开浸入，其中 $f$ 是在 $t=\infty$ 处的求值。这个断言的证明属于 Beauville 与 Laszlo（[BLa94] 中的命题 1.11）。它以 $X=\mathbb{P}^1$ 情形下 $G(K)/G(O)$ 的整体解释为基础（见 4.5.2），也以 $\mathbb{P}^1$ 上 $G$-丛的下述性质为基础：对 $S\times\mathbb{P}^1$ 上的 $G$-丛 $\mathcal{F}$，使 $\mathcal{F}$ 在 $s\times\mathbb{P}^1$ 上的限制平凡的那些点 $s\in S$ 构成 $S$ 的\emph{开}子集（事实上，$H^1(\mathbb{P}^1,\mathcal{O}\otimes\mathfrak{g})=0$，$\mathfrak{g}:=\operatorname{Lie} G$）。

由 (iii) 推出\footnote{事实上，不用 (iii) 也能证明 (ii)。} (ii)。$G(K)$ 形式光滑，故由 (iii) 知 $G(K)/G(O)$ 的每一点都有形式光滑的邻域。又 $G(K)/G(O)$ 是归纳有限型的，故得 (ii)。

还需考虑 (v)。$G(O)$ 是既约的。故 $G(K)$ 既约当且仅当 $G(K)/G(O)$ 既约。Laszlo 与 Sorger 证明：若 $\operatorname{Hom}(G,\mathbb{G}_m)=0$，则 $G(K)/G(O)$ 既约（见 [La-So] 中命题 4.6 的证明）；其证明基于\SourcePage{176}Shafarevich 的一个定理。若 $\operatorname{Hom}(G,\mathbb{G}_m)\ne 0$，则存在态射 $f:\mathbb{G}_m\to G$ 与 $\chi:G\to\mathbb{G}_m$，使 $\chi f=\varphi_n$，$n\ne 0$，其中 $\varphi_n(\lambda):=\lambda^n$。由 $\varphi_n$ 诱导的态射 $\mathbb{G}_m(K)\to\mathbb{G}_m(K)$ 的像不含于 $\mathbb{G}_m(K)_{red}$ 中，故 $G(K)$ 不是既约的。\hfill$\Box$

\bdp{4.5.2} 设 $X$ 是 $\mathbb{C}$ 上的连通光滑射影曲线，$x\in X(\mathbb{C})$，$O_x$ 是 $x$ 处的完备化局部环，$K_x$ 是其分式域。按 Beauville–Laszlo（见 2.3.4），fpqc-商 $G(K_x)/G(O_x)$ 可解释为对 $(\mathcal{F},\gamma)$ 的模空间。这里 $\mathcal{F}$ 是 $X$ 上的主 $G$-丛，$\gamma:X\setminus\{x\}\to\mathcal{F}$ 是它的截面（即平凡化）。把 $(\mathcal{F},\gamma)$ 对应到 $\gamma/\gamma_x$ 在 $G(K_x)/G(O_x)$ 中的像；其中 $\gamma_x$ 是 $\mathcal{F}$ 在 $\operatorname{Spec} O_x$ 上的截面，$\gamma/\gamma_x$ 表示满足 $\gamma=g\gamma_x$ 的元素 $g\in G(K_x)$。我们已在 $\mathbb{C}$-点层面把 $G(K_x)/G(O_x)$ 与对 $(\mathcal{F},\gamma)$ 的模空间等同；对任意 $\mathbb{C}$-代数 $R$，读者容易对 $R$-点完成同样的事。

\bdp{4.5.3} 下面回顾 $G$ 的拓扑基本群的代数定义。记 $\pi_1^{\text{ét}}(G)$ 为 $G$ 在 Grothendieck 意义下的基本群。特征 $f:G\to\mathbb{G}_m$ 诱导态射 $\pi_1^{\text{ét}}(G)\to\pi_1^{\text{ét}}(\mathbb{G}_m)=\hat{\mathbb{Z}}(1)$，从而诱导态射 $f_*:(\pi_1^{\text{ét}}(G))(-1)\to\hat{\mathbb{Z}}$。记 $\pi_1(G)$ 为 $(\pi_1^{\text{ét}}(G))(-1)$ 中满足下述条件的 $\alpha$ 的集合：对一切 $f\in\operatorname{Hom}(G,\mathbb{G}_m)$ 有 $f_*(\alpha)\in\hat{\mathbb{Z}}$。我们把 $\pi_1(G)$ 看作离散群。事实上，把 $G$ 换成它的极大约化商，$\pi_1(G)$ 不变。对约化 $G$，可把 $\pi_1(G)$ 等同于 $G$ 的余权群模去余根格所得的商。

对任意有限覆叠 $p:\tilde{G}\to G$，有上边缘映射 $G(K)\to H^1(K,A)=A(-1)$，$A:=\operatorname{Ker} p$。这些映射给出同态 $G(K)\to(\pi_1^{\text{ét}}(G))(-1)$。其像含于 $\pi_1(G)$ 中。于是我们构造了典范同态
\begin{equation}\tag{224}
\varphi:G(K)\to\pi_1(G),
\end{equation}
\SourcePage{177}其中 $G(K)$ 按朴素意义理解（即看作 $G$ 的 $K$-点群，或归纳概形 $G(K)$ 的 $\mathbb{C}$-点群）。(224) 在 $G(O)$ 上的限制平凡，故 (224) 诱导映射
\begin{equation}\tag{225}
G(K)/G(O)\to\pi_1(G),
\end{equation}
其中 $G(K)/G(O)$ 也按朴素意义理解。

现在把 $G(K)$ 与 $G(K)/G(O)$ 看作归纳概形。$G(K)/G(O)$ 的 $\mathbb{C}$-点集在 $G(K)/G(O)$ 中稠密；$G(K)$ 亦然。

\bdp{4.5.4} \emph{命题.}
\begin{enumerate}[label=(\roman*),leftmargin=2.2em]
\item 映射 (224) 与 (225) 局部常值。
\item 相应的映射
\begin{equation}\tag{226}
\pi_0(G(K))\to\pi_1(G),
\end{equation}
\begin{equation}\tag{227}
\pi_0(G(K)/G(O))\to\pi_1(G)
\end{equation}
是双射。
\end{enumerate}

\emph{证明.} 我们已用整体论证证明了 (i)（见 4.1.7 末的注）。同样的论证可改用 4.5.2 中 $G(K_x)/G(O_x)$ 的解释来表述：映射 (225) 等于自然映射 $G(K_x)/G(O_x)\to\operatorname{Bun}_G$ 与「第一陈类」映射 $c:\pi_0(\operatorname{Bun}_G)\to\pi_1(G)$ 的复合取负。(i) 的局部证明见 4.5.5。

现在证 (ii)。映射 $\pi_0(G(K))\to\pi_0(G(K)/G(O))$ 是双射（因为 $G$ 连通）。故只需考虑 (226)。因 $G$ 可表为约化群与幂幺群的半直积，可设 $G$ 约化。固定嘉当子群 $H\subset G$。有 $\pi_0(H(K))=\pi_1(H)$；复合 $\pi_0(H(K))\to\pi_0(G(K))\to\pi_1(G)$ 就是自然映射 $\pi_1(H)\to\pi_1(G)$，它是满射。故 (226) 也是满射。映射 $\pi_0(H(K))\to\pi_0(G(K))$ 是满射（对抽象群 $G(K)$ 用 Bruhat 分解）。因此，为证 (226) 单射，只需证自然态射\SourcePage{178}$f:\pi_0(H(K))\to\pi_1(G)$ 的核含于 $\pi_0(H(K))\to\pi_0(G(K))$ 的核中。因 $\operatorname{Ker} f$ 是余根格，只需证：对任意余根 $\gamma:\mathbb{G}_m\to H$，$\mathbb{G}_m(K)$ 在 $G(K)$ 中的像属于 $e\in G(K)$ 的连通分支。余根 $\mathbb{G}_m\to H$ 可延拓为态射 $SL(2)\to G$，故只需注意 $SL(2,K)$ 连通（因为 $SL(2,K)$ 中任一矩阵都可表为幂幺矩阵的乘积）。\hfill$\Box$

下一小节给出 4.5.4(i) 的局部证明。

\bdp{4.5.5} \emph{引理.} 设 $M=\operatorname{Spec} R$ 是连通仿射簇，$A$ 是有限交换群，$\alpha\in H^1_{\text{ét}}(\operatorname{Spec} R((t)),A)$。对 $x\in M(\mathbb{C})$，记 $\alpha(x)$ 为 $\alpha$ 在 $x$ 上方纤维 $\operatorname{Spec} R((t))\to\operatorname{Spec} R$ 上的限制，于是 $\alpha(x)\in H^1_{\text{ét}}(\operatorname{Spec}\mathbb{C}((t)),A)=A(-1)$。则 $\alpha(x)\in A(-1)$ 不依赖于 $x$。

\emph{证明.} 只需证：对任意光滑连通 $M'$ 与任意态射 $M'\to M$，$\alpha$ 到 $M'(\mathbb{C})$ 的拉回是常值\footnote{事实上，只须考虑那些光滑曲线 $M'$。}。于是可设 $M$ 光滑。置 $V:=\operatorname{Spec} R[[t]]$，$V':=\operatorname{Spec} R((t))$。可设 $A=\mu_n$。则 $\alpha$ 对应 $V'$ 上的 $\mu_n$-torsor，即 $V'$ 上的线丛 $\mathcal{A}$ 连同一个同构 $\psi:\mathcal{A}^{\otimes n}\xrightarrow{\ \sim\ }\mathcal{O}_{V'}$。因 $V$ 正则，$\mathcal{A}$ 延拓为 $V$ 上的线丛 $\tilde{\mathcal{A}}$。此时 $\psi$ 诱导同构 $\tilde{\mathcal{A}}^{\otimes n}\xrightarrow{\ \sim\ }t^k\mathcal{O}_V$，$k\in\mathbb{Z}$。显然 $\alpha(x)\in\mathbb{Z}/n\mathbb{Z}$ 就是 $k$ 的像。\hfill$\Box$

下面给出 4.5.4(i) 的局部证明。因 $G(K)/G(O)$ 是归纳有限型的，只需证：对每个连通仿射簇 $M=\operatorname{Spec} R$ 与任意态射 $f:M\to G(K)$，复合 $M(\mathbb{C})\to G(K)\to\pi_1(G)$ 是常值。对任意有限交换群 $A$，正合列 $0\to A\to\tilde{G}\to G\to 0$ 定义了映射 $\pi_1(G)\to A(-1)$；故只需证复合 $M(\mathbb{C})\to G(K)\to\pi_1(G)\to A(-1)$ 常值。为此把引理用于 $\alpha=\varphi^*\beta$，其中 $\varphi:\operatorname{Spec} R((t))\to G$ 对应于\SourcePage{179} $f:\operatorname{Spec} R\to G(K)$，而 $\beta\in H^1_{\text{ét}}(G,A)$ 是 $\tilde{G}$ 作为 $G$ 上 $A$-torsor 的类。

\emph{注.} 事实上可以证明：对 $\mathbb{C}$ 上任一仿射概形 $M=\operatorname{Spec} R$，所谓「Künneth 态射」
\begin{equation}\tag{228}
\begin{gathered}
H^1_{\text{ét}}(M,A)\oplus H^0(M,\mathbb{Z})\otimes H^1_{\text{ét}}(\operatorname{Spec}\mathbb{C}((t)),A)\to H^1_{\text{ét}}(M((t)),A),\\
M((t)):=\operatorname{Spec} R((t)),
\end{gathered}
\end{equation}
是同构（显然这蕴含引理）。对使 $A$ 的阶在 $R$ 中可逆的任意环 $R$，类似的陈述也成立。

\bdp{4.5.6} \emph{命题.} 设 $A\subset G$ 是有限中心子群，$G':=G/A$。
\begin{enumerate}[label=(\roman*),leftmargin=2.2em]
\item 态射 $G(K)/G(O)\to G'(K)/G'(O)$ 诱导 $G(K)/G(O)$ 与 $G'(K)/G'(O)$ 的若干连通分支之并之间的同构。
\item 态射 $G(K)\to G'(K)$ 是平展覆叠。
\end{enumerate}

\emph{注.} 由 4.5.4，(i) 中提到的分支由 $\operatorname{Im}(\pi_1(G)\to\pi_1(G'))$ 的元素标号。$G(K)$ 在 $G'(K)$ 中的像的连通分支也如此。

\emph{证明.} 显然 (i) 与 (ii) 等价。设 $G$ 半单（这等价于 $G'$ 半单）来证 (i)。此时态射 $f:G(K)/G(O)\to G'(K)/G'(O)$ 由 4.5.1(iv) 是归纳真的。由 4.5.4(i)，对满足 $f^{-1}(C)\ne\emptyset$ 的分支 $C\subset G'(K)/G'(O)$，$f$ 在其几何点\footnote{对 $\mathbb{C}$-点的陈述由 4.5.4(i) 立即得到。把 $\mathbb{C}$ 换成代数闭域 $E\supset\mathbb{C}$ 后 4.5.4 仍成立，故对 $E$-点的陈述也成立。}上的纤维恰含一点，且容易看出这些纤维既约。由 4.5.1(v)，$G'(K)/G'(O)$ 既约。故半单情形下 (i) 显然。

\SourcePage{180}

下面把 (ii) 的证明化到半单情形。可设 $A$ 循环。只需构造从 $G$ 到半单群 $G_1$ 的态射 $\rho$，使 $\rho|_A$ 单射且 $\rho(A)\subset G_1$ 是中心的（此时态射 $G(K)\to G'(K)$ 由 $G_1(K)\to G'_1(K)$ 经基变换得到，其中 $G'_1:=G_1/\rho(A)$）。构造 $G_1$ 与 $\rho$ 的办法如下。固定同构 $\chi:A\xrightarrow{\ \sim\ }\mu_n$。设 $V$ 是有限维 $G$-模，使 $Z$ 经 $\chi$ 作用在 $V$ 上。记 $W_{pq}$ 为 $p$ 份 $V$ 与 $q$ 份 $\operatorname{Sym}^{n-1}V^*$ 的直和。若 $p\cdot\dim V=q(n-1)\cdot\dim\operatorname{Sym}^{n-1}V$，则可置 $G_1:=SL(W_{pq})$（事实上，$GL(V)$ 在 $GL(W_{pq})$ 中的像含于 $SL(W_{pq})$。\hfill$\Box$
% [Tutor Note: 原书 p.180「(indeed, the image of … is contained in $SL(W_{pq})$.)」的左括号未闭合，底本如此，保留原样。]

\emph{注记.}
\begin{enumerate}[label=(\roman*),leftmargin=2.2em]
\item 命题 4.5.6 是 (228) 之双射性的直接推论。
\item 用 4.5.2 中 $G(K)/G(O)$ 的整体解释，容易证明命题 4.5.6。
\end{enumerate}

\bdp{4.5.7} 设 $G$ 约化。记 $G_{ad}$ 为 $G$ 模去其中心的商。置 $T:=G/[G,G]$，$G':=G_{ad}\times T$。则 $G'=G/A$，其中 $A\subset G$ 是某有限中心子群。故由 4.5.6，$G(K)/G(O)$ 可等同于 $G'(K)/G'(O)=G_{ad}(K)/G_{ad}(O)\times T(K)/T(O)$ 的某些连通分支之并。

$T(K)/T(O)$ 的结构相当简单。例如 $\mathbb{G}_m(K)/\mathbb{G}_m(O)$ 的既约部分是离散空间 $\mathbb{Z}$；$1\in\mathbb{G}_m(K)/\mathbb{G}_m(O)$ 的连通分支是以 $K/O$ 为李代数的形式群。

\bdp{4.5.8} 从现在起设 $G$ 约化，并置 $\mathcal{GR}:=G(K)/G(O)$。

回忆 $\mathcal{GR}$ 中的 $G(O)$-轨道由 $G$ 的支配余权标号；换言之，由 $P_+({}^LG):=$ ${}^LG$ 的支配权集标号。更确切地说，$\chi\in P_+({}^LG)$ 定义了态射 $\nu:\mathbb{G}_m\to G$ 的一个共轭类；按定义，$\operatorname{Orb}_\chi$ 是 $\nu(\pi)$ 在 $\mathcal{GR}$ 中的像的 $G(O)$-轨道，其中\SourcePage{181}$\pi$ 是 $O$ 的一个素元（这个像不依赖于 $\pi$ 的选取）。显然 $\operatorname{Orb}_\chi$ 不依赖于共轭类中 $\nu$ 的选取，故 $\operatorname{Orb}_\chi$ 有定义。按 [IM]，映射 $\chi\mapsto\operatorname{Orb}_\chi$ 是 $P_+({}^LG)$ 与 $\mathcal{GR}$ 中 $G(O)$-轨道集之间的双射。容易证明
\begin{equation}\tag{229}
\dim\operatorname{Orb}_\chi=(\chi,2\rho),
\end{equation}
其中 $2\rho$ 是 $G$ 的正根之和。

\emph{注.} 显然 $\operatorname{Orb}_\chi$ 是 $\operatorname{Aut}^0 O$-不变的。

\bdp{4.5.9} 已有 $\pi_0(\mathcal{GR})$ 与 $\pi_1(G)$ 之间的双射 (227)。设 $Z$ 是朗兰兹对偶群 ${}^LG$ 的中心。利用 $G$ 与 ${}^LG$ 的嘉当环面之间的对偶，把 $\pi_1(G)$ 等同于 $Z^\vee:=\operatorname{Hom}(Z,\mathbb{G}_m)$。于是 $\mathcal{GR}$ 的连通分支由 $Z^\vee$ 的元素标号。

\emph{注.} 含 $\operatorname{Orb}_\chi$ 的 $\mathcal{GR}$ 连通分支对应于 $\chi_Z\in Z^\vee$，其中 $\chi_Z$ 是 $\chi\in P^+({}^LG)$ 到 $Z$ 的限制。
% [Tutor Note: 原书 p.181 此处记作 $P^+({}^LG)$，而 p.180 同一对象记作 $P_+({}^LG)$，底本记号不一致，保留原样。]

\bdp{4.5.10} 存在典范态射 $\alpha:\mu_2\to Z$。当 $G$ 半单时，我们已用 (56) 定义过它。当 $G$ 约化时，这给出态射 $\mu_2\to Z'$，其中 $Z'$ 是 ${}^LG$ 的中心化子的中心；于是定义 $\alpha$ 为复合 $\mu_2\to Z'\hookrightarrow Z$。

按 4.4.4，对偶态射 $\alpha^\vee:\pi_1(G)\to\mathbb{Z}/2\mathbb{Z}$ 就是由伴随表示 $G\to SO(\mathfrak{g}_{ss})$ 给出的基本群之间的态射，其中 $\mathfrak{g}_{ss}:=[\mathfrak{g},\mathfrak{g}]$。

(227) 与 $\alpha^\vee$ 的复合定义了局部常值的\emph{奇偶函数}
\begin{equation}\tag{230}
p:\mathcal{GR}\to\mathbb{Z}/2\mathbb{Z}.
\end{equation}

若 (230) 把 $\mathcal{GR}$ 的连通分支映到 0（相应地 1），则称它是\emph{偶}的（相应地\emph{奇}的）。

\bdp{4.5.11} \emph{命题.} $\mathcal{GR}$ 的偶（相应地奇）连通分支上的一切 $G(O)$-轨道，维数都是偶（相应地奇）的。
% ===================== end 源码/units/ch04_u06.tex =====================

% ===================== begin 源码/units/ch04_u07.tex =====================
% ch04_u07 —— 原书 p.182–p.191（4.5.11 的证明、4.5.12、4.6.1–4.6.20）
% 本单元自 p.182 起，开篇为命题 4.5.11 的证明（该命题的陈述在 p.181，属上一单元）。
\SourcePage{182}

\emph{证明.} 设 $x=gG(O)\in\mathcal{GR}$。利用 $\alpha^{\vee}$ 与伴随表示之间的关系（见 4.5.10），以及 4.3.4 的注 (ii) 与 (iii)，可见 $x$ 属于 $\mathcal{GR}$ 的偶分支，当且仅当
\begin{equation}\tag{231}
\dim\mathfrak{g}_{ss}\otimes O/((\mathfrak{g}_{ss}\otimes O)\cap\operatorname{Ad}_g(\mathfrak{g}_{ss}\otimes O))
\end{equation}
为偶数。而（231）正是 $x$ 的 $G(O)$-轨道的维数。\hfill$\Box$

另有如下证明。利用（229）与 4.5.9 的注，可见该命题等价于公式 $\chi_Z(\alpha(-1))=(-1)^{\langle\chi,2\rho\rangle}$。此式显然：按（56），$\alpha:\mu_2\to Z$ 是相应于 $2\rho$ 的态射 $\lambda^{\#}:\mathbb{G}_m\to H\subset G$ 的限制。

\bdp{4.5.12} 下文关于 $\mathcal{GR}$ 中 $G(O)$-轨道的性质，本文不会用到，但我们认为仍值得一提。

$\overline{\operatorname{Orb}}_\chi$ 是诸 $\operatorname{Orb}_{\chi'}$（$\chi'\le\chi$）的并。事实上，设 $\rho:G\to GL(V)$ 是以 $\lambda$ 为最低权的表示，则对 $g\in\operatorname{Orb}_\chi$ 有 $\rho(g)\in t^{(\chi,\lambda)}\operatorname{End}(V\otimes O)$，而 $\rho(g)\notin t^{(\chi,\lambda)+1}\operatorname{End}(V\otimes O)$。于是若 $\operatorname{Orb}_{\chi'}\subset\overline{\operatorname{Orb}}_\chi$，则对 $G$ 的每个反支配权 $\lambda$ 都有 $(\chi-\chi',\lambda)\le0$。因此 $\chi-\chi'$ 是 $G$ 的单余根的线性组合，系数非负；由 4.5.4(i)，这些系数是整数，故 $\chi'\le\chi$。另一方面，一次 $GL(2)$ 计算表明：${}^{L}G$ 中满足 $\operatorname{Orb}_{\chi'}\subset\overline{\operatorname{Orb}}_\chi$ 的权 $\chi'$ 构成的集合，在 [Bour75] 第 VIII 章 §7 第 2 号的意义下是饱和的。于是 loc.cit 的命题 5 表明：对每个满足 $\chi'\le\chi$ 的支配权 $\chi'$，都有 $\operatorname{Orb}_{\chi'}\subset\overline{\operatorname{Orb}}_\chi$。

上述对 $\overline{\operatorname{Orb}}_\chi$ 的描述表明：$\operatorname{Orb}_\chi$ 闭，当且仅当 $\chi$ 极小。若 $G$ 单，则 $\chi$ 极小当且仅当 $\chi=0$，或 $\chi$ 是 ${}^{L}G$ 的微权（见 [Bour68] 第 VI 章 §2 习题 5）。于是 $\mathcal{GR}$ 的每个连通分支上恰有一个闭的 $G(O)$-轨道（用 4.5.4 与 loc.cit 习题的第一部分）。若 $\operatorname{Orb}_\chi$ 闭，则它是射影的；此时 $G(O)$ 经 $G=G(O/tO)$ 作用在 $\operatorname{Orb}_\chi$ 上，而 $\operatorname{Orb}_\chi$ 是 $G$ 对某个抛物子群的商。用 9.1.3 的话说，$\operatorname{Orb}_\chi=\operatorname{orb}_\chi=G/P_\chi^-$。

\SourcePage{183}

设 $G$ 单。则恰有一个 $\chi$，使得 $\overline{\operatorname{Orb}}_\chi\setminus\operatorname{Orb}_\chi$ 由单个点构成\footnote{当然，该点是 $e\in G(K)$ 的像。}；这个 $\chi$ 是 $\mathfrak{g}:=\operatorname{Lie}G$ 中相应于 $\mathfrak{g}$ 的最大根 $\alpha_{max}$ 的余根（见 [Bour75] 第 VIII 章 §7 习题 22）。此时 $\overline{\operatorname{Orb}}_\chi$ 可描述如下。置 $V:=\mathfrak{g}\otimes(\mathfrak{m}^{-1}/O)$，其中 $\mathfrak{m}$ 是 $O$ 的极大理想。记 $\bar V$ 为以 $V$ 为仿射子空间的射影空间。于是 $\bar V$ 是 $V\oplus\mathbb{C}$ 中直线构成的空间；特别地，$V^*=\mathfrak{g}^*\otimes(\mathfrak{m}/\mathfrak{m}^2)$ 作用在 $\bar V$ 上，并保持 $0\in V$。记 $C$ 为 $V$ 中那些 $G$-共轭于 $\mathfrak{g}_{\alpha_{max}}\otimes(\mathfrak{m}^{-1}/O)$ 的元素构成的集合。这是 $V$ 的闭子簇。其射影闭包 $\bar C\subset\bar V$ 是 $V^*$-不变的，因为 $C$ 是锥。容易证明，态射 $\exp:C\to G(K)/G(O)$ 可延拓为同构 $f:\bar C\xrightarrow{\ \sim\ }\overline{\operatorname{Orb}}_\chi$。显然 $f$ 既是 $\operatorname{Aut}^0 O$-等变的，也是 $G$-等变的。$\operatorname{Ker}(G(O)\to G(O/\mathfrak{m}))$ 在 $\bar C$ 上的作用由它在 $\overline{\operatorname{Orb}}_\chi$ 上的作用诱导，该作用来自 $V^*$ 在 $\bar C$ 上的作用，以及同构
\[
\operatorname{Ker}(G(O/\mathfrak{m}^2)\to G(O/\mathfrak{m}))\xrightarrow{\ \sim\ }\mathfrak{g}\otimes\mathfrak{m}/\mathfrak{m}^2\xrightarrow{\ \sim\ }V^*,
\]
其中最后一个箭头由 $\mathfrak{g}$ 上使 $(\alpha_{max},\alpha_{max})=2$ 的不变内积诱导。

\bdsec{4.6. 局部 Pfaffian 丛}

考虑仿射 Grassmann 流形 $\mathcal{GR}:=G(K)/G(O)$，其中 $O=\mathbb{C}[[t]]$，$K=\mathbb{C}((t))$。置 $Z:=\operatorname{Hom}(\pi_1(G),\mathbb{G}_m)$（由 4.1.1 的注，$Z$ 是 ${}^{L}G$ 的中心）。本小节将构造并研究一个函子 $\mathcal{L}\mapsto\lambda_{\mathcal{L}}=\lambda_{\mathcal{L}}^{\mathrm{loc}}$，它从群胚 $Z\operatorname{tors}_\theta(O)$（见 3.4.3）指向 $\mathcal{GR}$ 上线丛的范畴。称 $\lambda_{\mathcal{L}}$ 为相应于 $\mathcal{L}$ 的\emph{局部 Pfaffian 丛}。

我们建议读者眼下跳过本小节。

\bdp{4.6.1} 在 4.4.9 中我们定义了函子 $\mathcal{L}\mapsto\widetilde{G(K)}_{\mathcal{L}}$，它从 $Z\operatorname{tors}_\theta(O)$ 指向 $G(K)$ 被 $\mathbb{G}_m$ 的中心扩张的范畴。对 $\mathcal{L}\in Z\operatorname{tors}_\theta(O)$ 有分裂 $G(O)\to\widetilde{G(K)}_{\mathcal{L}}$，因而有主 $\mathbb{G}_m$-丛
\begin{equation}\tag{232}
\widetilde{G(K)}_{\mathcal{L}}/G(O)\to G(K)/G(O)=\mathcal{GR}.
\end{equation}

\SourcePage{184}

\bdp{4.6.2} \emph{定义.} $\lambda_{\mathcal{L}}$ 是 $\mathcal{GR}$ 上相应于 $\mathbb{G}_m$-丛（232）的线丛之逆。

显然 $\lambda_{\mathcal{L}}$ 函子性地依赖于 $\mathcal{L}\in Z\operatorname{tors}_\theta(O)$。

\bdp{4.6.3} \emph{注.} $\widetilde{G(K)}_{\mathcal{L}}$ 依赖于 $\mathfrak{g}$ 上非退化不变双线性型的选取（见 4.4.7）。故 $\lambda_{\mathcal{L}}$ 亦然。

\bdp{4.6.4} 记 $\bar e\in\mathcal{GR}$ 为单位元 $e\in G$ 的像。我们的 $\lambda_{\mathcal{L}}$ 是 $\mathcal{GR}$ 上唯一的 $\widetilde{G(K)}_{\mathcal{L}}$-等变线丛：它在 $\bar e$ 上平凡化，且任意 $c\in\mathbb{G}_m\subset\widetilde{G(K)}_{\mathcal{L}}$ 通过乘以 $c^{-1}$ 作用在 $\lambda_{\mathcal{L}}$ 上。唯一性由等式 $\operatorname{Hom}(G(O),\mathbb{G}_m)=0$ 推出。

\bdp{4.6.5} 由 4.4.11，$\widetilde{G(K)}_{\mathcal{L}}$ 在 $\lambda_{\mathcal{L}}$ 上的作用诱导出 $\widetilde{\mathfrak{g}}\otimes K$ 在 $\lambda_{\mathcal{L}}$ 上的作用，使得 $1\in\mathbb{C}\subset\widetilde{\mathfrak{g}}\otimes K$ 通过乘以 $-1$ 作用。它与 $\mathfrak{g}\otimes K$ 经左无穷小平移在 $\mathcal{GR}$ 上的作用相容。

\bdp{4.6.6} （63）经态射（56）的前推是正合列
\begin{equation}\tag{233}
0\to Z\to\operatorname{Aut}_Z O\to\operatorname{Aut}O\to0.
\end{equation}
对任意 $\mathcal{L}\in Z\operatorname{tors}(O)$，正合列
\begin{equation}\tag{234}
0\to Z\to\operatorname{Aut}(O,\mathcal{L})\to\operatorname{Aut}O\to0
\end{equation}
可与（233）典范地等同。这里 $\operatorname{Aut}(O,\mathcal{L})$ 是由对 $(\sigma,\varphi)$ 构成的群归纳概形，其中 $\sigma\in\operatorname{Aut}O$，$\varphi:\mathcal{L}\xrightarrow{\ \sim\ }\sigma_*\mathcal{L}$（读者或许更愿意把 $\mathcal{L}$ 视为 3.4.5 中范畴 $\widetilde{Z}\operatorname{tors}_\omega(O)$ 的对象）。（233）与（234）之间的同构由显然的态射 $\operatorname{Aut}_2 O:=\operatorname{Aut}(O,\omega_O^{1/2})\to\operatorname{Aut}(O,\mathcal{L})$ 诱导。

$\operatorname{Aut}_Z O=\operatorname{Aut}(O,\mathcal{L})$ 通过结构的转移作用在正合列（217）上；$\operatorname{Aut}_Z O$ 在 $\mathbb{G}_m$ 上的作用平凡，其在 $G(K)$ 上的作用来自 $\operatorname{Aut}O$ 在 $G(K)$ 上的通常作用。子群 $G(O)\subset\widetilde{G(K)}_{\mathcal{L}}$ 是 $\operatorname{Aut}_Z O$-不变的。

\SourcePage{185}

\bdp{4.6.7} 由 4.6.6，$\operatorname{Aut}O$ 在 $\mathcal{GR}$ 上的作用典范地提升为 $\operatorname{Aut}_Z O$ 在主丛（232）与线丛 $\lambda_{\mathcal{L}}$ 上的作用。$\operatorname{Aut}_Z O$ 在 $\lambda_{\mathcal{L}}$ 上的作用诱导出 $\operatorname{Der}O=\operatorname{Lie}\operatorname{Aut}_Z O$ 在 $\lambda_{\mathcal{L}}$ 上的作用。

\bdp{4.6.8} $Z=\operatorname{Aut}\mathcal{L}$ 在扩张（217）上的作用来自（215）。故 $Z$ 经由态射
\begin{equation}\tag{235}
Z\to H^0(\mathcal{GR},\mathcal{O}^*_{\mathcal{GR}})
\end{equation}
作用在 $\lambda_{\mathcal{L}}$ 上，该态射是（215）与自然嵌入 $\operatorname{Hom}(G(K),\mathbb{G}_m)\hookrightarrow H^0(\mathcal{GR},\mathcal{O}^*_{\mathcal{GR}})$ 之复合的逆。回顾 $\pi_0(\mathcal{GR})=Z^{\vee}$（见 4.5.9），故 $z\in Z$ 定义了 $f_z:\pi_0(\mathcal{GR})\to\mathbb{C}^*$，而（235）是映射 $z\mapsto f_z^{-1}$。

\bdp{4.6.9} \emph{注.}（这里需要它吗 ???）把 $\mathcal{GR}$ 上线丛的范畴视为 3.4.4 意义下的 $Z$-范畴，$Z$-结构由（235）定义。由 3.4.7(i) 有典范的 Picard 函子
\begin{equation}\tag{236}
Z\operatorname{tors}(O)=Z\operatorname{tors}\to\{\text{$\mathcal{GR}$ 上的线丛}\}.
\end{equation}
确切地说，（236）把 $\mathcal{E}\in Z\operatorname{tors}$ 对应到 $\mathcal{O}_{\mathcal{GR}}$ 的 $\mathcal{E}$-扭，并配以 $Z$-作用（235）。由 3.4.7(iv)，函子 $\mathcal{L}\mapsto\lambda_{\mathcal{L}}$（$\mathcal{L}\in Z\operatorname{tors}_\theta(O)$）关于 Picard 函子（236）是仿射的。
% [Tutor Note: 原书 p.185 4.6.9 中的「???」是作者自留的行内疑问记号，不是排印故障；底本如此，保留原样。]

\bdp{4.6.10} 由（56）定义的态射 $\alpha:\mu_2\to Z$ 诱导出 $\mu_2$ 在 $\lambda_{\mathcal{L}}$ 上的作用（$\mathcal{L}\in Z\operatorname{tors}_\theta(O)$）。它给出 $\lambda_{\mathcal{L}}$ 上的 $(\mathbb{Z}/2\mathbb{Z})$-分次。在 4.5.10 中我们引入了 $\mathcal{GR}$ 的偶分支与奇分支的概念。由 4.6.8，$(\mathbb{Z}/2\mathbb{Z})$-分次丛 $\lambda_{\mathcal{L}}$ 到 $\mathcal{GR}$ 的偶（相应地，奇）分支上的限制是偶（相应地，奇）的。

\bdp{4.6.11} 函子
\begin{equation}\tag{237}
Z\operatorname{tors}_\theta(O)\to\{\text{$\mathcal{GR}$ 上的线丛}\},\qquad \mathcal{L}\mapsto\lambda_{\mathcal{L}}
\end{equation}
是 3.4.4 意义下的 $Z$-函子，只要（237）右端的 $Z$-结构由（235）定义。由于 $Z\operatorname{tors}_\theta(O)$ 等价于 $\omega^{1/2}(O)\otimes_{\mu_2}Z$\SourcePage{186}（见 3.4.4），函子（237）可由相应的函子
\begin{equation}\tag{238}
\omega^{1/2}(O)\to\{\text{$\mathcal{GR}$ 上的线丛}\}
\end{equation}
重构，其中 $\omega^{1/2}(O)$ 是 $\omega(O)$ 的平方根构成的群胚。由于扩张（212）本质上来自「Clifford 扩张」（193），给出（238）的 Clifford 式描述是容易的。答案如下。

设 $\mathcal{L}\in\omega^{1/2}(O)$。我们固定了 $\mathfrak{g}$ 上非退化不变对称双线性型，故 Tate 空间 $V=V_{\mathcal{L}}:=\mathcal{L}\otimes_O(\mathfrak{g}\otimes K)$ 带有非退化对称双线性型（见 4.3.3），而 $L:=\mathcal{L}\otimes\mathfrak{g}\subset V$ 是 Lagrange c-格。置 $M=M_{\mathcal{L}}:=\operatorname{Cl}(V)/\operatorname{Cl}(V)L$；这是 $\operatorname{Cl}(V)$ 上不可约的 $(\mathbb{Z}/2\mathbb{Z})$-分次离散模。在 $V$ 中 Lagrange c-格构成的归纳概形 $\operatorname{Lagr}(V)$ 上有线丛 $\mathcal{P}_M$（见 4.3.2）。我们断言
\begin{equation}\tag{239}
\lambda_{\mathcal{L}}=\varphi^*\mathcal{P}_{M_{\mathcal{L}}},
\end{equation}
其中态射\footnote{容易证明 $\varphi$ 是闭嵌入，其像是满足 $O\Lambda=\Lambda$ 且 $\mathcal{L}^{-1}\otimes_O\Lambda$ 为 $\mathfrak{g}\otimes K$ 的李子代数的那些 $\Lambda\in\operatorname{Lagr}(V)$ 构成的归纳概形。} $\varphi:G(K)/G(O)\to\operatorname{Lagr}(V)$ 由 $\varphi(g):=gLg^{-1}$ 定义；换言之，
\begin{equation}\tag{240}
\parbox{0.82\linewidth}{\raggedright $\lambda_{\mathcal{L}}$ 在 $g\in G(K)/G(O)$ 处的纤维为 $M^{gLg^{-1}}:=\{m\in M_{\mathcal{L}}\mid(gLg^{-1})\cdot m=0\}$。}
\end{equation}
事实上，中心扩张（212）与由（193）诱导的扩张相反，因此 $\widetilde{O}(V)$ 在 $\mathcal{P}_{M_{\mathcal{L}}}$ 上的作用（见 4.3.2）诱导出 $\widetilde{G(K)}_{\mathcal{L}}$ 在 $\varphi^*\mathcal{P}_{M_{\mathcal{L}}}$ 上的作用，使得 $c\in\mathbb{G}_m\subset\widetilde{G(K)}_{\mathcal{L}}$ 通过乘以 $c^{-1}$ 作用；此外，$\varphi^*\mathcal{P}_{M_{\mathcal{L}}}$ 在 $\bar e$ 上的纤维是 $\mathbb{C}$。

显然同构（239）关于 $\mathcal{L}\in\omega^{1/2}(O)$ 是函子性的。

\bdp{4.6.12} \emph{注.}
\begin{enumerate}[label=(\roman*),leftmargin=2.2em]
\item 4.3.2 中的线丛 $\mathcal{P}_M$ 是 $(\mathbb{Z}/2\mathbb{Z})$-分次的。故（239）两边都是 $(\mathbb{Z}/2\mathbb{Z})$-分次的。（239）两边的分次都由 $\mu_2=\operatorname{Aut}\mathcal{L}$ 的作用诱导（对右端论证这一点时注意：\SourcePage{187}$\operatorname{Cl}(V)$ 上的 $(\mathbb{Z}/2\mathbb{Z})$-分次由 $\mu_2$ 在 $V$ 上的自然作用诱导）。因此（239）是分次同构。
\item 由 4.6.10，$-1\in\mu_2=\operatorname{Aut}\mathcal{L}$ 通过乘以 $(-1)^p$ 作用在（239）的右端，其中 $p$ 是奇偶性函数（230）。这也可由等式 $\chi=\theta$（见引理 4.3.4 的证明）与 4.3.4 末的注 (ii) 推出。
\end{enumerate}

\bdp{4.6.13} 我们应当思考超的方面，尤其是：一维超空间的逆是什么？（或许应在任意 Picard 范畴中表述；若它不是\emph{严格}交换的，可能会出问题。）

考虑 $G(O)$-轨道 $\operatorname{Orb}_\chi\subset\mathcal{GR}$，其中 $\chi\in P_+({}^{L}G)$（见 4.5.8）。我们将计算 $\lambda_{\mathcal{L},\chi}:=\lambda_{\mathcal{L}}$ 到 $\operatorname{Orb}_\chi$ 的限制，$\mathcal{L}\in Z\operatorname{tors}_\theta(O)$。由 4.6.4，$\lambda_{\mathcal{L},\chi}$ 是 $G(O)$-等变的。轨道 $\operatorname{Orb}_\chi$ 是 $\operatorname{Aut}^0 O$-不变的；由 4.6.7，$\lambda_{\mathcal{L},\chi}$ 是 $\operatorname{Aut}^0_Z O$-等变的，其中 $\operatorname{Aut}^0_Z O$ 是 $\operatorname{Aut}^0 O$ 在 $\operatorname{Aut}_Z O$ 中的原像（见（233））。最后，$\lambda_{\mathcal{L},\chi}$ 是 $\mathbb{Z}/2\mathbb{Z}$-分次的（但实际上 $\lambda_{\mathcal{L},\chi}$ 依 $\chi$ 而为偶或奇；此外，$\mathbb{Z}/2\mathbb{Z}$-分次可由 $Z\subset\operatorname{Aut}^0_Z O$ 的作用重构……）。群 $G(O)$ 与 $\operatorname{Aut}^0_Z O$ 也作用在典范层 $\omega_{\operatorname{Orb}_\chi}$ 上（$\operatorname{Aut}^0_Z O$ 经 $\operatorname{Aut}^0 O$ 作用）。在 4.6.17--4.6.19（???）中，我们将对某个一维向量空间 $\mathfrak{d}_{\mathcal{L},\chi}$ 构造典范同构
\begin{equation}\tag{241}
\lambda_{\mathcal{L},\chi}\xrightarrow{\ \sim\ }\omega_{\operatorname{Orb}_\chi}\otimes(\mathfrak{d}_{\mathcal{L},\chi})^{-1}.
\end{equation}
该空间带有 $G(O)$ 与 $\operatorname{Aut}^0_Z O$ 的作用，且（241）关于这些群等变。
% [Tutor Note: 原书 p.187 4.6.13 中的「???」是作者自留的交叉引用疑问记号（针对 4.6.17--4.6.19 的段号）；底本如此，保留原样。]

\bdp{4.6.14} 我们来定义 $\mathfrak{d}_{\mathcal{L},\chi}$。当然，$G(O)$ 在 $\mathfrak{d}_{\mathcal{L},\chi}$ 上的作用规定为平凡（$G(O)$ 没有非平凡特征标）。所以须对每个 $\chi$ 构造函子
\begin{equation}\tag{242}
Z\operatorname{tors}_\theta(O)\to\{\operatorname{Aut}^0_Z O\text{-mod}\},\qquad \mathcal{L}\mapsto\mathfrak{d}_{\mathcal{L},\chi}
\end{equation}

\SourcePage{188}

其中 $\{\operatorname{Aut}^0_Z O\text{-mod}\}$ 表示 $\operatorname{Aut}^0_Z O$-模的范畴。先定义函子
\begin{equation}\tag{243}
\omega^{1/2}(O)\to\{\operatorname{Aut}^0_Z O\text{-mod}\},\qquad \mathcal{L}\mapsto\mathfrak{d}_{\mathcal{L},\chi}.
\end{equation}
对 $\mathcal{L}\in\omega^{1/2}(O)$ 置
\begin{equation}\tag{244}
\mathfrak{d}_{\mathcal{L},\chi}:=(\mathcal{L}_0)^{\otimes d(\chi)},
\end{equation}
其中 $\mathcal{L}_0$ 是 $\mathcal{L}$ 在闭点 $0\in\operatorname{Spec}O$ 上的纤维，而
\begin{equation}\tag{245}
d(\chi):=(\chi,2\rho)=\dim\operatorname{Orb}_\chi.
\end{equation}
如下定义 $\operatorname{Aut}^0_Z O$ 在 $\mathfrak{d}_{\mathcal{L},\chi}$ 中的表示：$\operatorname{Aut}^0_2 O=\operatorname{Aut}^0(O,\mathcal{L})$ 以显然的方式作用，而 $Z\subset\operatorname{Aut}^0_Z O$ 经由
\begin{equation}\tag{246}
\chi_Z:Z\to\mathbb{G}_m
\end{equation}
作用，其中 $\chi_Z$ 是 $\chi\in P^+({}^{L}G)$ 到 $Z\subset{}^{L}G$ 的限制（这两个作用相容，因为 $\chi_Z$ 与态射（56）的复合把 $-1\in\mu_2$ 映到 $(-1)^{(\chi,2\rho)}$）。

于是我们构造了（243）。$\omega^{1/2}(O)$ 是 3.4.4 意义下的 $\mu_2$-范畴，$\{\operatorname{Aut}^0_Z O\text{-mod}\}$ 是 $Z$-范畴，而（243）是 $\mu_2$-函子（$\{\operatorname{Aut}^0_Z O\}$ 上的 $\mu_2$-结构来自态射（56），等价地，来自典范嵌入 $\mu_2\to\operatorname{Aut}^0_2 O$）。故（243）诱导出 $Z$-函子
\[
Z\operatorname{tors}_\theta(O)=\omega^{1/2}(O)\otimes_{\mu_2}Z\to\{\operatorname{Aut}^0_Z O\text{-mod}\}.
\]
这就是（242）的定义。

\bdp{4.6.15} 显然 $\operatorname{Lie}\operatorname{Aut}^0_Z O=\operatorname{Der}^0 O$ 如下作用在一维空间 $\mathfrak{d}_{\mathcal{L},\chi}$ 上：
\begin{equation}\tag{247}
L_0\mapsto(\chi,\rho)=-\frac{1}{2}\dim\operatorname{Orb}_\chi,\qquad L_n\mapsto0\quad(n>0).
\end{equation}
按惯例，$L_n:=-t^{n+1}\frac{d}{dt}\in\operatorname{Der}^0 O$。
% [Tutor Note: 原书 p.188 把该特征标记作 $(\chi,\rho)$，而 (245) 已给出 $d(\chi)=(\chi,2\rho)=\dim\operatorname{Orb}_\chi$，两式相差一个符号。数值 $-\frac{1}{2}\dim\operatorname{Orb}_\chi$ 本身与约定 $L_n:=-t^{n+1}\frac{d}{dt}$ 相符（$L_0$ 在 $\omega^{1/2}$ 的纤维上乘 $-1/2$），故偏差只在记号 $(\chi,\rho)$。底本如此，保留原样。]

\SourcePage{189}

\bdp{4.6.16} \emph{注.} 4.6.14 中 $\mathfrak{d}_{\mathcal{L},\chi}$ 的定义可改述如下。利用 3.4.5 的等价 $Z\operatorname{tors}_\theta(O)\xrightarrow{\ \sim\ }\widetilde{Z}\operatorname{tors}_\omega(O)$，我们把 $\mathcal{L}\in Z\operatorname{tors}_\theta(O)$ 按（59）解释为 $\mathbb{G}_m$-torsor $\omega_O$ 到 $\widetilde{Z}$-torsor 的提升。由（62）有典范态射 $\widetilde{Z}\to{}^{L}H$，其中 ${}^{L}H$ 是 ${}^{L}G$ 的嘉当环面；或者同样地，${}^{L}H$ 是 ${}^{L}G$ 的一个嘉当子群，并配有一个包含它的固定 Borel 子群。记 $\chi_{\widetilde{Z}}$ 为 $\widetilde{Z}\to{}^{L}H$ 与 $\chi:{}^{L}H\to\mathbb{G}_m$ 的复合。$\operatorname{Spec}O$ 上的 $\widetilde{Z}$-torsor $\mathcal{L}$ 与一维表示 $\chi_{\widetilde{Z}}:\widetilde{Z}\to\mathbb{G}_m$ 定义了 $\operatorname{Spec}O$ 上的线丛 $\mathfrak{d}^{\mathcal{O}}_{\mathcal{L},\chi}$。按 4.6.6，$\operatorname{Aut}_Z O=\operatorname{Aut}(O,\mathcal{L})$；故 $\operatorname{Aut}O$ 在 $\operatorname{Spec}O$ 上的作用典范地提升为 $\operatorname{Aut}_Z O$ 在 $\mathfrak{d}^{\mathcal{O}}_{\mathcal{L},\chi}$ 上的作用。因此 $\operatorname{Aut}^0_Z O$ 作用在 $\mathfrak{d}^{\mathcal{O}}_{\mathcal{L},\chi}$ 在 $0\in\operatorname{Spec}O$ 处的纤维上。读者容易把该纤维与 4.6.14 中的 $\mathfrak{d}_{\mathcal{L},\chi}$ 等同。

\bdp{4.6.17} 我们来对 $\mathcal{L}\in\omega^{1/2}(O)$ 构造同构（241）。我们用 $\lambda_{\mathcal{L}}$ 的 Clifford 式描述。照 4.6.11 置 $V=V_{\mathcal{L}}:=\mathcal{L}\otimes_O(\mathfrak{g}\otimes K)$，$L:=\mathcal{L}\otimes\mathfrak{g}\subset V$，$M=M_{\mathcal{L}}:=\operatorname{Cl}(V)/\operatorname{Cl}(V)L$。对 $x\in\mathcal{GR}=G(K)/G(O)$ 置 $L_x:=gLg^{-1}$，其中 $g$ 是 $x$ 在 $G(K)$ 中的原像。由（240），$\lambda_{\mathcal{L}}$ 在 $x$ 处的纤维等于
\begin{equation}\tag{248}
M^{L_x}:=\{m\in M_{\mathcal{L}}\mid L_x\cdot m=0\}.
\end{equation}
设 $x\in\operatorname{Orb}_\chi$。因 $\operatorname{Orb}_\chi$ 是 $x$ 的 $G(O)$-轨道，$\operatorname{Orb}_\chi$ 在 $x$ 处的切空间为 $(\mathfrak{g}\otimes O)/((\mathfrak{g}\otimes O)\cap g(\mathfrak{g}\otimes O)g^{-1})=L^{-1}\otimes_O(L/(L\cap L_x))$，其中 $g\in G(K)$ 是 $x$ 的原像。于是 $\omega^{-1}_{\operatorname{Orb}_\chi}$ 在 $x$ 处的纤维等于 $(\mathcal{L}_0)^{\otimes-d(\chi)}\otimes\det({L/(L\cap L_x)})$，其中 $d(\chi)=\dim\operatorname{Orb}_\chi$。计及（244），可见（241）右端在 $x$ 处的纤维等于
\begin{equation}\tag{249}
(\det({L/(L\cap L_x)}))^{-1}.
\end{equation}
于是只需构造同构
\begin{equation}\tag{250}
\det({L/(L\cap L_x)})\otimes M^{L_x}\xrightarrow{\ \sim\ }\mathbb{C}.
\end{equation}
% [Tutor Note: 原书 p.189 用 $\omega^{-1}_{\operatorname{Orb}_\chi}$ 的纤维代入 (241)（(241) 本身写的是 $\omega_{\operatorname{Orb}_\chi}$）；再与 (244) 的 $\mathfrak{d}_{\mathcal{L},\chi}=(\mathcal{L}_0)^{\otimes d(\chi)}$ 联立，右端纤维为 $\det(L/(L\cap L_x))\otimes(\mathcal{L}_0)^{\otimes-2d(\chi)}$，与 (249) 所记的 $(\det(L/(L\cap L_x)))^{-1}$ 表面不符。疑原文 $\omega$ 与 $\omega^{-1}$ 的上标混用；底本如此，保留原样。]

\SourcePage{190}

\bdp{4.6.18} \emph{引理.} 考虑带非退化对称双线性型的 Tate 空间 $V$。设 $L,\Lambda\subset V$ 是 Lagrange c-格，$M$ 是 Clifford 代数 $\operatorname{Cl}(V)$ 上的不可约离散模。考虑由自然映射 $\bigwedge^dL\to\bigwedge^dV\to\operatorname{Cl}(V)$ 诱导的算子
\begin{equation}\tag{251}
\bigwedge\nolimits^{d}L\otimes M\to M.
\end{equation}
若 $d=\dim L/(L\cap\Lambda)$，则（251）诱导同构
\begin{equation}\tag{252}
\bigwedge\nolimits^{d}(L/(L\cap\Lambda))\otimes M^{\Lambda}\xrightarrow{\ \sim\ }M^{L}.
\end{equation}
证明可化归到 $\dim V<\infty$ 且 $V=L\oplus\Lambda$ 的情形。

\bdp{4.6.19} 我们定义（250）为 $\Lambda=L_x$ 时的同构（252）（在 4.6.17 的情形下 $M^L=\mathbb{C}$）。于是对 $\mathcal{L}\in\omega^{1/2}(O)$，我们构造了同构（241），它关于 $G(O)$ 与 $\operatorname{Aut}^0_2 O=\operatorname{Aut}^0(O,\mathcal{L})$ 等变。

记 $\mathcal{C}_\chi$ 为 $\operatorname{Orb}_\chi$ 上线丛的范畴。（241）两边都是 $\mu_2$-函子 $\omega^{1/2}(O)\to\mathcal{C}_\chi$，并延拓为 $Z$-函子
\[
Z\operatorname{tors}_\theta(O)=\omega^{1/2}(O)\otimes_{\mu_2}Z\to\mathcal{C}_\chi
\]
（$\mathcal{C}_\chi$ 上的 $Z$-结构由与（246）相逆的 $Z$ 的特征标定义）；对（241）左端，这由 4.6.8 推出。显然（241）是函子 $\omega^{1/2}(O)\to\mathcal{C}_\chi$ 的同构。因此（241）是函子 $Z\operatorname{tors}_\theta(O)\to\mathcal{C}_\chi$ 的同构。同构（241）是 $\operatorname{Aut}^0_Z O$-等变的，因为它既 $\operatorname{Aut}^0_2 O$-等变又 $Z$-等变。

\bdp{4.6.20} 回顾 $\lambda_{\mathcal{L}}$ 依赖于 $\mathfrak{g}$ 上非退化不变双线性型的选取（见 4.6.3 与 4.4.7）。如 4.4.7 的脚注所解释，$\lambda_{\mathcal{L}}$ 有更典范的版本。当 $G$ 单时，该版本 $\lambda_{\mathcal{L}}^{can}$ 依赖于 $\beta^{1/2}$ 的选取，其中 $\beta$ 是 $\mathfrak{g}$ 上不变双线性型构成的线（参看 4.4.5）；$\lambda_{\mathcal{L}}^{can}$ 来自下述版本的（212）：\SourcePage{191}%
以 $SO(\mathfrak{g}\otimes\beta^{1/2})$ 代替 $SO(\mathfrak{g})$。容易看出，$\lambda_{\mathcal{L}}^{can}$ 上相应于 $-1\in\operatorname{Aut}\beta^{1/2}$ 的作用的 $(\mathbb{Z}/2\mathbb{Z})$-分次，与 4.6.10 的分次一致。（241）的「典范」版本是同构
\begin{equation}\tag{253}
\lambda^{can}_{\mathcal{L},\chi}\xrightarrow{\ \sim\ }\omega_{\operatorname{Orb}_\chi}\otimes(\mathfrak{d}_{\mathcal{L},\chi})^{-1}\otimes(\beta^{1/2})^{\otimes-d(\chi)},
\end{equation}
其中 $d(\chi)$ 由（245）定义。细节留给读者。
% ===================== end 源码/units/ch04_u07.tex =====================
